\section{The complementarity theorem}
\label{sec:complementarity-theorem}

% Macro safety net for the BRST functor of chap:universal-conductor.
% Guarded by \providecommand for portability across standalone extraction.
\providecommand{\BRST}{\mathrm{BRST}}

Classical Koszul duality over a point already distinguishes the bar
coalgebra $B(\cA)$, the inversion $\Omega(B(\cA)) \simeq \cA$, the
dual coalgebra $\cA^i = H^*(B(\cA))$, and the dual algebra
$\cA^! = ((\cA^i)^\vee)$.
Already at genus~$0$ on a curve, the Fulton--MacPherson
compactifications $\overline{C}_n(X)$ and the Arnold forms
$\eta_{ij} = d\log(z_i - z_j)$ supply configuration-space geometry
that has no counterpart over a point: the collision residues that
define the bar differential, the scalar projection of the spectral
$r$-matrix which gives the modular characteristic
$\kappa(\cA)$ on the standard landscape, and the higher operations
$m_n$ for $n \geq 3$ that determine whether the shadow tower
terminates. At genus~$g \geq 1$, the Hodge bundle
$\mathbb{E}_g=\pi_*\omega_\pi$ on $\overline{\mathcal{M}}_g$
introduces a curvature two-form
$\omega_{\mathbb{E}}=c_1(\det\mathbb{E}_g)$, and on the scalar
diagonal of the curved fiber bar model
\[
  \dfib^{\,2}
  =
  \kappa(\cA)\cdot\omega_{\mathbb{E}}.
\]
This scalar curvature identity is not Theorem~C. It is the genus-$g$
curvature surface on which Theorem~D projects to the Hodge class
$\lambda_g=c_g(\mathbb{E}_g)$ and to the Faber--Pandharipande number
$\lambda_g^{\mathrm{FP}}$. Theorem~C acts on a different object: the
strict flat ambient complex
$\mathbf{C}_g(\cA)=R\Gamma(\overline{\mathcal{M}}_g,\mathcal{Z}(\cA))$.
When the Verdier--Koszul involution is represented on this complex,
its eigenspaces give
\begin{equation}\label{eq:complementarity-chapter-opening}
 H^*\bigl(\overline{\mathcal{M}}_g,\,\mathcal{Z}(\cA)\bigr)
 \;=\;
 Q_g(\cA) \;\oplus\; Q_g(\cA^!),
\end{equation}
with $Q_g(\cA)=\ker(\sigma-\mathrm{id})$ and
$Q_g(\cA^!)=\ker(\sigma+\mathrm{id})$ after passage to cohomology.
If the Verdier pairing is non-degenerate and anti-invariant, these
eigenspaces are Lagrangian; without that pairing they are only
complementary eigenspaces.

The scalar complementarity constants are the shadow of this
eigenspace theorem after Theorem~D has been applied on its scalar
lane. On the computed standard landscape the constants are
\[
\begin{array}{c|c}
\text{family} & \kappa(\cA)+\kappa(\cA^!) \\ \hline
\text{free fields and affine Kac--Moody} & 0 \\
\mathrm{Vir}_c=\cW_2 & 13 \\
\cW_3 & 250/3 \\
\text{Bershadsky--Polyakov} & 98/3 \\
\text{Vol III Mukai-enhanced K3 Heisenberg witness} & 8
\end{array}
\]
For principal $\cW_N$ the general formula is
$\kappa+\kappa^!=\varrho_NK_N$, with
$\varrho_N=H_N-1$ and
$K_N=2(N-1)(2N^2+2N+1)$. These are scalar identities. They do
not identify the bar coalgebra with the Koszul dual, do not identify
the center local system with the chiral derived center, and do not
construct the bar-side BV pairing. Those are separate objects and
separate open mathematical conditions.

Theorem~C (Theorem~\ref{thm:quantum-complementarity-main})
therefore has three layers with distinct hypotheses.
\begin{enumerate}[label=\textup{(C\arabic*)},start=0]
\item \emph{Fiber--center identification
 \textup{(}unconditional in coderived form on the Koszul locus\textup{)}.}\;
 For every chiral algebra~$\cA$ carrying a modular pre-Koszul datum,
 the curved fiber bar family
 $(\barB^{(g)}(\cA),\dfib)$ determines a well-defined coderived
 object and is read through its strict flat comparison model
 $(\barB^{(g)}(\cA),\Dg{g})$
 \textup{(}Proposition~\textup{\ref{prop:gauss-manin-uncurving-chain}}\textup{)}.
 On the flat perfect locus, and in particular when
 $\kappa(\cA)=0$, the ordinary derived pushforward satisfies
 $\mathcal{H}^q(R\pi_{g*}\bar{B}^{(g)}_{\mathrm{flat}}(\cA)) = 0$
 for $q \neq 0$, with
 $\mathcal{H}^0(R\pi_{g*}\bar{B}^{(g)}_{\mathrm{flat}}(\cA))
 \cong \mathcal{Z}_{\cA}$
 \textup{(}Theorem~\textup{\ref{thm:fiber-center-identification})}.
 This produces the ambient complex on which~\textup{(C1)}
 and~\textup{(C2)} operate.

\item \emph{Verdier eigenspace decomposition \textup{(}strict flat
 representative\textup{)}.}\;
 For every chiral Koszul pair $(\cA, \cA^!)$ on a smooth projective
 curve~$X$ and every genus $g \geq 0$ for which the Verdier
 involution is represented on the strict flat ambient complex, the
 cohomology of that ambient complex splits as
 \begin{equation}\label{eq:complementarity-summary}
 H^*(\overline{\mathcal{M}}_g,\,\mathcal{Z}(\cA))
 \;=\;
 Q_g(\cA) \;\oplus\; Q_g(\cA^!),
 \end{equation}
where $Q_g(\cA) = \ker(\sigma - \mathrm{id})$ and
$Q_g(\cA^!) = \ker(\sigma + \mathrm{id})$ are the eigenspaces of
 the Verdier involution~$\sigma$. For $g \geq 1$, the Verdier pairing
 identifies $Q_g(\cA) \cong Q_g(\cA^!)^\vee$; at $g=0$, the point
 class is $\sigma$-fixed, so
 $Q_0(\cA)=H^0(\overline{\mathcal{M}}_0,\mathcal{Z}(\cA))$ and
 $Q_0(\cA^!)=0$
 \textup{(}Theorem~\textup{\ref{thm:quantum-complementarity-main})}.

\item \emph{Shifted-symplectic Lagrangian upgrade
 \textup{(}conditional on perfectness and nondegeneracy\textup{)}.}\;
 When the strict flat relative bar family
 $R\pi_{g*}\barB^{(g)}_{\mathrm{flat}}(\cA)$ is perfect
 \textup{(}Lemma~\textup{\ref{lem:perfectness-criterion})} and the
 Verdier pairing $\langle -,- \rangle_{\mathbb{D}}$ on the ambient
 complex $\mathbf{C}_g(\cA) = R\Gamma(\overline{\mathcal{M}}_g,
 \mathcal{Z}(\cA))$ is non-degenerate, the complex
 $\mathbf{C}_g(\cA)$ carries a $({-}(3g{-}3))$-shifted symplectic
 structure \textup{(}PTVV
 \cite{PTVV13}\textup{)}, and the summands $\mathbf{Q}_g(\cA)$ and
 $\mathbf{Q}_g(\cA^!)$ are Lagrangian subspaces: isotropic of half
 rank, with $\mathbf{Q}_g(\cA) \simeq
 \mathbf{Q}_g(\cA^!)^\vee[-(3g{-}3)]$.
 This is the geometric home of the finite-window shadow
 tower~$\Theta_{\cA}^{\leq r}$
 \textup{(}Appendix~\textup{\ref{app:nonlinear-modular-shadows})}.
\end{enumerate}
The conditionality in~(C2) is substantive:
perfectness requires PBW filterability and finite-dimensional
flat fiber cohomology (Lemma~\ref{lem:perfectness-criterion}), and
nondegeneracy of the Verdier pairing is a hypothesis verified
family by family
(Proposition~\ref{prop:standard-examples-modular-koszul}).
The fiber--center identification~(C0) holds unconditionally in
coderived form on the Koszul locus, and its ordinary-derived
realization holds on the flat perfect locus; the
S-level decomposition~(C1) is read on that strict flat
representative; the H-level Lagrangian upgrade~(C2) is the
geometric content of the nonlinear theory.

\begin{proposition}[Hodge curvature and scalar projection;
\ClaimStatusConditional]
\label{prop:hodge-curvature-scalar-projection}
\index{Hodge bundle!curvature and scalar projection}
\index{cross-channel correction!scalar projection}
Let $\mathbb{E}_g=\pi_*\omega_\pi$ be the Hodge bundle on
$\overline{\mathcal{M}}_g$. The curvature class entering the
fiberwise bar identity is the degree-$2$ class
\[
  \omega_{\mathbb{E}}=c_1(\det\mathbb{E}_g),
  \qquad
  \dfib^{\,2}=\kappa(\cA)\omega_{\mathbb{E}}
\]
on the scalar diagonal. The genus-$g$ scalar obstruction is the
top-Hodge projection
\[
  \operatorname{obs}_g(\cA)=\kappa(\cA)\lambda_g,
  \qquad
  \lambda_g=c_g(\mathbb{E}_g),
\]
on the uniform-weight lane. After integration against the
Faber--Pandharipande class this gives
\[
  F_g(\cA)
  =
  \kappa(\cA)\lambda_g^{\mathrm{FP}},
  \qquad
  \lambda_g^{\mathrm{FP}}
  =
  \frac{2^{2g-1}-1}{2^{2g-1}}\,
  \frac{|B_{2g}|}{(2g)!}.
\]
For a multi-weight modular Koszul algebra the all-weight scalar
free energy is
\[
  F_g(\cA)
  =
  \kappa(\cA)\lambda_g^{\mathrm{FP}}
  +
  \delta F_g^{\mathrm{cross}}(\cA),
\]
with $\delta F_1^{\mathrm{cross}}=0$ and with nonzero
cross-channel terms possible at $g\geq2$.
\end{proposition}

\begin{proof}
The Hodge bundle has determinant line $\det\mathbb{E}_g$ and top
Chern class $c_g(\mathbb{E}_g)$. Chern--Weil theory distinguishes
the curvature two-form $c_1(\det\mathbb{E}_g)$ from the top class
$c_g(\mathbb{E}_g)$. The curved fiber identity is a chain-level
degree-$2$ statement, hence uses the former. Theorem~D applies the
tautological genus-$g$ projection and produces the top-Hodge class
$\lambda_g=c_g(\mathbb{E}_g)$; integrating gives the
Faber--Pandharipande number displayed above
\cite{FP03}. The multi-weight formula is exactly the
cross-channel refinement of
Theorem~\ref{thm:multi-weight-genus-expansion}; it reduces to the
diagonal term at genus~$1$ and on the uniform-weight lane.
\end{proof}

\begin{proposition}[Object firewall for Theorem~C;
\ClaimStatusProvedHere]
\label{prop:theorem-c-object-firewall}
\index{object firewall!Theorem C}
\index{derived center!distinguished from center local system}
\index{chiral Hochschild!distinguished from Theorem C}
For a chiral Koszul pair $(\cA,\cA^!)$ the following objects are
pairwise typed differently:
\[
\cA,\qquad
\barB^{\mathrm{ch}}(\cA),\qquad
\cA^i:=H^*(\barB^{\mathrm{ch}}(\cA)),\qquad
\cA^!,\qquad
\mathcal{Z}(\cA),\qquad
Z^{\mathrm{der}}_{\mathrm{ch}}(\cA)
=\ChirHoch^\bullet(\cA,\cA).
\]
Theorem~C uses the center local system
$\mathcal{Z}(\cA)$ through
$R\Gamma(\overline{\mathcal{M}}_g,\mathcal{Z}(\cA))$.
Theorem~H uses the chiral Hochschild complex
$\ChirHoch^\bullet(\cA,\cA)$, the chiral derived center.
The fiber--center identification~\textup{(C0)} identifies the
degree-$0$ flat fiber cohomology sheaf with $\mathcal{Z}(\cA)$ on
the flat perfect locus; it does not identify the raw curved bar
object with the chiral derived center, it does not identify
$\cA^i$ with $\cA^!$ before dualization, and it does not identify
$\Omega\barB(\cA)$ with $\cA^!$.
\end{proposition}

\begin{proof}
The six objects have different functor signatures. The algebra
$\cA$ is a chiral algebra on~$X$; $\barB^{\mathrm{ch}}(\cA)$ is a
factorization coalgebra; $\cA^i$ is the bar cohomology coalgebra;
$\cA^!$ is the Verdier/Koszul dual algebra obtained from
dualizing the finite-type bar-cohomology object in the appropriate
completed category; the sheaf $\mathcal{Z}(\cA)$ is the degree-$0$
center local system appearing after flat fiber concentration; and
$Z^{\mathrm{der}}_{\mathrm{ch}}(\cA)$ is the full chiral
Hochschild cochain object with its brace operations. The maps
among them are comparison functors or projections. A scalar equality
for $\kappa$, a Verdier eigenspace decomposition of
$R\Gamma(\overline{\mathcal{M}}_g,\mathcal{Z}(\cA))$, and
Theorem~H's Hochschild concentration therefore live in different
categories and cannot be substituted for one another.
\end{proof}

\begin{remark}[Calabi--Yau exchange and the perfectness hypothesis]
\label{rem:hr24-cy-interchange}
Holstein--Rivera~\cite{HR24} prove that Koszul duality
between dg categories and pointed curved coalgebras exchanges
smooth and proper Calabi--Yau structures. On a finite or completed
Koszul surface satisfying their hypotheses, smooth
$\mathrm{CY}_1$ data on~$\cA$, given by a nondegenerate invariant
bilinear form, correspond to proper $\mathrm{CY}_1$ data on
$\barB(\cA)$. This supports finite cohomology at each fixed finite
weight. It does not prove family-level perfectness, and it does not
make the dual regularity hypothesis~(P3) redundant. The
perfectness required by~(C2) still uses PBW filterability, total
flat-fiber cohomology, bounded pseudo-coherence, and the corresponding
regularity hypotheses for~$\cA^!$
(Lemma~\ref{lem:perfectness-criterion}).
Calaque--Safronov~\cite{CS24} develop the relative AKSZ
framework for shifted Lagrangian morphisms that provides
the natural geometric setting for the~(C2) upgrade.
\end{remark}

\subsection{Physical and mathematical motivation}

\begin{remark}[Physical origin]
The genus-$g$ partition function
$Z_g[\mathcal{A}] = \int_{\mathcal{M}_g} \langle \mathcal{A} \rangle_{\Sigma_g}
\cdot e^{-S[\Sigma_g]}$
splits quantum corrections into deformations (marginal operators)
and obstructions (anomalies). The complementarity theorem gives
this split algebraic content: what $\mathcal{A}$ sees as obstruction,
$\mathcal{A}^!$ sees as deformation. The conformal
anomaly of matter is the ghost-number violation of ghosts.
\end{remark}

\begin{remark}[Kodaira--Spencer action on tautological classes]
The center $Z(\cA)$ acts on $H^*(\overline{\mathcal{M}}_g)$
($\lambda$-classes, $\psi$-classes, boundary classes) via the
Kodaira--Spencer map
$\rho\colon Z(\cA)\to\mathrm{End}(H^*(\overline{\mathcal{M}}_g))$,
and the eigenspace decomposition gives
$H^*(\overline{\mathcal{M}}_g,Z(\cA))
=Q_g(\cA)\oplus Q_g(\cA^!)$.
\end{remark}

\begin{remark}[Algebraic structure]
The involution $\cA \mapsto \cA^!$ with $(\cA^!)^! \simeq \cA$
does not force arbitrary functorial invariants to split. The
complementarity theorem supplies a specific involution~$\sigma$ on
the ambient center complex
$R\Gamma(\overline{\mathcal{M}}_g,\mathcal{Z}(\cA))$; its
idempotents give the direct and exhaustive decomposition. The word
Lagrangian is reserved for the additional non-degenerate
anti-invariant Verdier pairing.
\end{remark}

\begin{remark}[Computational perspective: Heisenberg at genus~$1$]
For $\mathcal{H}_\kappa$ at genus~$1$
(Convention~\ref{conv:heisenberg-kappa-notation}: level $=$ modular
characteristic $\kappa(\mathcal{H}_\kappa) = \kappa$):
$H^*(\overline{\mathcal{M}}_{1,1})=\mathbb{Q}\oplus\mathbb{Q}\lambda$,
$Q_1(\mathcal{H}_\kappa)=\mathbb{C}\cdot\kappa$ (central extension),
$Q_1(\mathcal{H}_\kappa^!)=\mathbb{C}\cdot\lambda$ (curved structure),
and $Q_1\oplus Q_1^!=H^0\oplus H^2=\mathbb{C}^2$.
\end{remark}

\begin{remark}[Complementarity as transversality of boundary conditions]
\label{rem:complementarity-transversality}
\index{complementarity!transversality of boundaries}%
\index{Dimofte!transversality and complementarity}%
In the holomorphic-topological framework of
Costello--Dimofte--Gaiotto~\cite{CDG2024,CG17}, a $3$d HT theory
on $\bC \times [0,1]$ places boundary conditions $\cA$
and~$\cA^!$ at the two ends of the interval. Complementarity is
the algebraic manifestation of the transversality of these
boundary conditions: the genus-$g$ obstruction space
$H^*(\overline{\mathcal{M}}_g, \mathcal{Z}(\cA))$ decomposes
into the deformation space of one boundary and the obstruction
space of the other. What~$\cA$ sees as anomaly at genus~$g$,
the transverse boundary~$\cA^!$ absorbs as deformation.
The Lagrangian upgrade~(C2) is the statement that the two
boundaries meet in a shifted-symplectic ambient space at
half rank after perfectness and nondegeneracy have been verified:
transversality in the sense of
Pantev--To\"en--Vaqui\'e--Vezzosi~\cite{PTVV13}.
\end{remark}

\subsection{Statement of the theorem}

\begin{definition}[Deformation-obstruction complexes: H-level]
\label{def:complementarity-complexes}
\index{deformation-obstruction complexes|textbf}
Let $(\cA, \cA^!)$ be a chiral Koszul pair on~$X$.
The \emph{ambient complex} at genus~$g$ is
\[
\mathbf{C}_g(\cA) := R\Gamma(\overline{\mathcal{M}}_g, \mathcal{Z}(\cA)),
\]
where $\mathcal{Z}(\cA)$ is the center local system.
The Verdier involution $\sigma$ acts on $\mathbf{C}_g(\cA)$ as a
cochain-level endomorphism with $\sigma^2 = \mathrm{id}$
(Theorem~\ref{thm:verdier-bar-cobar}). Let
$p^\pm=\frac{1}{2}(\mathrm{id}\pm\sigma)$.
Define the \emph{deformation-obstruction complexes} as the strict
idempotent eigensummands:
\begin{equation}\label{eq:complementarity-fiber}
\mathbf{Q}_g(\cA) := \operatorname{im}(p^+ \colon
\mathbf{C}_g \to \mathbf{C}_g), \qquad
\mathbf{Q}_g(\cA^!) := \operatorname{im}(p^- \colon
\mathbf{C}_g \to \mathbf{C}_g).
\end{equation}
The \emph{cohomological shadows} (S-level) are defined by:
\begin{align}\label{eq:shadow-spaces}
Q_g(\cA) &:= H^*(\mathbf{Q}_g(\cA))
= \ker(\sigma - \mathrm{id} \mid H^*(\overline{\mathcal{M}}_g, \mathcal{Z}(\cA))),\\
Q_g(\cA^!) &:= H^*(\mathbf{Q}_g(\cA^!))
= \ker(\sigma + \mathrm{id} \mid H^*(\overline{\mathcal{M}}_g, \mathcal{Z}(\cA))).
\end{align}
\end{definition}

\begin{remark}[H/M/S layers; Convention~\ref{conv:hms-levels}]
The idempotent summands~\eqref{eq:complementarity-fiber} are the
H-level objects; the bar complexes
$(\bar{B}^{(g)}(\cA), \Dg{g})$ are M-level models; the kernel
decomposition~\eqref{eq:shadow-spaces} is the S-level shadow.
The involution-splitting lemma
(Lemma~\ref{lem:involution-splitting}) makes the passage automatic
in characteristic~$0$, with no cone contribution.
\end{remark}

\begin{lemma}[Involution splitting in characteristic~\texorpdfstring{$0$}{0}; \ClaimStatusProvedHere]
\label{lem:involution-splitting}
\index{involution splitting lemma|textbf}
Let $V$ be a cochain complex over a field of characteristic $\neq 2$,
and let $\sigma \colon V \to V$ be a cochain involution
($\sigma^2 = \mathrm{id}$, $\sigma \circ d = d \circ \sigma$). Then:
\begin{enumerate}[label=\textup{(\alph*)}]
\item The idempotent eigensummands
 $V^+ := \operatorname{im}\bigl(\frac{1}{2}(\mathrm{id}+\sigma)\bigr)$ and
 $V^- := \operatorname{im}\bigl(\frac{1}{2}(\mathrm{id}-\sigma)\bigr)$
 split $V$ as a direct sum of complexes:
 $V^+ \oplus V^- \xrightarrow{\;\cong\;} V$.
\item On cohomology:
 $H^*(V) = \ker(\sigma - \mathrm{id}) \oplus \ker(\sigma + \mathrm{id})$.
\item If $\langle -,- \rangle$ is a cochain-level pairing on~$V$
 with $\langle \sigma x, \sigma y \rangle = -\langle x, y \rangle$
 \textup{(}anti-symmetric under $\sigma$\textup{)}, then $V^+$ and
 $V^-$ are isotropic for $\langle -,- \rangle$; if the pairing is
 non-degenerate, they are Lagrangian.
\end{enumerate}
\end{lemma}

\begin{proof}
The projectors $p^\pm := \tfrac{1}{2}(\mathrm{id} \pm \sigma)$
are cochain maps with $p^+ + p^- = \mathrm{id}$,
$(p^\pm)^2 = p^\pm$, and $p^+ \circ p^- = 0$.
Part~(a) is the direct-sum decomposition
$V=\operatorname{im}(p^+)\oplus\operatorname{im}(p^-)$.
Part~(b) follows by applying $H^*$ to~(a).
Part~(c): for $x, y \in V^+$ we have $\sigma x = x$, $\sigma y = y$,
so $\langle x, y \rangle = \langle \sigma x, \sigma y \rangle
= -\langle x, y \rangle$, hence $\langle x, y \rangle = 0$.
The same argument applies to $V^-$. If the pairing is non-degenerate,
then the induced pairing between $V^+$ and $V^-$ is perfect, since any
vector pairing trivially with the opposite summand pairs trivially
with all of~$V$. Thus each summand is Lagrangian.
\end{proof}

\begin{definition}[Relative bar fiber complex]\label{def:relative-bar-fiber}
\index{bar complex!relative fiber|textbf}
Let $(\cA, \cA^!)$ be a chiral Koszul pair on~$X$, and let
\[
\pi_g\colon \mathcal{C}_g(\cA) \longrightarrow
\overline{\mathcal{M}}_g
\]
be the family of genus-$g$ compactified configuration spaces
carrying the relative bar complex $\bar{B}^{(g)}(\cA)$.
Its stalk at $[\Sigma] \in \overline{\mathcal{M}}_g$ is the
\emph{genus-$g$ fiber bar complex}
$C^{(g)}_X(\cA)\!\mid_\Sigma$, i.e.\ the bar complex of~$\cA$
on the fixed curve~$\Sigma$ with the fiberwise curved
differential~$\dfib$
\textup{(}Convention~\textup{\ref{conv:higher-genus-differentials})}.
\end{definition}

\begin{lemma}[Perfectness criterion for the strict flat relative bar family;
\ClaimStatusProvedHere]
\label{lem:perfectness-criterion}
\index{perfectness criterion|textbf}
\index{bar complex!relative perfectness|textbf}
Let $(\cA, \cA^!)$ be a chiral Koszul pair on a smooth projective
curve~$X$, carrying a modular pre-Koszul datum
\textup{(}Definition~\textup{\ref{def:modular-koszul-chiral})}.
Write
\[
\barB^{(g)}_{\mathrm{flat}}(\cA)
\;:=\;
\bigl(\barB^{(g)}(\cA), \Dg{g}\bigr)
\]
for the strict flat comparison family of
Convention~\textup{\ref{conv:higher-genus-differentials}}.
Suppose the following two conditions hold:
\begin{enumerate}[label=\textup{(\roman*)}]
\item \emph{PBW filterability.}
 The genus-$g$ bar complex $\bar{B}^{(g)}(\cA)$ admits an
 exhaustive multiplicative filtration~$F_\bullet$ such that the
 associated graded $\operatorname{gr}_F \bar{B}^{(g)}(\cA)$
 sees only the genus-$0$ collision differential~$\dzero$ and is
 Koszul-acyclic \textup{(}axioms \textup{MK1} and
 \textup{MK3} of
 Definition~\textup{\ref{def:modular-koszul-chiral})}.

 \item \emph{Finite-dimensional flat fiber cohomology.}
 For each closed point
 $[\Sigma] \in \overline{\mathcal{M}}_g$,
 $\dim_\mathbb{C} H^n\bigl(
 \bar{B}^{(g)}_{\mathrm{flat}}(\cA)\big|_\Sigma\bigr)
 < \infty$ for all~$n$.

\item \emph{Pseudo-coherent bounded family.}
 The filtered strict flat family is pseudo-coherent over
 $\overline{\mathcal{M}}_g$, bounded in a uniform finite
 cohomological interval, and its filtration is finite on every
 cohomological truncation.
\end{enumerate}
Then the derived pushforward
$R\pi_{g*}\bar{B}^{(g)}_{\mathrm{flat}}(\cA)$ is a perfect complex on
$\overline{\mathcal{M}}_g$.
\end{lemma}

\begin{proof}
The argument has three steps.

\emph{Step~1: Uniform cohomological bound from the associated graded.}
The PBW filtration~$F_\bullet$ is flat over the base
$\overline{\mathcal{M}}_g$ (it is a bar-degree filtration,
hence defined fiber-by-fiber by the combinatorial bar grading);
hypothesis~(iii) supplies the finite truncations needed to apply the
filtered spectral sequence in a bounded cohomological range.
At the associated graded level the quantum correction terms in
$\Dg{g}$ sit in positive filtration degree, so
$\operatorname{gr}_F$ carries only the genus-$0$
differential~$\dzero$. By Koszul acyclicity~(i), the
cohomology of $\operatorname{gr}_F$ is concentrated in total
degree~$0$, giving a uniform vanishing range
$H^q(\operatorname{gr}_F\bar{B}^{(g)}_{\mathrm{flat}}(\cA)\big|_\Sigma) = 0$
for $q \neq 0$ at every point of the base.

\emph{Step~2: Transfer to the filtered complex.}
The bar-degree filtration spectral sequence
$E_1^{p,q} =
H^q(\operatorname{gr}_p\bar{B}^{(g)}_{\mathrm{flat}}(\cA)\big|_\Sigma)$
has $E_1^{p,q} = 0$ for $q \neq 0$ by Step~1. All higher
differentials $d_r$ ($r \geq 2$) therefore vanish (they map into
or out of the zero row), and the spectral sequence collapses at
$E_2$ with $E_\infty^{p,0} = E_2^{p,0}$. In particular, the
full flat fiber complex $\bar{B}^{(g)}_{\mathrm{flat}}(\cA)\big|_\Sigma$ has
cohomology concentrated in degree~$0$ for every~$\Sigma$.
The uniform vanishing range $q \neq 0$ is independent
of~$[\Sigma]$.

\emph{Step~3: Cohomology and base change.}
The family
$\pi_g\colon \mathcal{C}_g(\cA) \to \overline{\mathcal{M}}_g$
is a proper morphism of finite type, and hypothesis~(iii) makes the
strict flat family pseudo-coherent and uniformly bounded. By
Steps~1--2, the fiber
cohomology sheaves
$\mathcal{H}^q(R\pi_{g*}\bar{B}^{(g)}_{\mathrm{flat}}(\cA))$ vanish for
$q \neq 0$ uniformly over the base, and the surviving
degree-$0$ cohomology sheaf has
finite-dimensional fibers by hypothesis~(ii). The standard
cohomology and base change theorem
(EGA~III, Th\'eor\`eme~7.7.5)
then gives coherent cohomology. Since
$\overline{\mathcal{M}}_g$ is smooth, bounded coherent complexes are
perfect, so $R\pi_{g*}\bar{B}^{(g)}_{\mathrm{flat}}(\cA)$ is perfect.
\end{proof}

\begin{theorem}[Fiber--center identification \textup{(Theorem~\texorpdfstring{$\mathrm{C}_0$}{C0})}; \ClaimStatusConditional]
\label{thm:fiber-center-identification}
\index{fiber--center identification|textbf}
\textup{[}Regime: curved-central
\textup{(}Convention~\textup{\ref{conv:regime-tags})].}

\smallskip\noindent
Assume $\cA$ carries a modular pre-Koszul datum
\textup{(}Definition~\textup{\ref{def:modular-koszul-chiral})}.
Write
\[
\barB^{(g)}_{\mathrm{curv}}(\cA)
\;:=\;
\bigl(\barB^{(g)}(\cA), \dfib\bigr),
\qquad
\barB^{(g)}_{\mathrm{flat}}(\cA)
\;:=\;
\bigl(\barB^{(g)}(\cA), \Dg{g}\bigr)
\]
for the curved fiberwise and strict flat genus-$g$ bar models.
Then:
\begin{enumerate}[label=\textup{(\roman*)}]
\item \emph{Coderived form.}\;
 The curved fiber model
 $\barB^{(g)}_{\mathrm{curv}}(\cA)$ determines a well-defined
 coderived object. Proposition~\textup{\ref{prop:gauss-manin-uncurving-chain}}
 supplies its strict flat comparison representative
 $\barB^{(g)}_{\mathrm{flat}}(\cA)$.

\item \emph{Ordinary-derived realization on the flat perfect locus.}\;
 If $R\pi_{g*}\barB^{(g)}_{\mathrm{flat}}(\cA)$ is perfect over
 $\overline{\mathcal{M}}_g$
 \textup{(}Lemma~\textup{\ref{lem:perfectness-criterion})},
 then
 \begin{equation}\label{eq:fiber-center}
 \mathcal{H}^q\!\bigl(
 R\pi_{g*}\bar{B}^{(g)}_{\mathrm{flat}}(\cA)\bigr) = 0
 \quad\text{for } q \neq 0,
 \qquad
 \mathcal{H}^0\!\bigl(
 R\pi_{g*}\bar{B}^{(g)}_{\mathrm{flat}}(\cA)\bigr)
 \;\cong\; \mathcal{Z}_{\cA}
 \end{equation}
 as sheaves on $\overline{\mathcal{M}}_g$, where
 $\mathcal{Z}_{\cA}$ is the center local system.

\item \emph{Zero-curvature recovery.}\;
 If $\kappa(\cA)=0$, then $\dfib^{\,2}=0$, so the curved fiber model
 is itself an ordinary complex. In that case the conclusion of
 \textup{(ii)} may be read directly on
 $\barB^{(g)}_{\mathrm{curv}}(\cA)$.
\end{enumerate}
In particular, $\mathrm{C}_0$ holds unconditionally in coderived
form, and its ordinary-cohomology manifestation holds on the flat
perfect locus, with the original fiberwise reading recovered when
$\kappa(\cA)=0$.
\end{theorem}

\begin{proof}
The proof proceeds in four steps.

\emph{Step~1: Separate the curved and strict models.}
By Convention~\ref{conv:higher-genus-differentials}, the fiberwise
differential satisfies
$\dfib^{\,2} = \kappa(\cA)\cdot\omega_g$, so
$\barB^{(g)}_{\mathrm{curv}}(\cA)$ is a curved object and ordinary
cohomology is not the invariant when $\kappa(\cA)\neq 0$.
Proposition~\ref{prop:gauss-manin-uncurving-chain} supplies the strict
comparison model
$\barB^{(g)}_{\mathrm{flat}}(\cA)$ with $\Dg{g}^{\,2}=0$.
This proves \textup{(i)} and shows that the ordinary-derived
calculation must be performed on the flat representative.

\emph{Step~2: Filter the strict flat representative.}
Filter the flat fiber bar complex
$C^{(g)}_{X,\mathrm{flat}}(\cA)\!\mid_\Sigma$ by bar degree, writing
\[
F_p := \bigoplus_{n \leq p} \bar{B}^{(g),n}(\cA).
\]

\emph{Step~3: Identify the associated graded and collapse.}
At the associated graded level, the quantum corrections in
$\Dg{g}$ lie in positive filtration degree, so
$\operatorname{gr}_F$ sees only the genus-$0$
collision differential~$\dzero$.
By genus-$0$ Koszulity (axiom~MK1), the associated graded is
the classical Koszul/Ext complex of
$\operatorname{gr}_F \cA$, which has cohomology concentrated
in total degree~$0$.

The bar-degree filtration spectral sequence
\[
E_1^{p,q} = H^q\bigl(\operatorname{gr}_p
C^{(g)}_{X,\mathrm{flat}}(\cA)\big|_\Sigma\bigr)
\]
has $E_1^{p,q} = 0$ for $q \neq 0$ by Step~2
(the associated graded sees only the
genus-$0$ collision differential, which is Koszul-acyclic).
Therefore all differentials $d_r$ with $r \geq 2$ vanish
(their target lies in a row $q' \leq -1$, which is zero),
and $d_1$ acts within the $q = 0$ row, so
$E_\infty^{p,q} = E_2^{p,q}$. In particular,
$H^q(C^{(g)}_{X,\mathrm{flat}}(\cA)\!\mid_\Sigma) = 0$
for $q \neq 0$: the strict flat fiber complex has cohomology
concentrated in bar-differential degree~$0$.
On the Koszul locus, Theorem~\ref{thm:cobar-resolution-scoped}
identifies the strict flat bar family as a bar resolution of~$\cA$,
and the flat-side center computation of
Theorem~\ref{thm:obstruction-quantum} identifies the surviving
degree-$0$ sheaf with the center local system~$\mathcal{Z}_{\cA}$.

\emph{Step~4: Base change and the $\kappa=0$ recovery.}
Apply cohomology and base change to the family
$\pi_g\colon \mathcal{C}_g(\cA) \to \overline{\mathcal{M}}_g$.
The fiber-level concentration from Step~3 gives
$\mathcal{H}^q(R\pi_{g*}\bar{B}^{(g)}_{\mathrm{flat}}(\cA)) = 0$
for $q \neq 0$.
Perfectness of $R\pi_{g*}\bar{B}^{(g)}_{\mathrm{flat}}(\cA)$
(Lemma~\ref{lem:perfectness-criterion}) ensures base change applies
uniformly over $\overline{\mathcal{M}}_g$.
The surviving $\mathcal{H}^0$ is identified with
$\mathcal{Z}_{\cA}$ by the fiber-level identification
from Step~3, which holds uniformly over the base.
If $\kappa(\cA)=0$, then $\dfib^{\,2}=0$ and the curved model is
already an ordinary complex; the same conclusion can therefore be
read directly on $\barB^{(g)}_{\mathrm{curv}}(\cA)$.
\end{proof}

\begin{remark}[Theorem~C decomposition]\label{rem:theorem-C-decomposition}
\index{Theorem C!decomposition}
Theorem~C has three named pieces: $\mathrm{C}_0$ (fiber-center
identification, this theorem) produces the ambient complex;
$\mathrm{C}_1$ (Theorem~\ref{thm:quantum-complementarity-main})
gives the Verdier eigenspace decomposition; the Lagrangian statement
inside $\mathrm{C}_1$ requires the non-degenerate anti-invariant
Verdier pairing. $\mathrm{C}_2$
(Theorem~\ref{thm:shifted-symplectic-complementarity}) is the
conditional bar-side BV / shifted-symplectic upgrade.
The chain-level mechanism for~$\mathrm{C}_1$ is the Verdier involution
$\sigma$ exchanging the two eigenspaces; at the scalar level
the complementarity formulas below are consequences of the
$\mathrm{C}_1$ decomposition together with Theorem~D, and not a
separate $\mathrm{C}_2$ label:
$\kappa(\cA)+\kappa(\cA^!)=0$ for Kac--Moody and free-field pairs,
while $\kappa(\cA)+\kappa(\cA^!)=\varrho\cdot K\neq 0$ for
$\mathcal{W}$-algebras (Theorem~\ref{thm:complementarity-root-datum};
Table~\ref{tab:complementarity-landscape}).
\end{remark}

The Heisenberg complementarity
of~\S\ref{sec:frame-complementarity} generalizes to a Verdier
eigenspace decomposition, with a Lagrangian polarization under the
pairing hypotheses below:

\begin{theorem}[Quantum complementarity: Verdier eigenspaces and
Lagrangian upgrade; \ClaimStatusConditional]
\label{thm:quantum-complementarity-main}
\label{thm:quantum-complementarity}
\index{deformation-obstruction complementarity|textbf}
\index{Lagrangian!complementarity|textbf}
\textup{[}Regime: curved-central
\textup{(}Convention~\textup{\ref{conv:regime-tags})].}

Let $(\mathcal{A}, \mathcal{A}^!)$ be a chiral Koszul pair on a smooth projective
curve $X$ over $\mathbb{C}$, with $\mathcal{A}$ a sheaf of chiral algebras in
the sense of Beilinson--Drinfeld \cite[Chapter 3]{BD04} and $\mathcal{A}^!$
its Koszul dual (Theorem~\ref{thm:chiral-koszul-duality}).
Write $\mathbf{C}_g(\cA)$, $\mathbf{Q}_g(\cA)$, $\mathbf{Q}_g(\cA^!)$
for the ambient complex and homotopy eigenspaces of
Definition~\textup{\ref{def:complementarity-complexes}}.
Fix a genus $g\geq0$ for which the ordinary flat representative of
\textup{(C0)} is available, the center-comparison datum
$\iota_Z$ is fixed, and the Verdier--Koszul involution
$\sigma$ is represented by a cochain involution on
$\mathbf{C}_g(\cA)$.

\smallskip\noindent
\textbf{H-level \textup{(}homotopy\textup{)}.}
\begin{enumerate}[label=\textup{(\roman*)}]
\item \emph{Homotopy eigenspace decomposition.}
 By Lemma~\textup{\ref{lem:involution-splitting}(a)},
 \begin{equation}\label{eq:complementarity-decomp}
 \mathbf{C}_g(\cA) \;\simeq\;
 \mathbf{Q}_g(\cA) \;\oplus\; \mathbf{Q}_g(\cA^!),
 \end{equation}
 functorially in~$(\cA,\cA^!)$.

\item \emph{Perfect duality for $g \geq 1$.}
 Assume $g \geq 1$ and assume that Verdier duality induces a
 non-degenerate cochain-level pairing
 $\langle -,- \rangle_{\mathbb{D}}$ on~$\mathbf{C}_g(\cA)$
 of cohomological degree $-(3g{-}3)$, anti-symmetric under the
 Verdier involution~$\sigma$. Under this pairing, each summand in
 \eqref{eq:complementarity-decomp} is Lagrangian, and
 $\langle -, - \rangle_{\mathbb{D}}$ restricts to a
 non-degenerate pairing
 $\mathbf{Q}_g(\cA) \otimes \mathbf{Q}_g(\cA^!) \to \mathbb{C}[-{(3g{-}3)}]$,
 so $\mathbf{Q}_g(\cA) \simeq \mathbf{Q}_g(\cA^!)^\vee[-(3g{-}3)]$.

\item \emph{Functoriality.}
 The decomposition is natural in morphisms of Koszul pairs
 and compatible with the conformal weight and
 cohomological degree gradings.
\end{enumerate}

\smallskip\noindent
\textbf{S-level \textup{(}cohomological shadow\textup{)}.}
Setting $\mathcal{H}_g(\cA) := H^*(\overline{\mathcal{M}}_g,
\mathcal{Z}(\cA))$, Lemma~\textup{\ref{lem:involution-splitting}(b)}
gives:
\begin{equation}\label{eq:verdier-symplectic-pairing}
\mathcal{H}_g(\cA) \;=\;
Q_g(\cA) \;\oplus\; Q_g(\cA^!),
\end{equation}
where the cohomological shadows
$Q_g(\cA) = \ker(\sigma - \mathrm{id})$ and
$Q_g(\cA^!) = \ker(\sigma + \mathrm{id})$
\textup{(}Definition~\textup{\ref{def:complementarity-complexes}}\textup{)}
are the $\pm 1$ eigenspaces of~$\sigma$.
If the pairing hypothesis in the H-level clause holds, then for
$g \geq 1$ they are Lagrangians for the induced pairing
on~$\mathcal{H}_g$, and
\begin{equation}\label{eq:quantum-duality}
Q_g(\cA) \cong Q_g(\cA^!)^\vee.
\end{equation}
At genus $0$, the unique point class is $\sigma$-fixed, so
\[
Q_0(\cA)=H^0(\overline{\mathcal{M}}_0,\mathcal{Z}(\cA)),
\qquad
Q_0(\cA^!)=0.
\]
\end{theorem}

\begin{proof}
The ordinary flat representative supplied by~\textup{(C0)} places the
ambient object in the category of cochain complexes, and the hypothesis
on~$\sigma$ says that this representative carries a cochain involution.
The idempotents
\[
 p^\pm=\tfrac12(\mathrm{id}\pm\sigma)
\]
are therefore cochain maps. Lemma~\ref{lem:involution-splitting}(a)
gives the direct-sum decomposition
\[
\mathbf{C}_g(\cA)
=
\operatorname{im}(p^+) \oplus \operatorname{im}(p^-)
=
\mathbf{Q}_g(\cA)\oplus\mathbf{Q}_g(\cA^!),
\]
which is the H-level assertion. Applying cohomology gives
\[
H^*(\mathbf{C}_g(\cA))
=
\ker(\sigma-\mathrm{id})\oplus\ker(\sigma+\mathrm{id})
=
Q_g(\cA)\oplus Q_g(\cA^!)
\]
by Lemma~\ref{lem:involution-splitting}(b), and
$H^*(\mathbf{C}_g(\cA))=
H^*(\overline{\mathcal{M}}_g,\mathcal{Z}(\cA))$ by definition of the
ambient complex.

For $g\geq1$, assume the Verdier pairing is non-degenerate and
anti-invariant. Lemma~\ref{lem:involution-splitting}(c) makes the two
summands isotropic and maximal; the non-degenerate cross-pairing has
degree $-(3g{-}3)$, hence
$\mathbf{Q}_g(\cA)\simeq
\mathbf{Q}_g(\cA^!)^\vee[-(3g{-}3)]$ and, on cohomology,
$Q_g(\cA)\cong Q_g(\cA^!)^\vee$. At $g=0$ the moduli space is a point
in this ambient convention and the Verdier orientation class is fixed;
the projector $p^+$ is the identity on the point class and $p^-$
vanishes. Naturality follows because morphisms of Koszul pairs commute
with the represented Verdier involution and hence with the projectors.
\end{proof}

\begin{remark}\label{rem:lagrangian-decomp}
What one algebra sees as deformation, its dual sees as obstruction.
\end{remark}

\begin{remark}[External independence of Theorem~C]
\label{rem:theorem-c-mok-independence}
Theorem~C depends on three inputs: the genus
filtration spectral sequence (standard), Poincar\'e--Verdier
duality on FM compactifications (standard algebraic
geometry, no external preprint), and the involution-splitting
lemma (linear algebra in characteristic~$0$).
In particular, Theorem~C does \emph{not} depend on the
ambient-level $D^2 = 0$
(Theorem~\textup{\ref{thm:ambient-d-squared-zero}},
which uses \textup{[}Mok25\textup{]}) or on the full
bar-intrinsic MC element $\Theta_\cA$
(Theorem~\textup{\ref{thm:mc2-bar-intrinsic}}).
The only curvature input is the scalar $\kappa(\cA)$,
which comes from the convolution-level
$D^2 = 0$ (unconditional).
\end{remark}

\begin{remark}[Elementary model presentation;
Convention~\ref{conv:proof-architecture}]
\label{rem:theorem-c-model}
\label{rem:complementarity-model}% backward-compatible label
\index{complementarity!model presentation}
\emph{Step~B} (M-level): The homotopy eigenspace
$\mathbf{Q}_g(\cA)=\operatorname{fib}(\sigma-\mathrm{id})$ is modeled by
$\operatorname{im}(p^+)$ with $p^+=\tfrac{1}{2}(\mathrm{id}+\sigma)$
on $C_g=R\Gamma(\overline{\mathcal{M}}_g,\mathcal{Z}(\cA))$.
The Lagrangian property is $\langle p^+x,p^+y\rangle=0$
(Lemma~\ref{lem:involution-splitting}(c)).
Taking cohomology gives the S-level decomposition
$\mathcal{H}_g=Q_g(\cA)\oplus Q_g(\cA^!)$; model independence
(Proposition~\ref{prop:model-independence}) lifts to H-level.
For the Heisenberg at genus~$1$: $Q_1=H^0$, $Q_1^!=H^2$.
\end{remark}

\begin{remark}[Comparison with literature]
Beilinson--Drinfeld~\cite[Chapter~4]{BD04} proved the $g=0$ case;
we extend to $g\geq 1$.
Gui--Li--Zeng~\cite{GLZ22} (curved Koszul duality for non-quadratic
operads) is applied here to the chiral setting.
Costello--Gwilliam~\cite{CG17} (factorization homology for TFTs)
is the topological analogue of our holomorphic construction.
Arakawa~\cite{Ara12} ($\mathcal{W}$-algebra representation theory)
is explained by our complementarity at critical level.
\end{remark}

\begin{remark}[Homotopy-native formulation]\label{rem:homotopy-native-c}
The complex $\mathbf{C}_g(\cA)=R\Gamma(\overline{\mathcal{M}}_g,
\mathcal{Z}(\cA))$ carries a $-(3g{-}3)$-shifted symplectic structure
(PTVV~\cite{PTVV13}; Proposition~\ref{prop:ptvv-lagrangian})
under the perfectness and nondegeneracy hypotheses recorded there.
The decomposition $\mathbf{C}_g\simeq\mathbf{Q}_g(\cA)\oplus
\mathbf{Q}_g(\cA^!)$ is a Lagrangian fibration with $\sigma$ acting
as anti-symplectomorphism only in that paired regime; the tangent
and obstruction complexes are then Serre dual. The S-level
eigenspace decomposition is the decategorification.
\end{remark}

\begin{remark}[PTVV comparison]
The comparison with
Pantev--To\"en--Vaqui\'e--Vezzosi~\cite{PTVV13}
has four parts.

First, the ambient formal moduli problem
$\mathcal{M}_{\mathrm{comp}}(\cA)$ of
Theorem~\ref{thm:ambient-complementarity-fmp} fits the standard
shifted-symplectic pattern: the degree-$(-1)$ invariant pairing on its
tangent dg~Lie algebra produces a $(-1)$-shifted symplectic form on the
formal moduli problem. This is the correct PTVV comparison for the
bar--cobar moduli space; on the linear side it is the constant-form
model of \cite[Example~1.4]{PTVV13}.

Second, the scalar complementarity identity
$\kappa(\cA)+\kappa(\cA^!)=K$ is not a direct consequence of PTVV.
PTVV controls shifted symplectic forms, isotropic structures, and
Lagrangian maps
\textup{(}Definition~\textup{2.8}, Theorem~\textup{2.9}\textup{)};
it does not by itself produce the scalar conductor~$K$.
In this chapter the scalar formula is obtained only after the
Verdier eigenspace decomposition of~$\mathrm{C}_1$ is projected to the
scalar lane by Theorem~D and the family computations of
Table~\ref{tab:complementarity-landscape}.

Third, the Lagrangian content of~$\mathrm{C}_1$ is consistent with
PTVV, but the precise comparison is polarization, not derived
intersection. Proposition~\ref{prop:ptvv-lagrangian} shows that
$\mathbf{Q}_g(\cA)$ and $\mathbf{Q}_g(\cA^!)$ satisfy the PTVV
Lagrangian conditions inside the ambient
$(-(3g{-}3))$-shifted space $\mathbf{C}_g(\cA)$.
The derived-intersection theorem of PTVV
\textup{(}Theorem~\textup{2.9}\textup{)} enters the
cotangent-normal-form package
\textup{(}Theorem~\textup{\ref{thm:shifted-cotangent-normal-form}}\textup{)},
not as the statement of~$\mathrm{C}_1$ itself.

Fourth, PTVV transgression
\textup{(}Theorem~\textup{2.5}\textup{)} and the AKSZ refinements of
Calaque--Haugseng--Scheimbauer~\cite{CHS25} and
Calaque--Safronov~\cite{CS24} are compatible with the shift conventions
of our fiberwise construction, but they do not justify the shortcut
``$(-1)$ on $\mathcal{M}_{\mathrm{comp}}$ integrates to $0$ on the
genus-$g$ field space.'' The correct comparison is that a
$0$-shifted pairing on the center local system, pushed forward along a
base of dimension $3g{-}3$, gives the
$(-(3g{-}3))$-shifted Verdier pairing of
Proposition~\ref{prop:ptvv-lagrangian}. The present proof reaches that
pairing through fiber--center identification and Verdier duality; the
AKSZ route is an alternative geometric packaging of the same
orientation data.
\end{remark}

\subsection{Verdier splitting mechanism}

Theorem~C is a statement about one represented involution on one
ambient complex. The relative configuration-space machinery supplies
that complex and proves that the cohomology sheaf is the center
local system on the flat perfect locus. Verdier duality supplies the
involution~$\sigma$ and, when finite-type duality is available, the
anti-invariant pairing. The direct sum itself is then the idempotent
splitting of $\sigma$ in characteristic~$0$; no dimension count,
scalar conductor identity, or bar-side BV structure is part of the
logical proof of the eigenspace decomposition.

\subsection{Relative Leray spectral sequence}

\begin{proof}[Construction of the relative Leray package]

\emph{Relative genus summand.}

\begin{lemma}[Genus filtration; \ClaimStatusProvedHere]
\label{lem:genus-filtration}
Let
\[
\pi_{g,n}\colon \overline{\operatorname{Conf}}_{g,n}(X)
\longrightarrow \overline{\mathcal{M}}_g
\]
be the relative Fulton--MacPherson compactification of $n$ ordered
points on the universal stable curve. The strict flat genus-$g$ bar
summand is the derived global section complex
\[
\barB^{(g)}_{\mathrm{flat}}(\cA)
:=
\bigoplus_{n\geq0}
R\Gamma\!\left(
\overline{\mathcal{M}}_g,\,
R\pi_{g,n*}\bigl(\cA^{\boxtimes n}\otimes\Omega^\bullet_{\log},
\Dg{g}\bigr)\right).
\]
For a fixed curve there is no filtration by varying genus. On the
modular bar object
$\barB^{\mathrm{mod}}_{\mathrm{flat}}(\cA)
=\prod_{h\geq0}\barB^{(h)}_{\mathrm{flat}}(\cA)$ there is the
external genus truncation
\[
F^{\leq G}\barB^{\mathrm{mod}}_{\mathrm{flat}}(\cA)
:=
\prod_{0\leq h\leq G}\barB^{(h)}_{\mathrm{flat}}(\cA).
\]
The differential, coproduct, and stable-graph clutching maps preserve
arithmetic genus, so this truncation is by subcomplexes and coalgebra
subobjects.
\end{lemma}

\begin{proof}[Proof of Lemma~\ref{lem:genus-filtration}]
For fixed~$g$, the relative compactification
$\overline{\operatorname{Conf}}_{g,n}(X)$ is over
$\overline{\mathcal{M}}_g$; its fiber at a stable curve~$\Sigma$ is
the Fulton--MacPherson compactification
$\overline{\operatorname{Conf}}_{n}(\Sigma)$. This defines the displayed relative
bar summand. A collision residue contracts a subset of marked points
inside the same stable curve and therefore preserves the arithmetic
genus of the curve. A separating boundary clutching replaces one
vertex of genus~$g$ by vertices of genera $g_1,g_2$ with
$g_1+g_2+h^1(\Gamma)=g$; a non-separating clutching lowers a vertex
genus by one and raises the graph loop number by one. In both cases
the arithmetic genus is unchanged. Deconcatenation of bar words
changes only the point partition. Hence the differential, coproduct,
and clutching maps preserve $F^{\leq G}$.
\end{proof}

\begin{remark}[Genus grading and loop degree]
In the physics normalization the external genus grading is the
loop-degree grading:
\begin{equation}
Z = Z^{(0)} + \hbar Z^{(1)} + \hbar^2 Z^{(2)} + \cdots
\end{equation}
where $Z^{(g)}$ is the contribution of the genus-$g$ relative bar
summand. This is a grading of the modular object, not a filtration of
the bar complex on one fixed curve.
\end{remark}

\emph{Relative Leray spectral sequence.}

\begin{theorem}[Spectral sequence for quantum corrections; \ClaimStatusProvedHere]
\label{thm:ss-quantum}
The relative Leray filtration on the strict flat comparison family
$\barB^{(g)}_{\mathrm{flat}}(\mathcal{A})
= (\barB^{(g)}(\mathcal{A}), \Dg{g})$
over $\overline{\mathcal{M}}_g$ induces a spectral sequence:
\begin{equation}
E_1^{p,q,g} = H^q\left(\barB^p_{g,\mathrm{flat}}(\mathcal{A})\right)
\Longrightarrow
H^{p+q}\left(\barB^{(g)}_{\mathrm{flat}}(\mathcal{A})\right)
\end{equation}
where:
\begin{itemize}
\item $p$ = configuration space degree (number of points)
\item $q$ = form degree (dimension of logarithmic forms)
\item $g$ = genus degree
\item $\barB^p_{g,\mathrm{flat}}(\mathcal{A})$ = the $p$-point piece of the
 flat genus-$g$ bar complex with square-zero differential~$\Dg{g}$
\item the base differential is the de~Rham/Cech differential on
 $\overline{\mathcal{M}}_g$ coupled to the Gauss--Manin connection
\end{itemize}
The curved fiber operator~$\dfib$ remains only as a coderived
comparison surface when $\kappa(\mathcal{A}) \neq 0$; ordinary
cohomology is taken after passage to the flat representative.

The $E_2$ page is:
\begin{equation}
E_2^{p,q,g} = H^p\left(\overline{\mathcal{M}}_g, \mathcal{H}^q_{\mathrm{flat}}(
\mathcal{A})\right)
\end{equation}
where $\mathcal{H}^q_{\mathrm{flat}}(\mathcal{A})$ is the sheaf of
flat fiber cohomologies.
\end{theorem}

\begin{proof}[Proof of Theorem~\ref{thm:ss-quantum}]
For each genus~$g$, the Leray spectral sequence for the fibration
below is applied to the strict flat comparison family
$\barB^{(g)}_{\mathrm{flat}}(\mathcal{A})$, hence converges to
$H^{p+q}(\barB^{(g)}_{\mathrm{flat}}(\mathcal{A}))$.
The external genus truncation of
Lemma~\ref{lem:genus-filtration} assembles the fixed-genus spectral
sequences into the completed modular bar object, but the displayed
convergence statement is genuswise.

The Leray spectral sequence is applied to the fibration:
\begin{equation}
\begin{tikzcd}
\overline{\operatorname{Conf}}_{g,n}(X) \arrow[d, "\pi_{g,n}"] \\
\overline{\mathcal{M}}_g
\end{tikzcd}
\end{equation}

\emph{$E_1$~page.}
By definition, $E_1^{p,q,g}$ is the cohomology
of the fiber complex.
The fiber over
$[(\Sigma_g;\allowbreak p_1, \ldots, p_n)]$ is:
\begin{equation}
\barB^p_{g,\mathrm{flat}}(\mathcal{A})\big|_{[\Sigma_g]}
= \Gamma(\overline{C}_p(\Sigma_g), \mathcal{A}^{\boxtimes p}
\otimes \Omega^*_{\log})
\end{equation}
\[
\text{with differential } \Dg{g}^{(\Sigma_g)}.
\]
By Convention~\ref{conv:higher-genus-differentials},
$\Dg{g}^{\,2} = 0$ on this flat comparison complex. The Arnold
relations still govern the local collision term in the associated
graded, but they do not imply $\dfib^{\,2} = 0$ when
$\kappa(\mathcal{A}) \neq 0$; the curved fiber operator remains on the
coderived side.
\begin{equation}
E_1^{p,q,g}
= H^q\bigl(\barB^p_{g,\mathrm{flat}}(\mathcal{A})\big|_{[\Sigma_g]}\bigr)
\end{equation}

\emph{$d_1$ differential}: This is induced by the differential on $\overline{
\mathcal{M}}_g$. It measures how the fiber cohomology varies as we move in moduli space.

\emph{$E_2$ page}: After taking cohomology with respect to $d_1$, we obtain:
\begin{equation}
E_2^{p,q,g} = H^p(\overline{\mathcal{M}}_g, \mathcal{H}^q_{\mathrm{flat}})
\end{equation}
where $\mathcal{H}^q_{\mathrm{flat}}$ is the sheaf on
$\overline{\mathcal{M}}_g$ whose stalk at $[(\Sigma_g; \vec{p})]$ is
$H^q\bigl(\barB^p_{g,\mathrm{flat}}(\mathcal{A})\big|_{[\Sigma_g]}\bigr)$.

This sheaf is \emph{locally constant} away from boundary strata, by the local
triviality of the fibration. On boundary strata, it has monodromy captured by the
\emph{Picard--Lefschetz formula}.
\end{proof}

\begin{lemma}[Convergence of the quantum spectral sequence;
\ClaimStatusConditional]
\label{lem:quantum-ss-convergence}
\index{spectral sequence!convergence}
The spectral sequence of
Theorem~\textup{\ref{thm:ss-quantum}} converges genuswise to
$H^{p+q}(\barB^{(g)}_{\mathrm{flat}}(\cA))$. Since
$\overline{\mathcal{M}}_g$ has finite cohomological dimension, it
stabilizes at a finite page; in the crude cohomological-dimension
bound one may take $E_r=E_\infty$ for $r>6g-5$.
\end{lemma}

\begin{proof}
Three conditions must be verified.

\emph{(1) Finite cohomological dimension of the base.}
The Deligne--Mumford--Knudsen stack
$\overline{\mathcal{M}}_g$ is a smooth proper
Deligne--Mumford stack of dimension $3g - 3$. For any
constructible sheaf $\mathcal{F}$ on
$\overline{\mathcal{M}}_g$,
$H^i(\overline{\mathcal{M}}_g, \mathcal{F}) = 0$ for
$i > 2(3g-3) = 6g - 6$
(Artin vanishing for proper DM stacks,
\cite[\S4.1]{Olsson16}).

\emph{(2) Constructibility of the fiber sheaves.}
The fiber cohomology sheaf
$\mathcal{H}^q_{\mathrm{flat}}(\cA)$ on
$\overline{\mathcal{M}}_g$ is constructible with respect to
the stratification by topological type of stable curve.
On the smooth locus $\mathcal{M}_g$, it is a local system
by Lemma~\ref{lem:fiber-cohomology-center} (constancy of
the center $Z(\cA)$ under variation of complex structure).
On each boundary stratum $D_\Gamma$
(indexed by stable graphs~$\Gamma$), the monodromy
is quasi-unipotent: the Picard--Lefschetz formula gives
monodromy operators $T_e = \mathrm{id} + N_e$ where $N_e$
is nilpotent (one Jordan block per node
$e \in E(\Gamma)$). A local system with quasi-unipotent
monodromy along a normal crossing divisor is constructible
(\cite[\S8.1]{KS90}).

\emph{(3) Finite Leray filtration.}
For fixed~$g$, the Leray filtration on
$R\Gamma(\overline{\mathcal{M}}_g,R\pi_{g,*}(-))$ is finite in each
total degree because the base has finite cohomological dimension and
the fiber complex is bounded in the flat perfect regime. The external
genus truncation of Lemma~\ref{lem:genus-filtration} is complete and
separated on the modular product, but convergence of the displayed
spectral sequence is the fixed-genus Leray convergence.

\emph{Convergence.} By (1)--(3), the spectral sequence
satisfies the hypotheses of the filtered convergence theorem
for finite filtrations with constructible graded pieces
(Grothendieck~\cite[\S II.4.17]{Tohoku}). The $d_r$
differential shifts the Leray degree by~$r$ and the complementary
degree by~$-r+1$. Since $E_r^{p,q,g}=0$ for $p>6g-6$ by (1), every
$d_r$ with $r>6g-5$ has either zero source or zero target. Thus the
genus-$g$ Leray spectral sequence stabilizes by that page.
\end{proof}

\emph{Leray corrections and Verdier projection.}

\begin{lemma}[Leray correction object and Verdier projection; \ClaimStatusConditional]
\label{lem:quantum-from-ss}
The relative Leray spectral sequence defines a genus-$g$ correction
graded object
\begin{equation}
G_g(\mathcal{A})
:= E_\infty^{*,*,g}
= \bigoplus_{p+q=*}
\operatorname{gr}^{p}_{\mathrm{Leray}}
H^{p+q}\bigl(\barB^{(g)}_{\mathrm{flat}}(\mathcal{A})\bigr).
\end{equation}
After the flat fiber--center comparison and the represented Verdier
involution~$\sigma$ are in place, the Theorem~C space
$Q_g(\mathcal{A})$ is the $+1$ Verdier projection
$p^+H^*(\mathbf{C}_g(\mathcal{A}))$, not the definition of the
Leray graded object. The comparison map
$G_g(\mathcal{A})\to H^*(\mathbf{C}_g(\mathcal{A}))$ lands in
$Q_g(\mathcal{A})$ precisely when the class is represented by the
$j_*$-extended $\mathcal{A}$-bar sector.
\end{lemma}

\begin{proof}[Proof of Lemma~\ref{lem:quantum-from-ss}]
By definition of spectral sequences, $E_\infty$ is the associated graded of the
filtered cohomology:
\begin{equation}
E_\infty^{p,q,g} \cong
\frac{L^p H^{p+q}(\barB^{(g)}_{\mathrm{flat}}(\mathcal{A}))}
{L^{p+1} H^{p+q}(\barB^{(g)}_{\mathrm{flat}}(\mathcal{A}))}
\end{equation}

The genus-$g$ Leray correction object records those cohomology
classes that arise from genus-$g$ relative configuration data:
\begin{equation}
G_g(\mathcal{A}) :=
\operatorname{gr}^{\bullet}_{\mathrm{Leray}}
H^*(\barB^{(g)}_{\mathrm{flat}}(\mathcal{A}))
= E_\infty^{*,*,g}.
\end{equation}
Definition~\ref{def:complementarity-complexes} reserves the notation
$Q_g(\mathcal{A})$ for the Verdier eigenspace
$\ker(\sigma-\mathrm{id})$. The comparison between these two objects
uses the fiber--center identification and
Lemma~\ref{lem:eigenspace-decomposition-complete}: $j_*$-extended
bar cochains from~$\mathcal{A}$ project to the $+1$ eigenspace, while
$j_!$-extended bar cochains from~$\mathcal{A}^!$ project to the
$-1$ eigenspace.

\emph{Explicit description.} An element of $G_g(\mathcal{A})$ is
represented by a $\Dg{g}$-closed relative class
$\omega \in \barB^{(g)}_{\mathrm{flat}}(\mathcal{A})$ modulo the
next Leray filtration step. Its image in $Q_g(\mathcal{A})$ is
$p^+[\omega]$ after the center comparison. On the curved fiber model
one retains the corresponding coderived class whenever
$\kappa(\mathcal{A}) \neq 0$.

\emph{Example.} For the Heisenberg algebra at $g=1$, the $p^+$
projection of the genus-$1$ correction class is
\begin{equation}
Q_1(\mathcal{H}_\kappa)
= \operatorname{span}\{\kappa \cdot \lambda_1\} \cong \mathbb{C}
\end{equation}
Here $\kappa$ is the level parameter and $\lambda_1 = c_1(\mathbb{E})$ is the Hodge class;
the obstruction class $\kappa \cdot \lambda_1$ is a genus-$1$ quantum correction
that does not appear at genus~$0$.
\end{proof}

\emph{Flat fiber cohomology and center.}

\begin{lemma}[Flat fiber cohomology and center; \ClaimStatusConditional]
\label{lem:fiber-cohomology-center}
For a Koszul chiral algebra~$\mathcal{A}$
\textup{(}Definition~\textup{\ref{def:chiral-koszul-pair})},
the flat fiber cohomology sheaf over the smooth
locus~$\overline{\mathcal{M}}_g^{\mathrm{sm}}$ satisfies:
\begin{equation}
\mathcal{H}^*_{\mathrm{flat}}(\mathcal{A})
\big|_{\overline{\mathcal{M}}_g^{\mathrm{sm}}}
\cong Z(\mathcal{A}) \otimes \underline{\mathbb{C}}
\end{equation}
where $\underline{\mathbb{C}}$ is the constant sheaf.
\end{lemma}

\begin{proof}[Proof of Lemma~\ref{lem:fiber-cohomology-center}]
This is the ordinary-derived clause of
Theorem~\ref{thm:fiber-center-identification} restricted to the smooth
locus. On the strict flat comparison family
$\barB^{(g)}_{\mathrm{flat}}(\mathcal{A})$, the differential
$\Dg{g}$ squares to zero, the higher flat fiber cohomology vanishes,
and the surviving degree-$0$ sheaf is the center local system.
Hence
\[
\mathcal{H}^q_{\mathrm{flat}}(\mathcal{A}) = 0
\quad\text{for } q \neq 0,
\qquad
\mathcal{H}^0_{\mathrm{flat}}(\mathcal{A})
\cong Z(\mathcal{A}) \otimes \underline{\mathbb{C}}.
\]
When $\kappa(\mathcal{A}) \neq 0$, the curved fiber differential
$\dfib$ is retained only in coderived form; if $\kappa(\mathcal{A})=0$,
the curved and flat models coincide by
Theorem~\ref{thm:fiber-center-identification}(iii).
\end{proof}

The relative Leray spectral sequence, the correction object
$G_g(\mathcal{A})$, and the flat fiber cohomology sheaf are separated
from the Verdier eigenspaces $Q_g(\mathcal{A})$ and
$Q_g(\mathcal{A}^!)$.
\end{proof}

\subsection{Verdier duality on fibers}

\begin{proof}[Verdier duality package]

\emph{Poincaré--Verdier duality on configuration spaces.}

\begin{theorem}[Verdier duality for compactified configuration spaces; \ClaimStatusProvedHere]
\label{thm:verdier-duality-config-complete}
\index{Grothendieck--Verdier duality}
Let $X$ be a smooth projective curve of genus $g$. The Fulton--MacPherson
compactification $\overline{C}_n(X)$ satisfies Poincaré--Verdier duality:
\begin{equation}
\mathbb{D}: \mathcal{H}^k(\overline{C}_n(X)) \xrightarrow{\sim} \mathcal{H}^{d-k}(
\overline{C}_n(X))^\vee[d]
\end{equation}
where $d = \dim_{\mathbb{R}} \overline{C}_n(X) = 2n$ and $\mathbb{D}$ is the Verdier
dualizing functor.
\end{theorem}

\begin{proof}[Proof of Theorem~\ref{thm:verdier-duality-config-complete}]
\emph{Setup}: Recall from Section~\ref{sec:FM-compactification} that $\overline{C}_n(X)$
is constructed by iterated blow-ups along diagonal strata. The key properties are:
\begin{enumerate}
\item $\overline{C}_n(X)$ is a smooth complex manifold (real dimension $2n$)
\item The boundary $\partial \overline{C}_n(X) = \overline{C}_n(X) \setminus C_n(X)$
is a normal crossing divisor
\item The compactification is functorial in $X$ and natural with respect to the
symmetric group $\Sigma_n$
\end{enumerate}

\emph{Verdier duality}: For any smooth proper variety $Y$ over $\mathbb{C}$ with
normal crossing boundary, the Verdier dualizing complex is:
\begin{equation}
\mathbb{D}_Y \mathcal{F} = \mathcal{RH}om(\mathcal{F}, \omega_Y[\dim Y])
\end{equation}
where $\omega_Y$ is the dualizing sheaf (canonical bundle).

\emph{Application to $\overline{C}_n(X)$}: Since $\overline{C}_n(X)$ is smooth and
proper, we have:
\begin{equation}
\omega_{\overline{C}_n(X)} = K_{\overline{C}_n(X)} = \Omega^{n}_{\overline{C}_n(X)}
\end{equation}
Here $\Omega^n_{\overline{C}_n(X)} = \bigwedge^n \Omega^1_{\overline{C}_n(X)}$ denotes the holomorphic line bundle of top-degree holomorphic forms (the canonical bundle), not the sheaf of smooth $n$-forms.
Since $\overline{C}_n(X)$ is a smooth proper variety over~$\mathbb{C}$, the algebraic and analytic canonical bundles coincide by~GAGA.

The duality pairing is given by integration:
\begin{equation}
\langle \alpha, \beta \rangle = \int_{\overline{C}_n(X)} \alpha \wedge \beta
\end{equation}
for $\alpha \in H^k(\overline{C}_n(X))$ and $\beta \in H^{2n-k}(\overline{C}_n(X))$.

\emph{Perfect pairing}: By Poincaré duality for compact oriented manifolds:
\begin{equation}
H^k(\overline{C}_n(X)) \times H^{2n-k}(\overline{C}_n(X)) \xrightarrow{\wedge}
H^{2n}(\overline{C}_n(X)) \xrightarrow{\int} \mathbb{C}
\end{equation}
is a perfect pairing. This is the geometric incarnation of Verdier duality.

\emph{Logarithmic forms}: When we include logarithmic forms $\Omega^*_{\log}(
\overline{C}_n(X))$ (forms with logarithmic poles along $\partial \overline{C}_n(X)$),
the duality becomes:
\begin{equation}
\Omega^k_{\log}(\overline{C}_n(X)) \times \Omega^{2n-k}_{\log}(\overline{C}_n(X))
\to \mathbb{C}
\end{equation}
given by:
\begin{equation}
\langle \eta, \xi \rangle = \text{Res}_{\partial \overline{C}_n(X)} (\eta \wedge \xi)
\end{equation}
where $\text{Res}$ denotes the Poincaré residue map.

This pairing is also perfect, by the logarithmic Poincar\'e lemma
(Deligne~\cite{Deligne71}, Hodge~II, Proposition~3.1.8), which
identifies $H^*(\Omega^*_{\log}(\overline{C}_n(X) \setminus D))
\cong H^*(C_n(X))$; the perfectness then follows from
Poincar\'e--Lefschetz duality for the smooth quasi-projective
variety~$C_n(X)$.
\end{proof}

\begin{corollary}[Duality for bar complexes; \ClaimStatusProvedHere]
\label{cor:duality-bar-complexes-complete}
The Verdier duality on $\overline{C}_n(X)$ induces a perfect pairing:
\begin{equation}
\langle -, - \rangle: \bar{B}^n(\mathcal{A}) \otimes \bar{B}^n(\mathcal{A}^!) \to
\mathbb{C}
\end{equation}
where $\mathcal{A}^!$ is the Koszul dual of $\mathcal{A}$.
\end{corollary}

\begin{proof}[Proof of Corollary~\ref{cor:duality-bar-complexes-complete}]
Recall that:
\begin{align}
\bar{B}^n(\mathcal{A}) &= \Gamma(\overline{C}_n(X), \mathcal{A}^{\boxtimes n} \otimes
\Omega^*_{\log})\\
\bar{B}^n(\mathcal{A}^!) &= \Gamma(\overline{C}_n(X), (\mathcal{A}^!)^{\boxtimes n}
\otimes \Omega^*_{\log})
\end{align}

By Koszul duality (Definition~\ref{def:koszul-dual-chiral}), there is a natural pairing:
\begin{equation}
\mathcal{A} \otimes \mathcal{A}^! \to \mathcal{O}_X
\end{equation}
which extends to:
\begin{equation}
\mathcal{A}^{\boxtimes n} \otimes (\mathcal{A}^!)^{\boxtimes n} \to \mathcal{O}_{X^n}
\end{equation}

Combining with the Verdier pairing on $\Omega^*_{\log}$ from Theorem~\ref{thm:verdier-duality-config-complete}, we obtain:
\begin{equation}
\langle s, t \rangle = \int_{\overline{C}_n(X)} (s \otimes t) \wedge (-)
\end{equation}
for $s \in \bar{B}^n(\mathcal{A})$ and $t \in \bar{B}^n(\mathcal{A}^!)$.

This pairing is perfect because both the Koszul pairing and the Verdier pairing are
perfect.
\end{proof}

\emph{Duality interchange for spectral sequences.}

\begin{lemma}[Spectral sequence duality; \ClaimStatusProvedHere]
\label{lem:ss-duality-complete}
The Verdier duality of Theorem~\ref{thm:verdier-duality-config-complete} induces an isomorphism
of spectral sequences:
\begin{equation}
(E_r^{p,q,g})_{\mathcal{A}} \cong ((E_r^{p,d-q,g})_{\mathcal{A}^!})^\vee
\end{equation}
for all $r \geq 1$, where $d = \dim_{\mathbb{R}} \overline{C}_n(X) = 2n$.
\end{lemma}

\begin{proof}[Proof of Lemma~\ref{lem:ss-duality-complete}]
\emph{$E_1$ page}: By definition,
\begin{align}
(E_1^{p,q,g})_{\mathcal{A}} &= H^q(\barB^p_{g,\mathrm{flat}}(\mathcal{A}))\\
(E_1^{p,d-q,g})_{\mathcal{A}^!} &= H^{d-q}(\barB^p_{g,\mathrm{flat}}(\mathcal{A}^!))
\end{align}

By Corollary~\ref{cor:duality-bar-complexes-complete}, the pairing:
\begin{equation}
\langle -, - \rangle: H^q(\barB^p_{g,\mathrm{flat}}(\mathcal{A})) \otimes
H^{d-q}(\barB^p_{g,\mathrm{flat}}(\mathcal{A}^!)) \to \mathbb{C}
\end{equation}
is perfect. Thus $(E_1^{p,q,g})_{\mathcal{A}} \cong ((E_1^{p,d-q,g})_{\mathcal{A}^!})^\vee$.

\emph{Differential $d_1$}: The differential $d_1: E_1^{p,q,g} \to E_1^{p+1,q,g}$ is
induced by the moduli space differential. Under Verdier duality:
\begin{equation}
\mathbb{D} \circ d_1 = (-1)^{p+q} d_1^\vee \circ \mathbb{D}
\end{equation}
where $d_1^\vee$ is the dual differential.

This sign is precisely the Koszul sign convention (see Appendix~\ref{app:signs}, \S\ref{app:sign-conventions}).
Thus the differential on $(E_1)_{\mathcal{A}}$ is dual to the differential on
$(E_1)_{\mathcal{A}^!}$, up to the appropriate sign.

\emph{Higher pages}: By induction, if $(E_r)_{\mathcal{A}} \cong ((E_r)_{\mathcal{A}^!})^\vee$,
then taking cohomology with respect to $d_r$ preserves this duality:
\begin{equation}
(E_{r+1})_{\mathcal{A}} = H(E_r, d_r)_{\mathcal{A}} \cong (H(E_r, d_r)_{\mathcal{A}^!})^\vee
= ((E_{r+1})_{\mathcal{A}^!})^\vee
\end{equation}

\emph{$E_\infty$ page}: Taking the limit $r \to \infty$:
\begin{equation}
(E_\infty^{p,q,g})_{\mathcal{A}} \cong ((E_\infty^{p,d-q,g})_{\mathcal{A}^!})^\vee
\end{equation}

But $E_\infty^{*,*,g}$ is the associated graded for the fixed-genus
Leray filtration, so:
\begin{equation}
\operatorname{gr}^{p}_{\mathrm{Leray}}
H^{p+q}(\barB^{(g)}_{\mathrm{flat}}(\mathcal{A}))
\cong
\left(
\operatorname{gr}^{p}_{\mathrm{Leray}}
H^{p+d-q}(\barB^{(g)}_{\mathrm{flat}}(\mathcal{A}^!))
\right)^\vee .
\end{equation}
\end{proof}

\begin{corollary}[Leray corrections and Verdier eigenspaces; \ClaimStatusConditional]
\label{cor:quantum-dual-complete}
For each genus $g \geq 1$ and every Koszul dual chiral pair
$(\mathcal{A}, \mathcal{A}^!)$, Verdier duality gives a perfect
pairing of Leray correction objects
\begin{equation}
G_g(\mathcal{A}) \otimes G_g(\mathcal{A}^!) \longrightarrow \mathbb{C}.
\end{equation}
If the flat fiber--center comparison identifies the relevant
$j_*$ and $j_!$ sectors with the Verdier eigenspaces and if the
non-degenerate anti-invariant pairing hypothesis of
Theorem~\textup{\ref{thm:quantum-complementarity-main}} holds, then
\[
Q_g(\mathcal{A}) \cong Q_g(\mathcal{A}^!)^\vee .
\]
At genus $0$, the cohomological-shadow eigenspaces are instead
$Q_0(\mathcal{A}) = H^0(\overline{\mathcal{M}}_0, Z(\mathcal{A}))$
and $Q_0(\mathcal{A}^!) = 0$
\textup{(}Theorem~\textup{\ref{thm:quantum-complementarity-main}}\textup{)}.
\end{corollary}

\begin{proof}[Proof of Corollary~\ref{cor:quantum-dual-complete}]
The pairing on $G_g$ is Lemma~\ref{lem:ss-duality-complete} after
passing to $E_\infty$ and summing over total degree. The final
statement is exactly the Lagrangian clause of
Theorem~\ref{thm:quantum-complementarity-main}, applied after the
comparison map of Lemma~\ref{lem:quantum-from-ss}.
\end{proof}

The Verdier package gives duality of the fixed-genus Leray spectral
sequences and, after the finite-type pairing hypotheses are imposed,
the perfect cross-pairing between the two Verdier eigenspaces:
\begin{equation}
Q_g(\mathcal{A}) \cong Q_g(\mathcal{A}^!)^\vee
\qquad (g\geq1).
\end{equation}
\end{proof}

\subsection{Kodaira--Spencer action and Verdier eigenspaces}

\begin{proof}[Kodaira--Spencer and eigenspace package]

\emph{Center action on moduli space cohomology.}

\begin{theorem}[Kodaira--Spencer map for chiral algebras; \ClaimStatusConditional]
\label{thm:kodaira-spencer-chiral-complete}
\index{Kodaira--Spencer map|textbf}
\textup{[Regime: curved-central on the Koszul locus; positive genus
\textup{(}Convention~\textup{\ref{conv:regime-tags})}.]}

Let $(\mathcal{A}, \mathcal{A}^!)$ be a chiral Koszul pair with
$g \geq 1$ carrying the center-comparison datum of
Lemma~\textup{\ref{lem:center-isomorphism}}, and let $\pi: \mathcal{C}_g
\to \overline{\mathcal{M}}_g$ be the universal curve. There is a natural action:
\begin{equation}
\rho: Z(\mathcal{A}) \to \mathrm{End}\bigl(H^*(\overline{\mathcal{M}}_g, Z(\mathcal{A}))\bigr)
\end{equation}
constructed from the Gauss--Manin connection on the chiral homology sheaf.

This action interacts with Verdier duality via anti-commutativity: let
$\mathbb{D}: H^*(\barB^{(g)}_{\mathrm{flat}}(\mathcal{A})) \xrightarrow{\sim}
H^*(\barB^{(g)}_{\mathrm{flat}}(\mathcal{A}^!))^\vee$ be the Verdier
isomorphism of Corollary~\ref{cor:quantum-dual-complete}.
Then for every $z \in Z(\mathcal{A})$ and $v \in T_{[\Sigma_g]}\overline{\mathcal{M}}_g$:
\begin{equation}\label{eq:verdier-ks-anticommute}
\mathbb{D} \circ \nabla_{\kappa(v)}^z =
-\nabla_{\kappa(v)}^{\iota_Z(z)} \circ \mathbb{D}
\end{equation}
where $\kappa(v)$ is the Kodaira--Spencer class associated to $v$ and $\nabla^z$
denotes the action of $z$ via the Gauss--Manin connection.
\end{theorem}

\begin{proof}[Proof of Theorem~\ref{thm:kodaira-spencer-chiral-complete}]
We construct the action in three stages: the Gauss--Manin connection on the chiral
homology sheaf, the center action on this sheaf, and the anti-commutativity with
Verdier duality.

\emph{Stage 1: Gauss--Manin connection on chiral homology.}
Consider the universal curve $\pi: \mathcal{C}_g \to \overline{\mathcal{M}}_g$.
The classical Kodaira--Spencer map is:
\begin{equation}
\kappa: T_{\overline{\mathcal{M}}_g} \to R^1\pi_* T_{\mathcal{C}_g/\overline{\mathcal{M}}_g}
\end{equation}
For a chiral algebra $\mathcal{A}$ on the fibers, the relative chiral homology
$R^q\pi_*^{\mathrm{ch}} \mathcal{A}$ forms a sheaf on $\overline{\mathcal{M}}_g$.
By Beilinson--Drinfeld \cite[Theorem 4.6.1]{BD04}, this sheaf carries a flat
connection (the \emph{Gauss--Manin connection}):
\begin{equation}
\nabla^{\mathrm{GM}}: R^q\pi_*^{\mathrm{ch}} \mathcal{A} \to
R^q\pi_*^{\mathrm{ch}} \mathcal{A} \otimes \Omega^1_{\overline{\mathcal{M}}_g}(\log \partial)
\end{equation}
where $\partial = \overline{\mathcal{M}}_g \setminus \mathcal{M}_g$ is the boundary
divisor.

At the fiber over $[\Sigma_g] \in \mathcal{M}_g$, a tangent vector $v \in
T_{[\Sigma_g]}\mathcal{M}_g$ determines a first-order deformation $\Sigma_{g,\varepsilon}$
of $\Sigma_g$. The KS class $\kappa(v) \in H^1(\Sigma_g, T_{\Sigma_g})$ acts on the
bar complex via the connecting homomorphism of the short exact sequence:
\begin{equation}
0 \to \mathcal{A}|_{\Sigma_g} \to \mathcal{A}|_{\Sigma_{g,\varepsilon}} \to
\mathcal{A}|_{\Sigma_g} \otimes \Omega^1_{\Sigma_g} \to 0
\end{equation}
(twisting by the normal bundle of the central fiber). The connecting homomorphism gives:
\begin{equation}
\delta_v: H^q(\Sigma_g, \mathcal{A}) \to H^{q+1}(\Sigma_g, \mathcal{A})
\end{equation}
and $\nabla_v^{\mathrm{GM}} = \partial_v + \delta_v$ where $\partial_v$ is the
ordinary derivative along $v$.

\emph{Stage 2: Center action on moduli cohomology.}
Since $z \in Z(\mathcal{A})$ commutes with all chiral operations, $z$ defines a
global section of the endomorphism sheaf
$\mathcal{E}nd(R^q\pi_*^{\mathrm{ch}}\mathcal{A})$ that is flat with respect to
$\nabla^{\mathrm{GM}}$:
\begin{equation}
[\nabla^{\mathrm{GM}}, \rho(z)] = 0
\end{equation}
(centrality of $z$ implies it commutes with the connecting homomorphism $\delta_v$,
and it trivially commutes with $\partial_v$).

Taking cohomology over $\overline{\mathcal{M}}_g$, we obtain the action:
\begin{equation}
\rho(z): H^p(\overline{\mathcal{M}}_g, R^q\pi_*^{\mathrm{ch}}\mathcal{A}) \to
H^p(\overline{\mathcal{M}}_g, R^q\pi_*^{\mathrm{ch}}\mathcal{A})
\end{equation}
By Lemma~\ref{lem:fiber-cohomology-center}, the fiber cohomology sheaf restricted to
$\mathcal{M}_g^{\mathrm{smooth}}$ is $Z(\mathcal{A}) \otimes \underline{\mathbb{C}}$,
so this specialises to:
\begin{equation}
\rho(z): H^*(\overline{\mathcal{M}}_g, Z(\mathcal{A})) \to
H^*(\overline{\mathcal{M}}_g, Z(\mathcal{A}))
\end{equation}

\emph{Stage 3: Anti-commutativity with Verdier duality.}
We prove the anti-commutativity~\eqref{eq:verdier-ks-anticommute} under the
center-comparison and perfectness hypotheses, using three independent ingredients.

\emph{Ingredient 1: Lie derivative anti-commutes with Verdier duality.}
The Verdier duality $\mathbb{D}$ on the FM compactification $\overline{C}_n(X)$
exchanges $j_*$ and $j_!$ extensions (Lemma~\ref{lem:verdier-extension-exchange}):
$\mathbb{D}(j_* \mathcal{F}) \cong j_!(\mathbb{D}\mathcal{F})$.

The KS class $\kappa(v)$ acts on $H^*(\overline{C}_n(\Sigma_g), j_*\mathcal{A}^{\boxtimes n})$
via the Lie derivative $\mathcal{L}_{\tilde{v}}$ along the horizontal lift $\tilde{v}$
of $v$ to $\overline{C}_n(\Sigma_g)$. Since Verdier duality is a contravariant functor
on the derived category of $\mathcal{D}$-modules (it sends left $\mathcal{D}$-modules
to right $\mathcal{D}$-modules, reversing the sign of infinitesimal actions; this
sign reversal is a consequence of the contravariance of the Verdier duality functor
$\mathbb{D}$: if $\mathcal{F}$ carries an infinitesimal $G$-action $\rho$, then
$\mathbb{D}(\mathcal{F})$ carries the contragredient action
$\rho^\vee = -\rho^t$, which reverses signs):
\begin{equation}\label{eq:verdier-lie-anticommute}
\mathbb{D} \circ \mathcal{L}_{\tilde{v}} = -\mathcal{L}_{\tilde{v}} \circ \mathbb{D}
\end{equation}

Concretely, $\mathbb{D}$ reverses the orientation of the fibers of the FM
compactification, and the Lie derivative reverses sign under orientation reversal.

\emph{Ingredient 2: Verdier duality intertwines center actions through the supplied comparison.}
For $z \in Z(\mathcal{A})$, the action $\rho(z)$ on the bar complex
$\bar{B}^{(g)}(\mathcal{A})$ is a $\mathcal{D}$-module endomorphism (since $z$ is
central). The datum $\iota_Z$ identifies the center local systems in a way
compatible with the Verdier--Koszul pairing. Hence the transpose of the
$Z(\mathcal{A})$-action is the $Z(\mathcal{A}^!)$-action of $\iota_Z(z)$:

\begin{equation}\label{eq:verdier-center-intertwine}
\mathbb{D} \circ \rho(z) = \rho(\iota_Z(z)) \circ \mathbb{D}
\quad \text{for all } z \in Z(\mathcal{A})
\end{equation}

For $\omega \in \bar{B}^{(g)}(\mathcal{A})$ and $\omega' \in \bar{B}^{(g)}(\mathcal{A}^!)$:
$\langle \rho(z)\omega, \omega' \rangle_{\mathrm{Verdier}} =
\langle \omega, \rho(\iota_Z(z))\omega' \rangle_{\mathrm{Verdier}}$
by compatibility of $\iota_Z$ with the Koszul pairing.

\emph{Ingredient 3: The involution absorbs the transpose.}
Let $\sigma = \mathbb{D} \circ ((-)^!)^{-1}$ be the Verdier--Koszul involution
(Lemma~\ref{lem:verdier-involution-moduli}). The inverse comparison
$\iota_Z^{-1}$ sends $\iota_Z(z) \in Z(\mathcal{A}^!)$ back to
$z \in Z(\mathcal{A})$. Therefore:
\begin{equation}
\sigma \circ \rho(z) = \rho(z) \circ \sigma
\quad \text{for all } z \in Z(\mathcal{A})
\end{equation}

\emph{Assembly.} The combined action $\nabla_{\kappa(v)}^z := \rho(z) \circ
\mathcal{L}_{\tilde{v}}$ satisfies:
\begin{align}
\mathbb{D} \circ \nabla_{\kappa(v)}^z
&= \mathbb{D} \circ \rho(z) \circ \mathcal{L}_{\tilde{v}}
&& \text{(definition)} \\
&= \rho(\iota_Z(z)) \circ \mathbb{D} \circ \mathcal{L}_{\tilde{v}}
&& \text{(by~\eqref{eq:verdier-center-intertwine})} \\
&= -\rho(\iota_Z(z)) \circ \mathcal{L}_{\tilde{v}} \circ \mathbb{D}
&& \text{(by~\eqref{eq:verdier-lie-anticommute})} \\
&= -\nabla_{\kappa(v)}^{\iota_Z(z)} \circ \mathbb{D}
&& \text{(definition)}
\end{align}
Composing with the center comparison (which sends $\iota_Z(z)$ back to $z$), the
involution $\sigma$ satisfies:
\begin{equation}
\sigma \circ \nabla_{\kappa(v)}^z = -\nabla_{\kappa(v)}^z \circ \sigma
\end{equation}
which is the anti-commutativity~\eqref{eq:verdier-ks-anticommute} with $\mathbb{D}$
replaced by the full involution $\sigma$. This is a theorem for Koszul chiral
pairs carrying the stated comparison datum, not for a bare forgetful shadow.

\emph{Verification for Heisenberg.} For $\mathcal{H}_\kappa$ with center
$Z(\mathcal{H}_\kappa) = \mathbb{C} \cdot \mathbf{1} \oplus \mathbb{C} \cdot \kappa$,
the comparison $\iota_Z$ is the identity on $\mathbb{C} \cdot \kappa$ (the level is
self-dual under the Koszul pairing), so the general argument specialises to the
scalar commutation $[\mathbb{D}, \rho(\kappa)] = 0$. The action of $\kappa$ on
$H^*(\overline{\mathcal{M}}_{1,1})$ is:
\begin{equation}
\rho(\kappa)(\alpha) = \lambda_1 \cup \alpha
\end{equation}
where $\lambda_1 = c_1(\mathbb{E})$ is the first Chern class of the Hodge bundle.
This is computed via the Gauss--Manin connection: the central charge $\kappa$ pairs with
the KS class $\partial/\partial\tau$ to give $\lambda_1$ (Mumford isomorphism).
\end{proof}

\emph{Verdier involution and eigenspace decomposition.}

The anti-commutativity~\eqref{eq:verdier-ks-anticommute} is the key to the
eigenspace decomposition. We construct a canonical involution on $H^*(\overline{
\mathcal{M}}_g, Z(\mathcal{A}))$ whose $\pm 1$ eigenspaces give $Q_g(\mathcal{A})$
and $Q_g(\mathcal{A}^!)$.

\begin{lemma}[Verdier involution on moduli cohomology; \ClaimStatusConditional]
\label{lem:verdier-involution-moduli}
For a chiral Koszul pair $(\mathcal{A}, \mathcal{A}^!)$, Verdier duality on
$\overline{C}_n(X)$ together with the Koszul pairing $\mathcal{A} \otimes
\mathcal{A}^! \to \mathcal{O}_X$ induces an involution:
\begin{equation}
\sigma: H^*(\overline{\mathcal{M}}_g, Z(\mathcal{A})) \xrightarrow{\sim}
H^*(\overline{\mathcal{M}}_g, Z(\mathcal{A}))
\end{equation}
satisfying $\sigma^2 = \mathrm{id}$.
\end{lemma}

\begin{proof}[Proof of Lemma~\ref{lem:verdier-involution-moduli}]
\emph{Construction.}
By Corollary~\ref{cor:duality-bar-complexes-complete}, the
Verdier-Koszul pairing identifies:
\begin{equation}
\mathbb{D}: H^*(\barB^{(g)}_{\mathrm{flat}}(\mathcal{A})) \xrightarrow{\sim}
H^*(\barB^{(g)}_{\mathrm{flat}}(\mathcal{A}^!))^\vee
\end{equation}
Since $(\mathcal{A}^!)^! \simeq \mathcal{A}$ (Koszul involutivity,
Theorem~\ref{thm:chiral-koszul-duality}), we also have:
\begin{equation}
\mathbb{D}': H^*(\barB^{(g)}_{\mathrm{flat}}(\mathcal{A}^!)) \xrightarrow{\sim}
H^*(\barB^{(g)}_{\mathrm{flat}}(\mathcal{A}))^\vee
\end{equation}
The composition $(\mathbb{D}')^\vee \circ \mathbb{D}: H^*(\barB^{(g)}_{\mathrm{flat}}(\mathcal{A}))
\to H^*(\barB^{(g)}_{\mathrm{flat}}(\mathcal{A}))$ is the identity (by involutivity of Verdier
duality on compact smooth varieties, see Theorem~\ref{thm:verdier-duality-config-complete}).

Now, the $E_2$ page of the Leray spectral sequence
(Theorem~\ref{thm:ss-quantum}) identifies both $Q_g(\mathcal{A})$ and
$Q_g(\mathcal{A}^!)$ as subquotients of $H^*(\overline{\mathcal{M}}_g,
Z(\mathcal{A}))$, provided the following center-comparison datum is supplied.

\begin{lemma}[Center comparison datum; \ClaimStatusConditional]
\label{lem:center-isomorphism}
Let $(\mathcal{A},\mathcal{A}^!)$ be a Koszul chiral pair. Suppose the strict
flat Koszul package contains an isomorphism of center local systems
\[
  \iota_Z \colon Z(\mathcal{A}) \xrightarrow{\sim} Z(\mathcal{A}^!)
\]
compatible with the Verdier--Koszul pairing. Then the two Leray $E_2$ pages for
$\mathcal{A}$ and $\mathcal{A}^!$ may be read with coefficients in the same
local system.
\end{lemma}

\begin{proof}[Proof of Sublemma]
This is a datum, not a consequence of module Koszul duality alone. The
$E_1$ module Koszul equivalence
(Theorem~\ref{thm:e1-module-koszul-duality}) gives, under its stated
hypotheses,
\begin{equation}
\Phi: D^b(\mathrm{Mod}^{\Eone}_{\mathcal{A}}) \xrightarrow{\;\sim\;}
D^b(\mathrm{Mod}^{\Eone}_{\mathcal{A}^!})
\end{equation}
and therefore identifies endomorphism algebras of corresponding modules,
\begin{equation}
\mathrm{End}_{D(\mathrm{Mod}_{\mathcal{A}})}(M)
\;\xrightarrow[\Phi]{\;\sim\;}\;
\mathrm{End}_{D(\mathrm{Mod}_{\mathcal{A}^!})}(\Phi M).
\end{equation}
For $M=\mathcal{A}$ this is the endomorphism algebra of the regular module.
It is not, without further centrality and flatness data, the center local
system used in Theorem~C. The forgetful passage from commutative to associative
structure therefore gives no center isomorphism by itself. Once $\iota_Z$ is
part of the strict flat package, it transports the coefficients of the
$\mathcal{A}^!$ Leray sequence to the coefficient system for $\mathcal{A}$,
which is the assertion.
\end{proof}

The Verdier isomorphism
$\mathbb{D}$ then exchanges these subquotients, inducing the involution
$\sigma$ on the total space. Explicitly, $\sigma$ acts on the $E_2$ page by:
\begin{equation}
\sigma: E_2^{p,q}(\mathcal{A}) \to E_2^{p,d-q}(\mathcal{A}^!) \xrightarrow{\iota_Z}
E_2^{p,d-q}(\mathcal{A})
\end{equation}
(the last isomorphism via the supplied center comparison). This descends to $E_\infty$:
Verdier duality is an exact functor on the derived category of constructible sheaves,
and the Leray spectral sequence is functorial with respect to exact functors.
Therefore $\sigma$, being induced by the composition of Verdier duality with the
center comparison, commutes with the spectral sequence differentials $d_r$
at each page (with sign absorbed by the Koszul sign convention, cf.\
Lemma~\ref{lem:ss-duality-complete}).

\emph{Eigenvalue identification.}
The key sign is: $j_*$ and $j_!$ differ by a shift of $2\dim(\partial)$ in the
dualizing complex, which introduces a factor of $(-1)^{\dim \partial}$ in the
eigenvalue identification.

\emph{Involutivity.} $\sigma^2 = \mathrm{id}$ follows from
$\mathbb{D}^2 = \mathrm{id}$ on the compact smooth variety $\overline{C}_n(X)$
(Poincar\'e duality is an involution).
\end{proof}

\begin{lemma}[Eigenspace decomposition; \ClaimStatusConditional]
\label{lem:eigenspace-decomposition-complete}
The involution $\sigma$ decomposes moduli cohomology into $\pm 1$
eigenspaces:
\begin{equation}
H^*(\overline{\mathcal{M}}_g, Z(\mathcal{A})) =
H^*(\overline{\mathcal{M}}_g, Z(\mathcal{A}))^+ \oplus
H^*(\overline{\mathcal{M}}_g, Z(\mathcal{A}))^-
\end{equation}
where $H^*(\overline{\mathcal{M}}_g, Z(\mathcal{A}))^\pm = \ker(\sigma \mp
\mathrm{id})$. These eigenspaces are identified with the quantum corrections:
\begin{align}
Q_g(\mathcal{A}) &= H^*(\overline{\mathcal{M}}_g, Z(\mathcal{A}))^+\\
Q_g(\mathcal{A}^!) &= H^*(\overline{\mathcal{M}}_g, Z(\mathcal{A}))^-
\end{align}
\end{lemma}

\begin{proof}[Proof of Lemma~\ref{lem:eigenspace-decomposition-complete}]
Since $\sigma^2 = \mathrm{id}$ on the finite-dimensional vector space $V :=
H^*(\overline{\mathcal{M}}_g, Z(\mathcal{A}))$, the eigenvalues of $\sigma$ are
$\pm 1$, and $V = V^+ \oplus V^-$ (standard linear algebra over $\mathbb{C}$;
the projectors are $\pi^\pm = \frac{1}{2}(\mathrm{id} \pm \sigma)$).

\emph{Identification of $V^+$ with $Q_g(\mathcal{A})$.}
By construction, $\sigma$ acts on $E_\infty^{*,*,g}$ by exchanging the contributions
from $\mathcal{A}$ and $\mathcal{A}^!$. A class $\alpha \in Q_g(\mathcal{A}) =
\mathrm{gr}^g H^*(\bar{B}(\mathcal{A}))$ maps under $\mathbb{D}$ to a class in
$Q_g(\mathcal{A}^!)^\vee$. Under the identification with moduli cohomology,
$\alpha$ contributes to $V^+$ because:
\begin{enumerate}
\item The bar complex of $\mathcal{A}$ uses $j_*$-extension.
\item The cobar complex of $\mathcal{A}$ (equivalently, the bar of $\mathcal{A}^!$)
uses $j_!$-extension.
\item Verdier duality exchanges $j_*$ and $j_!$
(Lemma~\ref{lem:verdier-extension-exchange}).
\item The involution $\sigma$ acts on the $E_2$ page by
$\sigma: E_2^{p,q}(\mathcal{A}) \to E_2^{p,d-q}(\mathcal{A})$.
By the fiber cohomology vanishing
(Lemma~\ref{lem:fiber-cohomology-center}), the classes
$Q_g(\mathcal{A})$ live on the $q = 0$ row of the $E_2$ page.
Poincar\'e duality sends $H^q$ to $H^{d-q}$ with sign $(-1)^q$;
at $q = 0$ the sign is $(-1)^0 = +1$. Therefore elements of
$Q_g(\mathcal{A})$ have eigenvalue~$+1$.
\end{enumerate}

\emph{Identification of $V^-$ with $Q_g(\mathcal{A}^!)$.}
The bar complex of $\mathcal{A}^!$ on configuration spaces uses
$j_!$-extension (compact support), dual to the $j_*$-extension used
by $\bar{B}^{(g)}(\mathcal{A})$
(Lemma~\ref{lem:verdier-extension-exchange}).
By the same fiber cohomology vanishing applied to the Koszul dual,
elements of $Q_g(\mathcal{A}^!)$ also live on the $q = 0$ row of
\emph{their own} $E_2$ page.

However, the Verdier involution $\sigma$ on
$\mathcal{H}_g(\mathcal{A})$ identifies $Q_g(\mathcal{A}^!)$ not
through its own spectral sequence but through the $j_! \leftrightarrow j_*$
interchange. Concretely, $\sigma$ sends
$j_!$-extended classes to $j_*$-extended classes.
The eigenvalue is determined by the anti-commutativity
of~$\sigma$ with the Kodaira--Spencer action
\eqref{eq:verdier-ks-anticommute}: for $\beta$ arising from
$\bar{B}^{(g)}(\mathcal{A}^!)$,
\[
\sigma(\beta)
= \sigma\bigl(\text{$j_!$-extension}\bigr)
= \text{$j_*$-extension}
= -\beta
\]
where the sign $-1$ arises because the $j_! \to j_*$ natural
transformation, composed with the Koszul identification
$(\mathcal{A}^!)^! \simeq \mathcal{A}$, picks up the sign from
the anti-commutativity~\eqref{eq:verdier-ks-anticommute}.
Therefore elements from $\bar{B}^{(g)}(\mathcal{A}^!)$ have
eigenvalue $-1$ and contribute to~$V^-$.

The geometric content of the eigenspace decomposition
is the $j_*/j_!$ dichotomy.
Bar cochains of~$\mathcal{A}$ on the FM compactification
$\overline{C}_n(X)$ are sections of
$j_*j^*\mathcal{A}^{\boxtimes n}$: they \emph{extend across the
boundary} with logarithmic poles (the $j_*$-extension, which
permits growth along boundary divisors).
Cobar cochains of~$\mathcal{A}$ (equivalently, bar cochains
of~$\mathcal{A}^!$) are sections of
$j_!j^*(\mathcal{A}^!)^{\boxtimes n}$: they have \emph{compact
support} away from the boundary (the $j_!$-extension, which
forces vanishing along boundary divisors).
The Verdier involution exchanges $j_*$ and $j_!$
(Lemma~\ref{lem:verdier-extension-exchange}), hence exchanges
$\mathcal{A}$-contributions and $\mathcal{A}^!$-contributions.
The identification of the $+1$-eigenspace with
$Q_g(\mathcal{A})$ and the $-1$-eigenspace with
$Q_g(\mathcal{A}^!)$ then follows from the Verdier intertwining
$\mathbb{D}_{\mathrm{Ran}}\, \bar{B}(\mathcal{A}) \simeq
\cA^!_\infty$ (Theorem~\ref{thm:verdier-bar-cobar}):
at the level of centers,
$\sigma(\mathcal{Z}(\mathcal{A})) = \mathcal{Z}(\mathcal{A}^!)$,
so $\sigma$ acts as $+1$ on $j_*$-extended classes from~$\mathcal{A}$
and as $-1$ on $j_!$-extended classes from~$\mathcal{A}^!$.

More precisely, the anti-commutativity
\eqref{eq:verdier-ks-anticommute} ensures that the Kodaira--Spencer
deformation maps $V^+$ to $V^-$ and vice versa. If $\alpha \in V^+$,
then
\begin{equation}
\sigma(\nabla_{\kappa(v)}^z \alpha) = -\nabla_{\kappa(v)}^z (\sigma \alpha) =
-\nabla_{\kappa(v)}^z \alpha
\end{equation}
so $\nabla_{\kappa(v)}^z \alpha \in V^-$. The two eigenspaces therefore
carry genuinely different deformation-theoretic content.
\end{proof}

\begin{lemma}[Obstructions vs.\ deformations; \ClaimStatusConditional]
\label{lem:obs-def-split-complete}
In the decomposition of Lemma~\ref{lem:eigenspace-decomposition-complete}:
\begin{enumerate}
\item Elements of $Q_g(\mathcal{A}) = V^+$ correspond to \emph{obstructions}
to extending $\mathcal{A}$ from genus $g-1$ to genus $g$
\item Elements of $Q_g(\mathcal{A}^!) = V^-$ correspond to \emph{deformations}
of $\mathcal{A}^!$ at genus $g$
\item What $\mathcal{A}$ sees as obstruction, $\mathcal{A}^!$ sees as deformation
\end{enumerate}
\end{lemma}

\begin{proof}[Proof of Lemma~\ref{lem:obs-def-split-complete}]
The bar complex $\bar{B}^{(g)}(\mathcal{A})$ computes obstruction
classes: a class
$\alpha \in H^*(\barB^{(g)}_{\mathrm{flat}}(\mathcal{A}))$
represents a cohomological obstruction to extending the
genus-$(g{-}1)$ structure to genus~$g$. By
Lemma~\ref{lem:quantum-from-ss}, their Leray graded pieces form
$G_g(\cA)$; after the fiber--center comparison and Verdier projection,
the $j_*$-extended $\cA$ sector lies in~$V^+$ by the previous lemma.

Conversely, the bar complex $\bar{B}^{(g)}(\mathcal{A}^!)$ computes obstructions
for $\mathcal{A}^!$, which correspond to \emph{deformations} of $\mathcal{A}$
by Koszul duality: the deformation complex of $\mathcal{A}$ is quasi-isomorphic
to the obstruction complex of $\mathcal{A}^!$ (Theorem~\ref{thm:chiral-koszul-duality}).
These lie in $V^-$.
\end{proof}

\emph{Eigenspace intersection.}

\begin{lemma}[Trivial eigenspace intersection; \ClaimStatusConditional]
\label{lem:trivial-intersection-complete}
The cohomological-shadow eigenspaces intersect trivially:
\begin{equation}
Q_g(\mathcal{A}) \cap Q_g(\mathcal{A}^!) = 0
\end{equation}
\end{lemma}

\begin{proof}[Proof of Lemma~\ref{lem:trivial-intersection-complete}]
By Lemma~\ref{lem:eigenspace-decomposition-complete}, $Q_g(\mathcal{A}) = V^+$ and
$Q_g(\mathcal{A}^!) = V^-$ are the $+1$ and $-1$ eigenspaces of the involution
$\sigma$. Since $+1 \neq -1$ (we work over $\mathbb{C}$, characteristic $\neq 2$),
eigenspaces for distinct eigenvalues of any linear operator intersect trivially:
\[
V^+ \cap V^- = 0. \qedhere
\]
\end{proof}

\emph{Eigenspace exhaustion.}

\begin{lemma}[Verdier eigenspace exhaustion; \ClaimStatusConditional]
\label{lem:exhaustion-complete}
The Verdier eigenspaces exhaust the moduli space cohomology:
\begin{equation}
Q_g(\mathcal{A}) \oplus Q_g(\mathcal{A}^!)
= H^*(\overline{\mathcal{M}}_g,Z(\mathcal{A})).
\end{equation}
\end{lemma}

\begin{proof}[Proof of Lemma~\ref{lem:exhaustion-complete}]
Since $\sigma^2 = \mathrm{id}$ on the finite-dimensional vector space
$V = H^*(\overline{\mathcal{M}}_g, Z(\mathcal{A}))$, the minimal polynomial of
$\sigma$ divides $t^2 - 1 = (t-1)(t+1)$. Over $\mathbb{C}$, every vector decomposes
uniquely into $\pm 1$ eigenvectors via the projectors
$\pi^+ = \frac{1}{2}(\mathrm{id} + \sigma)$ and
$\pi^- = \frac{1}{2}(\mathrm{id} - \sigma)$. Therefore:
\begin{equation}
V = V^+ \oplus V^- = Q_g(\mathcal{A}) \oplus Q_g(\mathcal{A}^!)
\end{equation}
and in particular $\dim V^+ + \dim V^- = \dim V$.

\emph{Verification at small genera.}

\emph{Genus $0$}: $\overline{\mathcal{M}}_0 = \mathrm{pt}$, so
$\dim H^* = 1$. The involution $\sigma$ acts as $+1$ on $H^0$ (identity class),
giving $Q_0(\mathcal{A}) = \mathbb{C}$, $Q_0(\mathcal{A}^!) = 0$. Dimension check:
$1 + 0 = 1$. This is the genus-$0$ exception to the duality clause in
Theorem~\ref{thm:quantum-complementarity-main}.

\emph{Genus $1$}: $\dim H^*(\overline{\mathcal{M}}_{1,1}) = 2$ ($H^0 \oplus H^2$).
For the Heisenberg algebra: $Q_1(\mathcal{H}_\kappa) = \mathbb{C} \cdot \kappa$
(central extension, eigenvalue $+1$) and $Q_1(\mathcal{H}_\kappa^!) = \mathbb{C}
\cdot \lambda$ (curvature, eigenvalue $-1$). Dimension check: $1 + 1 = 2$.

\emph{Genus $2$}: $\dim H^*(\overline{\mathcal{M}}_2) = 8$
(Poincar\'e polynomial $1 + 3t^2 + 3t^4 + t^6$).
The involution $\sigma$ splits $H^*$ into $\pm 1$-eigenspaces whose
dimensions depend on the chiral algebra through the $j_*/j_!$ extension
data. For balanced Koszul self-dual algebras (e.g.,
$\mathrm{Vir}_{13}$ after the anti-equivariant representative
hypothesis of Corollary~\ref{cor:virasoro-quantum-dim}),
Theorem~\ref{thm:self-dual-halving} gives
$\dim V^+ = \dim V^- = 4$. For every Koszul chiral algebra
$\cA$ with trivial center satisfying the hypotheses above,
$\dim Q_2(\cA) + \dim Q_2(\cA^!) = 8$.
Detailed genus-$2$ complementarity dimensions are verified in
Part~\ref{part:characteristic-datum} (\S\ref{subsec:genus2-complementarity-verification}).
\end{proof}

The Kodaira--Spencer and Verdier package gives:
\begin{enumerate}
\item a Kodaira--Spencer map via Gauss--Manin connection, with anti-commutativity
$\mathbb{D} \circ \nabla_{\kappa(v)}^z =
-\nabla_{\kappa(v)}^{\iota_Z(z)} \circ \mathbb{D}$
\item a Verdier involution $\sigma$ with $\sigma^2 = \mathrm{id}$
\item Direct sum: $Q_g(\mathcal{A}) \cap Q_g(\mathcal{A}^!) = 0$ as eigenspaces of
$\sigma$
\item Exhaustion: $Q_g(\mathcal{A}) \oplus Q_g(\mathcal{A}^!) = H^*(\overline{
\mathcal{M}}_g, Z(\mathcal{A}))$ because $\sigma^2 = \mathrm{id}$
\end{enumerate}

Therefore:
\begin{equation}
Q_g(\mathcal{A}) \oplus Q_g(\mathcal{A}^!) \cong H^*(\overline{\mathcal{M}}_g,
Z(\mathcal{A}))
\end{equation}

This completes the proof of Theorem~\ref{thm:quantum-complementarity-main}. \qedhere

\end{proof}

\begin{remark}[Cochain-level proof path for Theorem~C]
\label{rem:h-level-proof-summary}
\index{Theorem C!cochain-level proof path}
The argument above establishes both the
S-level (cohomological) and H-level (homotopy) claims of
Theorem~\ref{thm:quantum-complementarity-main}. The Verdier
involution~$\sigma$ is a
\emph{cochain-level} involution on
$\mathbf{C}_g(\cA) = R\Gamma(\overline{\mathcal{M}}_g, \mathcal{Z}(\cA))$:
it is induced by the composition of Verdier duality~$\mathbb{D}$
(an exact functor on $D^b_{\mathrm{hol}}(\mathcal{D})$) with the
supplied center comparison
$\mathcal{Z}(\cA) \cong \mathcal{Z}(\cA^!)$
(Lemma~\ref{lem:center-isomorphism} and
Theorem~\ref{thm:kodaira-spencer-chiral-complete}).
Since~$\sigma$ is induced by exact functors, it commutes with
differentials at every page of the spectral sequence
(Lemma~\ref{lem:verdier-involution-moduli}).
Lemma~\ref{lem:involution-splitting}(a) gives
the cochain-level splitting: the projectors
$p^\pm = \tfrac{1}{2}(\mathrm{id} \pm \sigma)$ are cochain maps
whose images provide a quasi-isomorphic direct sum decomposition
$\mathbf{C}_g \simeq \operatorname{im}(p^+) \oplus \operatorname{im}(p^-)$.
Model independence
(Proposition~\ref{prop:model-independence}) ensures the
decomposition depends only on the Koszul pair $(\cA, \cA^!)$,
not on auxiliary choices.
\end{remark}

\begin{remark}[Geometric substrate (Volume~II)]
\label{rem:theorem-c-lagrangian}
\index{Lagrangian self-intersection!Theorem C}
Volume~II interprets the complementary decomposition as two
Lagrangians $\mathcal{L}$ and $\mathcal{L}^!$ inside a
$(-2)$-shifted symplectic stack $\mathcal{M}$, whose derived
intersection $\mathcal{L} \times_{\mathcal{M}} \mathcal{L}^!$
carries the $(-1)$-shifted symplectic structure that governs the
Verdier pairing of Proposition~\ref{prop:lagrangian-eigenspaces}.
\end{remark}

\begin{proposition}[Verdier pairing and Lagrangian eigenspaces;
\ClaimStatusConditional]
\label{prop:lagrangian-eigenspaces}
\index{Lagrangian!eigenspaces|textbf}
\index{Verdier duality!pairing on moduli}
Let $(\cA, \cA^!)$ be a chiral Koszul pair and let
$V = H^*(\overline{\mathcal{M}}_g, \mathcal{Z}(\cA))$.
Assume that the Koszul pairing identifies
$\mathbb{D}\mathcal Z(\cA)$ with $\mathcal Z(\cA^!)$, that the resulting
pairing on $V$ is non-degenerate, and that the Verdier involution
$\sigma$ is anti-symplectic for this pairing.
\begin{enumerate}[label=\textup{(\roman*)}]
\item Verdier duality on the center local system $\mathcal{Z}(\cA)$
 induces a non-degenerate bilinear pairing
 \[
 \langle -, - \rangle_{\mathbb{D}} \colon V \otimes V \to \mathbb{C}
 \]
 of cohomological degree $-(3g-3)$ \textup{(}the shift coming from
 $\dim_{\mathbb{C}} \overline{\mathcal{M}}_g = 3g-3$\textup{)}.
\item The Verdier involution $\sigma$ is an \emph{anti-involution}
 for this pairing:
 $\langle \sigma(v), \sigma(w) \rangle_{\mathbb{D}} = -\langle v, w \rangle_{\mathbb{D}}$
 for all $v, w \in V$.
\item Consequently, the eigenspaces
 $V^+ = Q_g(\cA)$ and $V^- = Q_g(\cA^!)$
 are \emph{Lagrangian} with respect to
 $\langle -, - \rangle_{\mathbb{D}}$: the pairing restricts to
 zero on each eigenspace, and $V = V^+ \oplus V^-$ is a Lagrangian
 polarization.
\end{enumerate}
\end{proposition}

\begin{proof}
\emph{Part (i).}
The center local system $\mathcal{Z}(\cA)$ on
$\overline{\mathcal{M}}_g$ is a constructible sheaf, and Verdier
duality gives a perfect pairing
$H^k(\overline{\mathcal{M}}_g, \mathcal{Z}) \otimes
H^{(3g-3)-k}(\overline{\mathcal{M}}_g, \mathbb{D}\mathcal{Z})
\to \mathbb{C}$.
The Koszul pairing identifies $\mathbb{D}\mathcal{Z}(\cA) \cong
\mathcal{Z}(\cA^!) \cong \mathcal{Z}(\cA)$
(the last isomorphism by the center comparison of
Lemma~\ref{lem:center-isomorphism}), so the
pairing lands on $V \otimes V$. Non-degeneracy follows from Verdier
duality for constructible sheaves on the smooth Deligne--Mumford
stack $\overline{\mathcal{M}}_g$.

\emph{Part (ii).}
This is exactly the anti-symplectic compatibility assumed in the
statement. In applications it is verified by constructing the
Verdier representative of $\sigma$ and checking that the
Kodaira--Spencer connection anticommutes with Verdier duality; the
present proposition records the formal consequence once that
chain-level compatibility has been supplied.

\emph{Part (iii).}
For $v, w \in V^+$ (eigenvalue $+1$):
$\langle v, w \rangle_{\mathbb{D}}
= \langle \sigma(v), \sigma(w) \rangle_{\mathbb{D}}
= -\langle v, w \rangle_{\mathbb{D}}$,
so $\langle v, w \rangle_{\mathbb{D}} = 0$. The same argument
applies to $V^-$. Since $V = V^+ \oplus V^-$ and the pairing is
non-degenerate on $V$, the subspaces $V^+$ and $V^-$ are maximal
isotropic, i.e., Lagrangian.
\end{proof}

\begin{remark}[Lagrangian interpretation of complementarity]\label{rem:lagrangian-complementarity}
\index{Lagrangian!complementarity}
The decomposition of
Theorem~\ref{thm:quantum-complementarity-main} is an eigenspace
decomposition without further hypotheses. Under the additional
non-degenerate anti-invariant Verdier pairing of
Proposition~\ref{prop:lagrangian-eigenspaces}, the eigenspaces
$Q_g(\cA)$ and $Q_g(\cA^!)$ become complementary Lagrangians.
The $(-1)$-shifted bar-side symplectic structure requires the
separate BV package of
Theorem~\ref{thm:shifted-symplectic-complementarity}.
\end{remark}

\subsection{Shifted symplectic complementarity}
\label{sec:shifted-symplectic-complementarity}
\index{shifted symplectic!complementarity|textbf}

Proposition~\ref{prop:lagrangian-eigenspaces} establishes the
Lagrangian polarization at the level of cohomology, and
Proposition~\ref{prop:ptvv-lagrangian} upgrades the Verdier side to a
proved shifted-symplectic statement on the ambient complex~$C_g$. The
direct bar-side realization of a $(-1)$-shifted symplectic chart on
$L_g = \barB^{(g)}_{\mathrm{flat}}(\cA)[1]$ requires the
additional BV package of
Chapter~\ref{ch:v1-bv-brst}; the theorem below records precisely that
conditional bar-side upgrade. The transport from the ambient complex to a
bar chart factors through Theorem~\ref{thm:fiber-center-identification}.

\begin{lemma}[Bar chart transport of the ambient Lagrangian polarization;
\ClaimStatusConditional]
\label{lem:bar-chart-lagrangian-lift}
\index{Lagrangian!bar-chart transport}
\index{fiber--center identification!transport to bar chart}
Let $(\cA, \cA^!)$ be a chiral Koszul pair, let $g \geq 1$, and assume
the hypotheses of Theorem~\ref{thm:fiber-center-identification}.
Write
\[
\mathbf{C}_g(\cA) := R\Gamma(\overline{\mathcal{M}}_g, \mathcal{Z}(\cA)).
\]
Then Theorem~\ref{thm:fiber-center-identification} induces a canonical
quasi-isomorphism
\[
\gamma_g \colon
R\Gamma(\overline{\mathcal{M}}_g, R\pi_{g*}\barB^{(g)}_{\mathrm{flat}}(\cA))
\xrightarrow{\;\sim\;}
\mathbf{C}_g(\cA)
\]
intertwining the Verdier involution and the Verdier pairing. If
\[
\iota_g \colon L_g := \barB^{(g)}_{\mathrm{flat}}(\cA)[1]
\xrightarrow{\;\sim\;}
R\Gamma(\overline{\mathcal{M}}_g, R\pi_{g*}\barB^{(g)}_{\mathrm{flat}}(\cA))
\]
is a quasi-isomorphism of paired complexes, set
$\chi_g := \gamma_g \circ \iota_g$. Then:
\begin{enumerate}[label=\textup{(\roman*)}]
\item $\chi_g \colon L_g \xrightarrow{\sim} \mathbf{C}_g(\cA)$ is a
 quasi-isomorphism intertwining the involutions and the pairings.
\item If $\mathbf{Q}_g(\cA)$ and $\mathbf{Q}_g(\cA^!)$ are
 complementary Lagrangians in $\mathbf{C}_g(\cA)$, then the
 homotopy eigenspaces
 \[
 L_g^+ := \operatorname{fib}(\sigma_{L_g} - \mathrm{id}),
 \qquad
 L_g^- := \operatorname{fib}(\sigma_{L_g} + \mathrm{id})
 \]
 are complementary Lagrangians in $L_g$.
\item On cohomology,
 $H^*(L_g^+) = Q_g(\cA)$ and
 $H^*(L_g^-) = Q_g(\cA^!)$.
\end{enumerate}
\end{lemma}

\begin{proof}
Theorem~\ref{thm:fiber-center-identification} gives
the concentration
$\mathcal{H}^q(R\pi_{g*}\barB^{(g)}_{\mathrm{flat}}(\cA))=0$
for $q \neq 0$ and the identification
$\mathcal{H}^0(R\pi_{g*}\barB^{(g)}_{\mathrm{flat}}(\cA))
\cong \mathcal{Z}_{\cA}$. Applying
derived global sections to the perfect flat representative yields the
quasi-isomorphism~$\gamma_g$. The
Verdier involution and pairing are functorial under exact functors, so
$\gamma_g$ intertwines both structures.

If $\iota_g$ is a quasi-isomorphism of paired complexes, then so is
$\chi_g=\gamma_g\circ\iota_g$, proving~(i). For~(ii), transport the
projectors
$p^\pm=\tfrac{1}{2}(\mathrm{id}\pm\sigma)$ from
$\mathbf{C}_g(\cA)$ to $L_g$ along~$\chi_g$. Because~$\chi_g$
intertwines pairings, isotropy of
$\mathbf{Q}_g(\cA)$ and $\mathbf{Q}_g(\cA^!)$ pulls back to isotropy of
$L_g^+$ and $L_g^-$. Since $\chi_g$ is a quasi-isomorphism and
$\mathbf{C}_g(\cA)=\mathbf{Q}_g(\cA)\oplus \mathbf{Q}_g(\cA^!)$ is a
Lagrangian polarization, Lemma~\ref{lem:involution-splitting}(c)
implies that $L_g^+$ and $L_g^-$ are complementary Lagrangians. Part
(\textup{iii}) is Lemma~\ref{lem:involution-splitting}(b) applied to
$L_g$ and then identified through~$\chi_g$.
\end{proof}

\begin{theorem}[Conditional bar-side BV upgrade of complementarity; \ClaimStatusConditional]
\label{thm:shifted-symplectic-complementarity}
\index{shifted symplectic!(-1)-shifted|textbf}
\index{Lagrangian!complementarity!shifted symplectic|textbf}
\textup{[Regime: curved-central on the Koszul locus; all genera
\textup{(}Convention~\textup{\ref{conv:regime-tags})}.]}

Assume the conditional BV package of
Theorems~\ref{thm:config-space-bv} and~\ref{thm:bv-functor}. Assume
also the cyclic nondegeneracy input: the BV bracket adjoint
\[
x \longmapsto \{x,-\}_{\mathrm{BV}} \colon
\barB^{\mathrm{ch}}(\cA)
\longrightarrow
\mathbb{D}_{\operatorname{Ran}}(\barB^{\mathrm{ch}}(\cA))[1]
\]
is a quasi-isomorphism on the Koszul locus after the
bracket-compatible Verdier comparison. Then the
genus-$1$ complementarity
$Q_1(\mathcal{H}_\kappa) \oplus Q_1(\mathcal{H}_\kappa^!)
\cong H^*(\overline{\mathcal{M}}_1, Z(\mathcal{H}_\kappa))$ that we
computed in~\S\ref{sec:frame-complementarity} admits the following
bar-side $(-1)$-shifted symplectic refinement.
\index{BV algebra!shifted symplectic structure}
Let $(\cA, \cA^!)$ be a chiral Koszul pair on a smooth projective
curve $X$.
\begin{enumerate}[label=\textup{(\roman*)}]
\item \emph{BV antibracket is $(-1)$-shifted Poisson.}
 The BV bracket $\{-,-\}_{\mathrm{BV}}$ on
 $\barB^{\mathrm{ch}}(\cA)$
 \textup{(}Theorem~\textup{\ref{thm:config-space-bv}}\textup{)}
 has degree~$+1$, hence defines a $(-1)$-shifted Poisson structure.
 Under the bracket-compatible Verdier comparison of
 Theorem~\textup{\ref{thm:bv-functor}}, refined on the Koszul locus by
 Theorem~\textup{\ref{thm:verdier-bar-cobar}} to
 $\mathbb{D}_{\operatorname{Ran}}(\barB^{\mathrm{ch}}(\cA))
 \simeq \cA^!_\infty$, the adjoint map
 \[
 x \longmapsto \{x,-\}_{\mathrm{BV}} \colon
 \barB^{\mathrm{ch}}(\cA)
 \longrightarrow
 \mathbb{D}_{\operatorname{Ran}}(\barB^{\mathrm{ch}}(\cA))[1]
 \simeq
 \cA^!_\infty[1]
 \]
 is precisely the cyclic duality quasi-isomorphism assumed above. On
 the Koszul locus, the underlying complex
 of $\cA^!_\infty$ is equivalent to $\barB^{\mathrm{ch}}(\cA^!)$.
 Therefore this Poisson structure is non-degenerate.

\item \emph{Formal moduli is $(-1)$-shifted symplectic.}
 The dg Lie algebra $L_g := \barB^{(g)}(\cA)[1]$, with Lie bracket
 induced by the BV antibracket, carries a non-degenerate invariant
 pairing of degree~$-1$. By the Kontsevich--Pridham correspondence
 \cite{Pridham17}, the formal moduli problem
 $\mathrm{Def}_g(\cA) := \mathrm{MC}(L_g)$
 is $(-1)$-shifted symplectic.

 The dg~Lie algebras $L_g$ are genus-truncations (via the complete
 filtration of
 Proposition~\ref{prop:modular-deformation-truncation}) of the
 modular $L_\infty$-deformation object $\Definfmod(\cA)$
 (Theorem~\ref{thm:modular-homotopy-convolution}).
\item \emph{Bar-side eigenspaces lift the ambient Lagrangian polarization.}
 For $g \geq 1$, assume in addition that the bar chart $L_g$ fits into
 the transport square of
 Lemma~\textup{\ref{lem:bar-chart-lagrangian-lift}}. Then the
 eigenspace decomposition $L_g = L_g^+ \oplus L_g^-$ provides
 complementary Lagrangian subspaces lifting the ambient Verdier
 polarization, and
 $H^*(L_g^+) = Q_g(\cA)$, $H^*(L_g^-) = Q_g(\cA^!)$.
\end{enumerate}
\end{theorem}

\begin{proof}
\emph{Part (i).}
The BV bracket $\{-,-\}_{\mathrm{BV}}$ on $\barB^{\mathrm{ch}}(\cA)$
has degree~$+1$ by Theorem~\ref{thm:config-space-bv}. A degree~$+1$
Lie bracket is by definition a $(-1)$-shifted Poisson structure: the
associated bivector field on the formal moduli problem has degree~$-1$.
Theorem~\ref{thm:config-space-bv} identifies this bracket with the
configuration-space residue pairing. Theorem~\ref{thm:bv-functor}
supplies a bracket-compatible Verdier comparison, and
Theorem~\ref{thm:verdier-bar-cobar} identifies that comparison with
\[
\mathbb{D}_{\operatorname{Ran}}(\barB^{\mathrm{ch}}(\cA))
\xrightarrow{\;\sim\;}
\cA^!_\infty
\]
on the Koszul locus. By the compatibility clause in
Theorem~\ref{thm:bv-functor}, the adjoint map
$x \mapsto \{x,-\}_{\mathrm{BV}}$ is exactly this Verdier comparison,
shifted by~$[1]$. Since the comparison is a quasi-isomorphism, the
BV Poisson structure is non-degenerate. When one passes to the
underlying complex of~$\cA^!_\infty$, this recovers the equivalent
description in terms of $\barB^{\mathrm{ch}}(\cA^!)$.

\emph{Part (ii).}
Shifting $\barB^{(g)}(\cA)$ by $[1]$ converts the degree~$+1$ BV
bracket to a degree~$0$ Lie bracket on $L_g$; the dg Lie algebra
axioms are inherited from the BV algebra axioms. The BV pairing on
$\barB^{(g)}(\cA)$ produced in Part~(i) restricts to a degree-$+1$
Verdier pairing. Evaluating against the canonical Verdier pairing on
the dual object and then shifting by~$[1]$ on both inputs produces a
bilinear form
\[
\omega_g \colon L_g \otimes L_g \longrightarrow \mathbb{C}[-1].
\]
Its non-degeneracy is equivalent to the quasi-isomorphism of the
adjoint map from Part~(i).
Invariance of the pairing (the cyclic property
$\langle [x,y], z \rangle = \langle x, [y,z] \rangle$)
is exactly the bracket-compatibility built into the conditional BV
package of Theorem~\ref{thm:bv-functor}. By the
Kontsevich--Pridham correspondence, a non-degenerate invariant pairing
of degree~$n$ on a dg Lie algebra yields an $n$-shifted symplectic
structure on its Maurer--Cartan formal moduli problem~\cite{Pridham17}.
Therefore $\mathrm{Def}_g(\cA) = \mathrm{MC}(L_g)$ is
$(-1)$-shifted symplectic.

\emph{Part (iii).}
Assume $g \geq 1$ and the transport square of
Lemma~\ref{lem:bar-chart-lagrangian-lift}. Proposition~\ref{prop:lagrangian-eigenspaces}
provides the ambient Verdier Lagrangian polarization of
$\mathbf{C}_g(\cA)$, and
Lemma~\ref{lem:bar-chart-lagrangian-lift} pulls that polarization back
to the $\sigma$-eigenspaces of~$L_g$. Thus
$L_g^+$ and~$L_g^-$ are complementary Lagrangians, and their
cohomology groups are $Q_g(\cA)$ and $Q_g(\cA^!)$.
\end{proof}

\begin{remark}[Diagonal scalar locus]
\index{uniform-weight!minimal scalar condition}
\index{cross-channel correction!modified pairing problem}
The scalar lane of the BV package does not stop at the
uniform-weight hypothesis. Theorem~\ref{thm:multi-weight-genus-expansion}
shows that for every modular Koszul algebra
\[
F_g(\cA)
\;=\;
\kappa(\cA)\cdot\lambda_g^{\mathrm{FP}}
\;+\;
\delta F_g^{\mathrm{cross}}(\cA),
\]
with $\delta F_1^{\mathrm{cross}}(\cA)=0$ for all~$\cA$ and
$\delta F_g^{\mathrm{cross}}(\cA)=0$ for all~$g$ on the
uniform-weight lane. Therefore the exact scalar condition for the
genus-$g$ BV pairing to reduce to the diagonal
$\kappa(\cA)\lambda_g^{\mathrm{FP}}$ term is
\[
\delta F_g^{\mathrm{cross}}(\cA)=0.
\]

The bar-side pairing used in
Theorem~\ref{thm:shifted-symplectic-complementarity} is this diagonal
pairing. It proves the scalar form of~\textup{(C2)} exactly on the
locus above. For multi-weight classes with
$\delta F_g^{\mathrm{cross}}(\cA)\neq0$, the diagonal pairing is not
the all-weight pairing; the missing datum is a mixed
Verdier-anti-invariant form indexed by the same mixed-channel boundary
graphs as the correction term.
\end{remark}

\begin{lemma}[Finite-window complementarity and completed limit;
\ClaimStatusConditional]
\label{lem:finite-window-complementarity}
\index{finite-window completion!complementarity}
\index{completion!Verdier eigenspaces}
Let $\mathbf{C}_g(\cA)$ be the strict flat ambient complex of
Theorem~\ref{thm:quantum-complementarity-main}, equipped with a
complete separated decreasing filtration $F^\bullet$ by conformal
weight, bar length, or shadow order. Assume $\sigma(F^N)\subset F^N$
for all~$N$. For every finite window
\[
  \mathbf{C}_g^{\leq N}(\cA)
  :=
  \mathbf{C}_g(\cA)/F^{N+1}\mathbf{C}_g(\cA)
\]
the projectors
$p^\pm=\tfrac12(\mathrm{id}\pm\sigma)$ descend and give a direct
sum of complexes
\[
  \mathbf{C}_g^{\leq N}(\cA)
  =
  \mathbf{Q}_{g,\leq N}(\cA)
  \oplus
  \mathbf{Q}_{g,\leq N}(\cA^!).
\]
If the Verdier pairing is perfect modulo $F^{N+1}$ and
anti-invariant under $\sigma$, then the two finite-window summands
are complementary Lagrangians in
$\mathbf{C}_g^{\leq N}(\cA)$. If the tower
$\{\mathbf{C}_g^{\leq N}(\cA)\}_N$ satisfies the Mittag--Leffler
condition on cohomology, then
\[
  \mathbf{C}_g(\cA)
  \simeq
  \varprojlim_N \mathbf{C}_g^{\leq N}(\cA)
  =
  \varprojlim_N \mathbf{Q}_{g,\leq N}(\cA)
  \oplus
  \varprojlim_N \mathbf{Q}_{g,\leq N}(\cA^!)
\]
and the completed Lagrangian statement is the inverse limit of the
finite-window statements.
\end{lemma}

\begin{proof}
The filtration is preserved by~$\sigma$, so the involution and the
idempotents $p^\pm$ act on every quotient
$\mathbf{C}_g^{\leq N}(\cA)$. Lemma~\ref{lem:involution-splitting}(a)
applied to the quotient gives the direct-sum decomposition. If the
pairing is perfect modulo $F^{N+1}$ and
$\langle\sigma x,\sigma y\rangle_{\mathbb{D}}
=-\langle x,y\rangle_{\mathbb{D}}$,
Lemma~\ref{lem:involution-splitting}(c) gives isotropy and maximality
of the two quotient eigenspaces. The
Mittag--Leffler hypothesis kills the derived inverse-limit term
$\varprojlim^1$ on cohomology, so passage to
$\varprojlim_N$ preserves both the splitting and the perfect
cross-pairing.
\end{proof}

\begin{remark}[Genus-wise strictification]
\label{rem:genuswise-strictification}
\index{strictification!genus-wise}
The dg~Lie algebra $L_g = \barB^{(g)}(\cA)[1]$ appearing in
Theorem~\ref{thm:shifted-symplectic-complementarity} is a genus-$g$
strict chart inside the modular deformation package of
\S\ref{subsec:modular-deformation-complex}. More precisely, after
filtering $\operatorname{Def}^{\mathrm{mod}}(\cA)$ by genus, $L_g$
identifies with the genus-$g$ strict truncation of the complete
cyclic $L_\infty$-object $\Definfmod(\cA)$
(Remark~\ref{rem:deformation-complex-linfty}). Thus the
$(-1)$-shifted symplectic formal moduli problem $\mathrm{MC}(L_g)$ is
one conditional coordinate chart on the modular formal moduli problem
controlled by $\Definfmod(\cA)$.
\end{remark}

\begin{proposition}[PTVV Lagrangian embedding; \ClaimStatusConditional]
\label{prop:ptvv-lagrangian}
\index{PTVV!Lagrangian embedding|textbf}
\index{Lagrangian!PTVV sense|textbf}
Assume the cochain-level complementarity decomposition and the
non-degenerate Verdier pairing of
Theorem~\textup{\ref{thm:quantum-complementarity-main}}. For a chiral
Koszul pair $(\cA, \cA^!)$ and $g \geq 2$, the cochain
complex
$C_g := R\Gamma(\overline{\mathcal{M}}_g, \mathcal{Z}(\cA))$
carries a $(-(3g{-}3))$-shifted symplectic structure in the sense of
Pantev--To\"en--Vaqui\'e--Vezzosi~\cite{PTVV13}, and
\(\mathbf{Q}_g(\cA)\), \(\mathbf{Q}_g(\cA^!)\) embed as complementary
Lagrangians in the
PTVV sense.
\end{proposition}

\begin{proof}
\emph{Shifted symplectic structure.}
For a cochain complex $V$ with a non-degenerate symmetric pairing
$\langle -,- \rangle \colon V \otimes V \to \mathbb{C}[n]$,
the associated linear derived affine scheme
$\mathbf{V} := \operatorname{Spec}(\operatorname{Sym}(V^*))$
carries an $n$-shifted symplectic structure: the pairing defines a
constant $2$-form of degree~$n$ on $\mathbf{V}$, and the closure
condition $d_{\mathrm{dR}}\omega = 0$ is automatic for a constant
form on a linear space~\cite[Example~1.4]{PTVV13}.
Non-degeneracy of the pairing translates to non-degeneracy of the
$2$-form.
The Verdier pairing on $C_g$
(Proposition~\ref{prop:lagrangian-eigenspaces}(i)) is a
non-degenerate symmetric pairing of degree $-(3g{-}3)$, so $C_g$
carries a $(-(3g{-}3))$-shifted symplectic structure.

\emph{PTVV Lagrangian conditions.}
A \emph{Lagrangian} in a linear $n$-shifted symplectic space
$(V, \omega)$ is a subcomplex $L \hookrightarrow V$
satisfying~\cite[Definition~2.8]{PTVV13}:
\begin{enumerate}[label=(\alph*)]
\item \emph{Isotropy:} $\omega|_L = 0$ (the pairing restricts to
 zero on $L$).
\item \emph{Non-degeneracy:} the induced map
 $L \to (V/L)^{\vee}[n]$ is a quasi-isomorphism.
\end{enumerate}
For $L = \mathbf{Q}_g(\cA)$:
\begin{itemize}
\item Condition~(a) is
 Proposition~\ref{prop:lagrangian-eigenspaces}(iii):
 $\langle v, w \rangle_{\mathbb{D}} = 0$ for
 $v, w \in \mathbf{Q}_g(\cA)$.
\item Condition~(b): Since
 $C_g = \mathbf{Q}_g(\cA) \oplus \mathbf{Q}_g(\cA^!)$
 (Theorem~\ref{thm:quantum-complementarity-main}), we have
 $C_g / \mathbf{Q}_g(\cA) \cong \mathbf{Q}_g(\cA^!)$, and the induced map
 $\mathbf{Q}_g(\cA) \to
 \mathbf{Q}_g(\cA^!)^{\vee}[-(3g{-}3)]$ is a
 quasi-isomorphism: the Verdier pairing restricted to
 $\mathbf{Q}_g(\cA) \otimes \mathbf{Q}_g(\cA^!)$ is perfect of degree $-(3g{-}3)$
 by~\eqref{eq:quantum-duality}.
\end{itemize}
The same argument applies with $\cA$ and~$\cA^!$ exchanged.
\end{proof}

\begin{remark}[Two shifted symplectic structures]\label{rem:ptvv-relation}
\index{Pantev--To\"en--Vaqui\'e--Vezzosi}
Complementarity carries two shifted symplectic structures:
(i)~Verdier on the ambient complex
($(-(3g{-}3))$-shifted,
Proposition~\ref{prop:ptvv-lagrangian});
(ii)~ambient cyclic deformation theory
($(-1)$-shifted,
Theorem~\ref{thm:ambient-complementarity-fmp}).
Conditional on Theorem~\ref{thm:shifted-symplectic-complementarity},
there is also a direct bar-side $(-1)$-shifted strict chart on
$L_g = \barB^{(g)}(\cA)[1]$. In each proved incarnation,
$Q_g(\cA)$ and $Q_g(\cA^!)$ are complementary Lagrangians; the
different realizations are compatible via the spectral sequence
$E_2^{p,q,g}\Rightarrow H^{p+q}_{\mathrm{chiral}}$.
\end{remark}

\begin{theorem}[External genus grading of the modular flat bar object; \ClaimStatusProvedHere]\label{thm:ss-genus-stratification}
\textup{[Regime: curved-central; all genera
\textup{(}Convention~\textup{\ref{conv:regime-tags})}.]}

The completed modular flat bar object has an external arithmetic-genus
grading
\[
\barB^{\mathrm{mod}}_{\mathrm{flat}}(\cA)
=
\prod_{g\geq0}\barB^{(g)}_{\mathrm{flat}}(\cA),
\qquad
F^{\leq G}\barB^{\mathrm{mod}}_{\mathrm{flat}}(\cA)
=
\prod_{0\leq g\leq G}\barB^{(g)}_{\mathrm{flat}}(\cA).
\]
For each fixed~$g$, the relative Leray filtration of
Theorem~\textup{\ref{thm:ss-quantum}} gives
\[
E_1^{p,q,g}
=
H^q\!\left(\barB^p_{g,\mathrm{flat}}(\cA)\right)
\Longrightarrow
H^{p+q}\!\left(\barB^{(g)}_{\mathrm{flat}}(\cA)\right).
\]
The spectral sequence is genuswise: the displayed differentials do
not change~$g$, and the product over~$g$ is assembled only after the
fixed-genus convergence statement.
\end{theorem}

\begin{proof}
For fixed~$g$ and~$p$, the relative configuration space is
\[
\pi_{g,p}\colon
\overline{\operatorname{Conf}}_{g,p}(X)
\longrightarrow
\overline{\mathcal M}_g,
\]
whose fiber over a stable curve~$\Sigma$ is the
Fulton--MacPherson compactification
$\overline{\operatorname{Conf}}_p(\Sigma)$. Thus the $p$-point
summand of the flat bar family is
\[
\barB^p_{g,\mathrm{flat}}(\cA)
=
R\Gamma\!\left(
\overline{\mathcal M}_g,\,
R\pi_{g,p*}
\bigl(\cA^{\boxtimes p}\otimes\Omega^\bullet_{\log},\Dg{g}\bigr)
\right).
\]
The Leray filtration of this proper morphism gives the displayed
fixed-genus spectral sequence.

Collision residues contract marked points on the same stable curve.
Separating clutching replaces one vertex by two vertices with
$g_1+g_2+h^1(\Gamma)=g$, and non-separating clutching lowers a
vertex genus by one while raising the graph loop number by one. In
both cases the arithmetic genus is preserved. Hence the differential,
coproduct, and stable-graph operations preserve
$F^{\leq G}$, and no differential in the fixed-genus Leray spectral
sequence shifts~$g$.
\end{proof}

\begin{remark}[Loop grading]\label{rem:ss-feynman}
\index{Feynman transform}
\index{modular operad!Feynman transform}
The external genus grading is the modular-operadic loop grading.
The Feynman transform of a modular operad preserves arithmetic
genus under both separating and non-separating clutching
\cite{Kon99}. The fixed-genus spectral sequences above are therefore
the graded pieces of one completed modular bar object.
\end{remark}

\begin{remark}[Genuswise convergence]\label{rem:genus-spectral-connection}
The differential $d_r$ in the Leray spectral sequence changes Leray
degree, not arithmetic genus. Convergence is controlled genus by
genus by the finite cohomological dimension of
$\overline{\mathcal{M}}_g$
(Lemma~\ref{lem:quantum-ss-convergence}; compare
Costello--Gwilliam~\cite{CG17}, Volume~2, Chapter~5).
\end{remark}

\subsection{Corollaries and physical interpretation}

\begin{conjecture}[Physical comparison datum; \ClaimStatusConjectured]
\label{conj:physical-complementarity}
Let $(\cA,\cA^!)$ be a chiral Koszul pair satisfying the hypotheses of
Theorem~\ref{thm:quantum-complementarity-main}. A physical
realization of this pair consists of the following extra data:
\begin{enumerate}[label=\textup{(\roman*)}]
\item a BRST complex $\mathcal B_{\mathrm{phys}}(\cA)$ and a
 quasi-isomorphism from the strict flat bar representative
 $\barB^{(g)}_{\mathrm{flat}}(\cA)$ to
 $\mathcal B_{\mathrm{phys}}(\cA)$, compatible with sewing;
\item a trace functional on BRST cohomology agreeing with the
 Verdier trace on $\mathbf C_g(\cA)$;
\item a second BRST model for $\cA^!$ carrying the Verdier-dual
 trace.
\end{enumerate}
Under these comparison data, the scalar central extension of the
Heisenberg algebra $\mathcal H_\kappa$ is carried to the curved
Koszul-dual term
\[
m_0^! = -\kappa\,\omega_{\mathrm{cen}},
\]
and the genus-$g$ quantum corrections split as the two Verdier
eigenspaces
$\mathbf Q_g(\cA)\oplus\mathbf Q_g(\cA^!)$ of
Theorem~\ref{thm:quantum-complementarity-main}. In gauge-theoretic
examples, the conjectural electric/magnetic language is the assertion
that the two BRST trace models above are the two physical sectors of
the same Verdier pairing.
\end{conjecture}

\begin{remark}[Scope of Conjecture~\ref{conj:physical-complementarity}]
The algebraic complementarity theorem is independent of the physical
dictionary. The missing input is the BRST comparison map with trace:
at genus~$0$ this is available in the Beilinson--Drinfeld chiral
setting, while at $g\geq1$ the flat bar complex is constructed
(Theorems~\ref{thm:inductive-genus-determination} and
\ref{thm:general-hs-sewing}) but its identification with the physical
BRST complex is exactly Conjecture~\ref{conj:master-bv-brst}. The
Kapustin--Witten example additionally requires constructing the
topological twist as an actual chiral Koszul pair. The Heisenberg and
Kac--Moody scalar checks are the explicit algebraic instances in
Examples~\ref{ex:heisenberg-complementarity-explicit} and
\ref{ex:kac-moody-complementarity-explicit}.
\end{remark}

\begin{corollary}[Mapping-class equivariance; \ClaimStatusConditional]
\label{cor:modular-properties}
Assume the center local system $\mathcal Z(\cA)$ and the
Verdier--Koszul involution $\sigma$ in
Theorem~\ref{thm:quantum-complementarity-main} are equivariant for the
mapping-class-group monodromy on the smooth moduli stack
$\mathcal M_g$, with quasi-unipotent extension across the boundary of
$\overline{\mathcal M}_g$. Then the decomposition
\[
H^*(\overline{\mathcal M}_g,\mathcal Z(\cA))
=Q_g(\cA)\oplus Q_g(\cA^!)
\]
is preserved by monodromy. On the Siegel/Hodge projection this gives
compatibility with the induced $\operatorname{Sp}(2g,\mathbb Z)$
action on $R^1\pi_*\mathbb Z$ and on the Hodge classes
$\lambda_i=c_i(\mathbb E_g)$.
\end{corollary}

\begin{proof}[Proof of Corollary~\ref{cor:modular-properties}]
The mapping class group acts by monodromy on the universal curve over
$\mathcal M_g$ and hence on any equivariant center local system. By
hypothesis this action commutes with the cochain representative of
the Verdier involution~$\sigma$. Therefore monodromy preserves the
two eigenspaces $\ker(\sigma-\mathrm{id})$ and
$\ker(\sigma+\mathrm{id})$, which are
$Q_g(\cA)$ and $Q_g(\cA^!)$ by
Theorem~\ref{thm:quantum-complementarity-main}. The quasi-unipotent
boundary extension lets the same argument pass to
$\overline{\mathcal M}_g$ through the Deligne canonical extension.
The projection to $R^1\pi_*\mathbb Z$ is the standard symplectic
monodromy representation, so the induced action factors through
$\operatorname{Sp}(2g,\mathbb Z)$ on the Hodge bundle and its Chern
classes.
\end{proof}

\begin{corollary}[Uniqueness of quantum corrections; \ClaimStatusConditional]
\label{cor:uniqueness-quantum}
Given genus-$g$ corrections with $g \geq 1$ for a chiral algebra
$\mathcal{A}$, the
Koszul dual corrections $Q_g(\mathcal{A}^!)$ are \emph{uniquely determined} by:
\begin{equation}
Q_g(\mathcal{A}^!) \cong \left(H^*(\overline{\mathcal{M}}_g, Z(\mathcal{A})) / Q_g(
\mathcal{A})\right)^\vee
\end{equation}
where the dual is taken with respect to Verdier duality.

This identification is constructive.
\end{corollary}

\begin{proof}[Proof of Corollary~\ref{cor:uniqueness-quantum}]
By the direct sum property (Lemma~\ref{lem:trivial-intersection-complete}) and exhaustion
(Lemma~\ref{lem:exhaustion-complete}), we have:
\begin{equation}
H^*(\overline{\mathcal{M}}_g, Z(\mathcal{A})) = Q_g(\mathcal{A}) \oplus Q_g(\mathcal{A}^!)
\end{equation}

Thus:
\begin{equation}
Q_g(\mathcal{A}^!) = H^*(\overline{\mathcal{M}}_g, Z(\mathcal{A})) / Q_g(\mathcal{A})
\end{equation}
as vector spaces.

By Verdier duality (Corollary~\ref{cor:quantum-dual-complete}):
\begin{equation}
Q_g(\mathcal{A}^!) \cong Q_g(\mathcal{A})^\vee
\end{equation}

Combining these gives the stated formula.

The algorithm is: compute $H^*(\overline{\mathcal{M}}_g, Z(\mathcal{A}))$, compute $Q_g(\mathcal{A})$ via the bar spectral sequence, and take the Verdier-orthogonal complement.
See Examples \ref{ex:heisenberg-complementarity-explicit} and \ref{ex:kac-moody-complementarity-explicit}.
\end{proof}

\begin{corollary}[Vanishing results; \ClaimStatusProvedHere]
\label{cor:vanishing-quantum}
If $\mathcal{A}$ has no quantum corrections at genus $g$, meaning $Q_g(\mathcal{A}) = 0$,
then:
\begin{equation}
Q_g(\mathcal{A}^!) \cong H^*(\overline{\mathcal{M}}_g, Z(\mathcal{A}))
\end{equation}

Conversely, if \emph{both} $Q_g(\mathcal{A}) = 0$ and $Q_g(\mathcal{A}^!) = 0$, then:
\begin{equation}
H^*(\overline{\mathcal{M}}_g, Z(\mathcal{A})) = 0
\end{equation}
meaning the center acts trivially on moduli space cohomology.
\end{corollary}

\begin{proof}[Proof of Corollary~\ref{cor:vanishing-quantum}]
\emph{First statement}: By the decomposition theorem:
\begin{equation}
H^*(\overline{\mathcal{M}}_g, Z(\mathcal{A})) = Q_g(\mathcal{A}) \oplus Q_g(\mathcal{A}^!)
\end{equation}

If $Q_g(\mathcal{A}) = 0$, then $Q_g(\mathcal{A}^!) \cong H^*(\overline{\mathcal{M}}_g,
Z(\mathcal{A}))$.

\emph{Second statement}: If both vanish, then by exhaustion:
\begin{equation}
0 = \dim Q_g(\mathcal{A}) + \dim Q_g(\mathcal{A}^!) = \dim H^*(\overline{\mathcal{M}}_g,
Z(\mathcal{A}))
\end{equation}

Thus $H^*(\overline{\mathcal{M}}_g, Z(\mathcal{A})) = 0$.
\end{proof}

\begin{remark}[Examples of vanishing]
At genus~$0$: $Q_0(\cA)=\mathbb{C}$, $Q_0(\cA^!)=0$ (tree-level,
no obstruction).
For TFTs: $Q_g(\cA)=0$ for all~$g$.
For $bc$--$\beta\gamma$: opposite obstructions cancel in the
complementarity sum (Theorem~\ref{thm:fermion-all-genera}), but
individual $Q_g$ need not vanish.
\end{remark}

\subsection{Self-dual algebras and critical level}

\begin{theorem}[Balanced self-dual halving; \ClaimStatusConditional]\label{thm:self-dual-halving}
\index{quantum corrections!self-dual halving}
\textup{[Regime: curved-central on the Koszul locus; all genera
\textup{(}Convention~\textup{\ref{conv:regime-tags})}.]}

Let $(\cA, \cA^!)$ be a Koszul chiral pair
(Definition~\ref{def:chiral-koszul-pair}), and fix $g \geq 1$ in
the ordinary flat regime of Theorem~\ref{thm:quantum-complementarity-main}.
Let $\sigma$ be the represented Verdier involution on
\[
 \mathbf C_g(\cA)=R\Gamma(\overline{\mathcal M}_g,\mathcal Z(\cA)).
\]
Assume that a self-duality isomorphism
$f\colon \cA\xrightarrow{\sim}\cA^!$ induces a cochain automorphism
$\tau_f$ of $\mathbf C_g(\cA)$, after the center comparison, with
\begin{equation}\label{eq:balanced-self-dual-antiequivariance}
 \tau_f\sigma=-\sigma\tau_f .
\end{equation}
Then
\begin{equation}\label{eq:self-dual-halving}
\dim Q_g(\cA) \;=\; \dim Q_g(\cA^!)
\;=\; \tfrac{1}{2}\,\dim H^*(\overline{\mathcal{M}}_g, Z(\cA))
\end{equation}
whenever the displayed cohomology is finite-dimensional. Without
\eqref{eq:balanced-self-dual-antiequivariance}, self-duality of the
underlying chiral algebra implies no halving statement. The defect is
the trace of the Verdier involution:
\begin{equation}\label{eq:self-dual-trace-defect}
\dim Q_g(\cA)-\dim Q_g(\cA^!)
=
\operatorname{Tr}\!\left(\sigma\mid
H^*(\mathbf C_g(\cA))\right).
\end{equation}
\end{theorem}

\begin{proof}
Let $V=H^*(\mathbf C_g(\cA))$. Theorem~\ref{thm:quantum-complementarity-main}
identifies
\[
Q_g(\cA)=V^+=\ker(\sigma-\mathrm{id}),\qquad
Q_g(\cA^!)=V^-=\ker(\sigma+\mathrm{id}).
\]
If $x\in V^+$, then
$\sigma(\tau_f x)=-\tau_f(\sigma x)=-\tau_f x$, hence
$\tau_f x\in V^-$. The same calculation with $V^-$ gives
$\tau_f(V^-)=V^+$, so $\tau_f$ identifies the two eigenspaces and
\[
\dim Q_g(\cA)=\dim Q_g(\cA^!).
\]
The direct-sum decomposition gives
\[
\dim Q_g(\cA) + \dim Q_g(\cA^!)
= \dim H^*(\overline{\mathcal{M}}_g, Z(\cA))
\]
and therefore the halving formula.

For the defect formula, use the eigenspace decomposition again:
\[
\operatorname{Tr}(\sigma\mid V)
=\dim V^+ - \dim V^-
=\dim Q_g(\cA)-\dim Q_g(\cA^!).
\]
An isomorphism $\cA\simeq\cA^!$ alone does not determine how the
induced automorphism interacts with~$\sigma$; the anti-equivariance
condition is the missing datum.
\end{proof}

\begin{corollary}[Virasoro quantum corrections under balanced self-duality; \ClaimStatusConditional]
\label{cor:virasoro-quantum-dim}
\index{Virasoro algebra!quantum corrections}
The Virasoro algebra at $c = 13$ is chiral Koszul self-dual
($\mathrm{Vir}_{13}^! = \mathrm{Vir}_{26-13} = \mathrm{Vir}_{13}$,
Proposition~\ref{prop:virasoro-generic-koszul-dual}).
Assume that this self-duality is represented on
$\mathbf C_g(\mathrm{Vir}_{13})$ by a balanced automorphism
$\tau_{13}$ satisfying
$\tau_{13}\sigma=-\sigma\tau_{13}$. Then for all $g \geq 1$:
\begin{equation}\label{eq:virasoro-quantum-dim}
\dim Q_g(\mathrm{Vir}_{13})
= \tfrac{1}{2}\,\dim H^*(\overline{\mathcal{M}}_g).
\end{equation}
Explicitly: $\dim Q_1(\mathrm{Vir}_{13}) = 1$ and
$\dim Q_2(\mathrm{Vir}_{13}) = 4$.
Note: $c = 0$ is the unique \emph{uncurved} point ($\kappa = 0$),
but $\mathrm{Vir}_0^! = \mathrm{Vir}_{26} \neq \mathrm{Vir}_0$,
so the halving formula does not apply there.
\end{corollary}

\begin{proof}
Apply Theorem~\ref{thm:self-dual-halving} with
$\cA = \mathrm{Vir}_{13}$ and the balanced representative
$\tau_{13}$. Since $Z(\mathrm{Vir}_{13}) = \mathbb{C}$,
$H^*(\overline{\mathcal{M}}_g, Z(\cA))
= H^*(\overline{\mathcal{M}}_g)$.
At genus~1: $\dim H^*(\overline{\mathcal{M}}_{1,1}) = 2$,
giving $\dim Q_1 = 1$. At genus~2:
$\dim H^*(\overline{\mathcal{M}}_2) = 8$
(Poincar\'e polynomial $1 + 3t^2 + 3t^4 + t^6$),
giving $\dim Q_2 = 4$.
\end{proof}

\begin{corollary}[Critical level scalar uncurving; \ClaimStatusConditional]
\label{cor:critical-uncurving}
\index{critical level!uncurving}
At the critical level $k = -h^\vee$, the uniform-weight scalar
obstruction class vanishes for all $g \geq 1$:
\begin{equation}\label{eq:critical-uncurving}
\mathrm{obs}^{\mathrm{scal}}_g(V_{-h^\vee}(\mathfrak g)) = 0 .
\end{equation}
Equivalently, the scalar diagonal of the curved fiber identity has
zero coefficient,
\[
\dfib^{\,2}\big|_{\mathrm{scal}}
=
\kappa(V_{-h^\vee}(\mathfrak g))\,\omega_{\mathbb E}=0 .
\]
If, in addition, the ordinary flat representative of
Theorem~\textup{\ref{thm:quantum-complementarity-main}} is available
and the Feigin--Frenkel center gives the center local system, then
the critical-level ambient complex admits the Verdier eigenspace
decomposition
\begin{equation}\label{eq:critical-complementarity}
Q_g(V_{-h^\vee}(\mathfrak g))
\;\oplus\;
Q_g(V_{-h^\vee}(\mathfrak g)^!)
\;=\;
H^*\!\bigl(\overline{\mathcal{M}}_g,\,
\mathrm{Fun}(\mathrm{Op}_{\mathfrak{g}^\vee}(D))\bigr)
\end{equation}
where $\mathrm{Fun}(\mathrm{Op}_{\mathfrak{g}^\vee}(D))$ is the
Feigin--Frenkel center
(Remark~\ref{rem:feigin-frenkel-center}; \cite{Feigin-Frenkel}).
The scalar vanishing alone does not identify the full curved bar
object with a strict cochain complex, does not compute the individual
eigenspaces, and does not identify the Koszul dual with a
$\mathcal W$-algebra.
\end{corollary}

\begin{proof}
From Table~\ref{tab:obstruction-summary}, the genus-$g$ obstruction for
$V_k(\mathfrak g)$ on the uniform-weight scalar lane is
$\mathrm{obs}_g = \frac{(k+h^\vee)\dim(\mathfrak{g})}{2h^\vee}\,
\lambda_g$ \textup{(}UNIFORM-WEIGHT\textup{)}.
At $k = -h^\vee$, the coefficient $k + h^\vee$ vanishes, proving
\eqref{eq:critical-uncurving}. The displayed curvature identity is
the degree-$2$ scalar diagonal of
Proposition~\ref{prop:hodge-curvature-scalar-projection}. Under the
additional center-local-system and Verdier-involution hypotheses, the
decomposition~\eqref{eq:critical-complementarity} is exactly
Theorem~\ref{thm:quantum-complementarity-main} applied to the
critical-level ambient complex.
\end{proof}

\begin{remark}[Significance]\label{rem:self-dual-significance}
Self-duality and balanced halving are distinct conditions. A fixed
point of the Koszul involution gives $\cA\simeq\cA^!$; halving
requires an anti-equivariant representative
$\tau\sigma=-\sigma\tau$ on the genus-$g$ ambient complex. At the
critical affine level, the scalar curvature coefficient vanishes and
the center local system is governed by the Feigin--Frenkel center.
The Verdier splitting of
$H^*(\overline{\mathcal{M}}_g,
\mathrm{Fun}(\mathrm{Op}_{\mathfrak{g}^\vee}))$ is available only
after the hypotheses of Theorem~\ref{thm:quantum-complementarity-main}
are imposed. The fixed-point property by itself is a scalar
uncurving statement, not a halving theorem. Different balanced
representatives give different Verdier splittings of the same
cohomology group.
\end{remark}

\begin{conjecture}[String theory interpretation; \ClaimStatusConjectured]
\label{conj:string-theory-complementarity-explicit}
In topological string theory, the complementarity theorem explains:
\begin{itemize}
\item \emph{A-model/B-model duality}: The A-model chiral algebra and B-model chiral
algebra are Koszul dual, with quantum corrections satisfying complementarity.

\item \emph{Large $N$ duality}: At large $N$ (genus expansion parameter), the
planar ($g=0$) contributions of one theory match the non-planar ($g \geq 1$)
contributions of the dual theory.

\item \emph{Gopakumar--Vafa invariants}: The generating function for Gopakumar--Vafa
invariants packages both $Q_g(\mathcal{A})$ and $Q_g(\mathcal{A}^!)$ into a single
modular form.
\end{itemize}
\end{conjecture}

\begin{remark}[Scope of Conjecture~\ref{conj:string-theory-complementarity-explicit}]
Labeled \ClaimStatusHeuristic{} because: (a)~A/B-model Koszul pair
identification requires~\cite{CL16}; (b)~large~$N$ duality is
physics input; (c)~Gopakumar--Vafa integrality is open.
The mathematical content is that Theorem~\ref{thm:quantum-complementarity-main}
applied to the conjectural Koszul pair yields these identifications.
\end{remark}

\subsection{Explicit examples of complementarity}

\begin{example}[Heisenberg algebra: complete genus-1 computation]
\label{ex:heisenberg-complementarity-explicit}

\emph{Setup}: The Heisenberg algebra $\mathcal{H}_\kappa$ at level $\kappa$ has:
\begin{equation}
[a_m, a_n] = m\delta_{m+n,0} \kappa
\end{equation}

The center is $Z(\mathcal{H}_\kappa) = \mathbb{C} \cdot \mathbb{1} \oplus \mathbb{C}
\cdot \kappa$.

\emph{Genus 1 moduli space}: The cohomology ring of $\overline{\mathcal{M}}_{1,1}$ is generated by $\lambda = c_1(\mathbb{E})$:
\begin{equation}
H^*(\overline{\mathcal{M}}_{1,1}) = \mathbb{Q}[\lambda] / (\lambda^2) = \mathbb{Q}
\oplus \mathbb{Q}\lambda
\end{equation}

\emph{Step 1: Compute $Q_1(\mathcal{H}_\kappa)$}.

The genus-1 bar complex is:
\begin{equation}
\bar{B}^{(1)}(\mathcal{H}_\kappa) = \bigoplus_{n \geq 0} \Gamma(\overline{C}_n(E_\tau),
\mathcal{H}_\kappa^{\boxtimes n} \otimes \Omega^*_{\log})
\end{equation}
where $E_\tau$ is the elliptic curve with modulus $\tau$.

The differential has a genus-1 correction using the elliptic propagator:
\begin{equation}
d^{(1)} = \sum_{i<j} \operatorname{Res}_{D_{ij}} \circ K^{(1)}_{ij}
\end{equation}
where $K^{(1)}(z,w|\tau) = \theta_1'(z-w|\tau)/\theta_1(z-w|\tau)$ is the genus-$1$ chiral propagator (Remark~\ref{rem:wp-config-space}).

The failure of this \emph{fiberwise} collision differential $d_{\mathrm{fib}}$
to square to zero is measured by the quasi-periodicity of the propagator
around the $B$-cycle
(cf.\ Remark~\ref{rem:differential-notation} for the distinction between
$d_{\mathrm{fib}}$ and the total differential $d^{(1)}_{\mathrm{total}}$,
which \emph{does} square to zero by Theorem~\ref{thm:genus1-d-squared}).
The propagator satisfies $K^{(1)}(z+\tau) = K^{(1)}(z) - 2\pi i$, so:
\begin{equation}
(d_{\mathrm{fib}})^2 = \kappa \cdot \omega_1 \cdot \operatorname{id}
\end{equation}
where $\omega_1 \in H^2(\overline{\mathcal{M}}_{1,1})$ is the fundamental class (the $B$-cycle monodromy of the propagator, valued in $\mathbb{C} \cdot \kappa$).

This curvature identifies the obstruction coefficient
$\kappa(\mathcal{H}_\kappa) = \kappa$.
Taking cohomology of the total differential
$d^{(1)}_{\mathrm{total}}$:
\begin{equation}
Q_1(\mathcal{H}_\kappa) = H^*(\bar{B}^{(1)}(\mathcal{H}_\kappa),\,
d^{(1)}_{\mathrm{total}}) = \mathbb{C} \cdot \kappa
\end{equation}

\emph{Step 2: Compute $Q_1(\mathcal{H}_\kappa^!)$ using complementarity}.

The Koszul dual of the Heisenberg algebra is the curved Chevalley--Eilenberg algebra (cf.\ Part~\ref{part:characteristic-datum}):
\begin{equation}
\mathcal{H}_\kappa^! = \mathrm{CE}(\mathfrak{h}_{-\kappa}) = (\mathrm{Sym}^{\mathrm{ch}}(V^*), d_{\mathrm{CE}}, m_0 = -\kappa \cdot c)
\end{equation}
where $V^*$ is the linear dual of the generating space~$V$, the curvature
$m_0 = -\kappa \cdot c$ encodes the obstruction, and
$\kappa(\mathcal{H}_\kappa^!) = -\kappa$
(since $h^\vee = 0$ for the Heisenberg).
Note: $\mathcal{H}_\kappa^! = \mathrm{Sym}^{\mathrm{ch}}(V^*)
\neq \mathcal{H}_{-\kappa}$; the Koszul dual is generated by~$V^*$,
not by~$V$ at level~$-\kappa$
\textup{(}Remark~\textup{\ref{rem:km-complementarity-vs-feigin-frenkel})}.

By the complementarity theorem:
\begin{equation}
Q_1(\mathcal{H}_\kappa^!) = \left(H^*(\overline{\mathcal{M}}_{1,1}, Z(\mathcal{H}_\kappa))
/ Q_1(\mathcal{H}_\kappa)\right)^\vee
\end{equation}

Since $H^*(\overline{\mathcal{M}}_{1,1}) = \mathbb{C} \oplus \mathbb{C}\lambda$ and
$Q_1(\mathcal{H}_\kappa) = \mathbb{C} \cdot \kappa$ (which pairs with $H^0 = \mathbb{C}$),
we have:
\begin{equation}
Q_1(\mathcal{H}_\kappa^!) = (\mathbb{C}\lambda)^\vee = \mathbb{C} \cdot \lambda^\vee
\end{equation}

\emph{Interpretation}: $Q_1(\mathcal{H}_\kappa) = \mathbb{C} \cdot \kappa$ (the central extension appears as an obstruction) and $Q_1(\mathcal{H}_\kappa^!) = \mathbb{C} \cdot \lambda$ (the first Chern class appears as a deformation).

Together they span:
\begin{equation}
Q_1(\mathcal{H}_\kappa) \oplus Q_1(\mathcal{H}_\kappa^!) = \mathbb{C} \oplus \mathbb{C}
= H^*(\overline{\mathcal{M}}_{1,1})
\end{equation}

\emph{Verification}: We can verify this directly by computing the cobar complex of
the curved commutative algebra $\mathrm{CE}(\mathfrak{h}_{-\kappa})$ and showing its genus-$1$ contributions are $\mathbb{C} \cdot \lambda$.
\end{example}

\begin{example}[Kac--Moody scalar lane at critical level]
\label{ex:kac-moody-complementarity-explicit}

\emph{Setup}: The affine Kac--Moody algebra $\widehat{\mathfrak{g}}_k$ at level $k$
has:
\begin{equation}
[J^a_m, J^b_n] = \sum_c f^{abc} J^c_{m+n} + m\delta_{m+n,0} k \delta^{ab}
\end{equation}

For the universal affine family the scalar modular characteristic is
\[
\kappa(V_k(\mathfrak g))
=
\frac{(k+h^\vee)\dim\mathfrak g}{2h^\vee}.
\]
The affine Koszul-dual level on the scalar lane is
$k'=-k-2h^\vee$, and hence
\[
\kappa(V_{k'}(\mathfrak g))
=
\frac{(k'+h^\vee)\dim\mathfrak g}{2h^\vee}
=
-\,\kappa(V_k(\mathfrak g)).
\]
At the critical value $k=-h^\vee$ both scalar coefficients vanish.
Thus the genus-$g$ scalar obstruction
\[
\operatorname{obs}^{\mathrm{scal}}_g(V_k(\mathfrak g))
=
\frac{(k+h^\vee)\dim\mathfrak g}{2h^\vee}\,\lambda_g
\]
vanishes at the critical level, and the same is true for the dual
scalar coefficient.

The center of the universal affine vertex algebra jumps at the
critical level: the Feigin--Frenkel center is
$\mathrm{Fun}(\mathrm{Op}_{\mathfrak g^\vee}(D))$. Consequently the
critical-level complementarity space, when the C0/C1 hypotheses are
available, is
\[
H^*\!\left(
\overline{\mathcal M}_g,\,
\mathrm{Fun}(\mathrm{Op}_{\mathfrak g^\vee}(D))
\right)
=
Q_g(V_{-h^\vee}(\mathfrak g))
\oplus
Q_g(V_{-h^\vee}(\mathfrak g)^!).
\]
This displayed equality is the Verdier eigenspace splitting of the
ambient center complex. It does not compute the individual
eigenspaces and does not identify them with a quadratic Casimir line
or a Hodge line without an explicit critical-level bar computation.

\begin{remark}[Distinction from Feigin--Frenkel duality]
\label{rem:km-complementarity-vs-feigin-frenkel}
\index{Feigin--Frenkel!distinction from Koszul duality}
The \emph{Koszul} dual branch appearing in complementarity is
obtained by Verdier duality on the bar coalgebra. The
$\mathcal{W}$-algebra in the Feigin--Frenkel theorem is obtained
from the affine algebra by quantum Drinfeld--Sokolov reduction. The
DS functor may intertwine the two structures under its own
hypotheses, but it is not the bar-cobar inversion functor and it is
not the definition of $\cA^!$ in Theorem~C.
\end{remark}
\end{example}

\begin{example}[\texorpdfstring{$\beta\gamma$}{beta-gamma} system: Koszul dual to free fermions]
\label{ex:betagamma-fermion-koszul-duality}

\emph{Setup}: The $\beta\gamma$ system (symplectic bosons) with OPE:
\begin{equation}
\beta(z)\gamma(w) \sim \frac{1}{z-w}
\end{equation}

This system is \emph{Koszul dual to free fermions}: $(\beta\gamma)^! \cong \mathcal{F}$,
where $\mathcal{F}$ is the free fermion chiral algebra with generator $\psi$ satisfying $\psi^2=0$.

\emph{Koszul duality} (following Gui--Li--Zeng~\cite{GLZ22}).

\begin{theorem}[Fermion-boson Koszul duality; \ClaimStatusProvedHere]
\label{thm:fermion-boson-koszul-hg}
The $\beta\gamma$ system and free fermions form a Koszul dual pair:
\begin{equation}
\mathcal{F}^! \cong \beta\gamma \quad \text{and} \quad (\beta\gamma)^! \cong \mathcal{F}
\end{equation}
\end{theorem}

\begin{proof}
Via bar construction and duality:
\begin{enumerate}
\item Bar complex: $\bar{B}(\mathcal{F}) = \Lambda^{*,c}(\psi^*, \partial\psi^*, \ldots)$ (the exterior bar coalgebra on the desuspended generators).
\item Koszul dual algebra: $\mathcal{F}^! \cong \beta\gamma$ is obtained by dualizing the dual coalgebra computed from $\bar{B}(\mathcal{F})$; in this free case that amounts to the linear dual of the bar coalgebra. Bar-cobar inversion gives $\Omega(\bar{B}(\mathcal{F})) \cong \mathcal{F}$ (recovers the fermion, not $\beta\gamma$).
\item Conversely: $\bar{B}(\beta\gamma)$ has nonzero cohomology classes $[\beta \otimes \beta], [\gamma \otimes \gamma] \in H^2(\bar{B}(\beta\gamma))$, which dualize to the fermionic relation $\psi^2 = 0$.
\item Koszul dual: $(\beta\gamma)^! = \bar{B}(\beta\gamma)^\vee \cong \mathcal{F}$.
\end{enumerate}
\end{proof}

\emph{Genus 1 computation}:
\begin{equation}
Q_1(\beta\gamma) = \mathbb{C} \cdot [\beta\gamma]
\end{equation}
where $[\beta\gamma]$ is the first descendant of the identity operator.

Since $(\beta\gamma)^! \cong \mathcal{F}$ (free fermions):
\begin{equation}
Q_1((\beta\gamma)^!) = Q_1(\mathcal{F}) = \mathbb{C} \cdot [\psi \partial\psi]
\end{equation}

By complementarity:
\begin{equation}
Q_1(\beta\gamma) \oplus Q_1(\mathcal{F}) = \mathbb{C} \cdot [\beta\gamma] \oplus \mathbb{C} \cdot [\psi \partial\psi]
\end{equation}

\emph{Physical interpretation.} The bosonization correspondence exchanges fermionic $\psi^2 = 0$ with symplectic bosonic $[\beta,\gamma] = 1$, and exterior algebra with symmetric-type algebra.

\emph{Explicit verification.} The partition function on $E_\tau$ is:
\begin{equation}
Z_{E_\tau}[\beta\gamma] = \frac{1}{\eta(\tau)^2}
\end{equation}

Expanding in $q = e^{2\pi i \tau}$:
\begin{equation}
Z_{E_\tau}[\beta\gamma] = q^{-1/12} (1 + 2q + 5q^2 + 10q^3 + \cdots)
\end{equation}
where the coefficients are the number of bicoloured partitions
$p_{-2}(n) = \sum_{k=0}^n p(k)p(n-k)$.
The leading $q^0$ coefficient ($= 1$) corresponds to $Q_1(\beta\gamma)$,
confirming $\dim Q_1(\beta\gamma) = 1$.
\end{example}

\subsection{Higher genus: genus 2 explicit computations}

\begin{example}[Heisenberg at genus 2]
\label{ex:heisenberg-genus-2-complementarity}

\emph{Setup}: For genus $g=2$, the moduli space has dimension $\dim \overline{
\mathcal{M}}_2 = 3$. The cohomology is:
\begin{equation}
H^*(\overline{\mathcal{M}}_2) = \mathbb{Q}[\lambda_1, \delta_0, \delta_1, \lambda_2] / (\text{relations})
\end{equation}
where $\lambda_1 \in H^2$ is the first Chern class of the Hodge bundle,
$\delta_0, \delta_1 \in H^2$ are the boundary divisor classes, and
$\lambda_2 \in H^4$ is the second Chern class of the Hodge bundle.

The Poincaré polynomial is:
\begin{equation}
P_t(H^*(\overline{\mathcal{M}}_2)) = 1 + 3t^2 + 3t^4 + t^6
\end{equation}

\emph{Genus-$2$ complementarity for Heisenberg.}
By Theorem~\ref{thm:quantum-complementarity-main},
\begin{equation}
Q_2(\mathcal{H}_\kappa) \oplus Q_2(\mathcal{H}_\kappa^!)
= H^*(\overline{\mathcal{M}}_2),
\qquad
\dim H^*(\overline{\mathcal{M}}_2) = 8.
\end{equation}
The individual dimensions
$\dim Q_2(\mathcal{H}_\kappa)$ and $\dim Q_2(\mathcal{H}_\kappa^!)$
are determined by the Verdier involution on
$H^*(\overline{\mathcal{M}}_2, Z(\mathcal{H}_\kappa))$; their
computation requires the explicit genus-$2$ bar complex
(\S\ref{subsec:genus2-complementarity-verification}).
The complementarity sum $\dim Q_2 + \dim Q_2^! = 8$ is
unconditional.
\end{example}

\begin{remark}[Open: non-abelian quantum corrections at higher genus]%
\label{rem:nonabelian-Qg-open}%
\index{quantum corrections!non-abelian open problem}%
The explicit $Q_g$ computations above are for the abelian
Heisenberg algebra, whose bar complex
$T^c(s^{-1}\bar{\mathcal H}_k)$ is the cofree exterior coalgebra
on the desuspended generators with no further relations. For
non-abelian algebras ($\widehat{\mathfrak{sl}}_{2,k}$,
principal $\mathcal{W}_N$, and lattice VOAs), the
complementarity decomposition $Q_g(\cA) \oplus Q_g(\cA^!)
= C_g(\cA)$ is \emph{proved} at all genera
(Theorem~\ref{thm:quantum-complementarity-main}),
but the individual dimensions $\dim Q_g(\cA)$ are
\emph{not computed} for $g \geq 2$.

The sharpest test case is $Q_1(\widehat{\mathfrak{sl}}_{2,k})$
as a function of the level~$k$: by comparison with the Verlinde
formula (equation~\eqref{eq:verlinde-general}), one expects
$\dim H^0(\barB^{(1)}(\widehat{\mathfrak{sl}}_{2,k}), \Dg{1})
= k{+}1$, algebra-dependent data invisible to the scalar
$\kappa = 3(k{+}2)/4$. This is Level~$1$ of the modular
invariant hierarchy (Remark~\ref{rem:four-levels}).
\end{remark}

\subsection{Algorithmic computation of quantum corrections}

We conclude with a practical algorithm for computing $Q_g(\mathcal{A})$ and verifying
complementarity.

\begin{remark}[Computing quantum corrections]\label{rem:quantum-corrections}
Given a chiral algebra $\mathcal{A}$ on a curve $X$ of genus $g$, the quantum correction space $Q_g(\mathcal{A})$ is computed as follows
(see Remark~\ref{rem:differential-notation} for notation).
\begin{enumerate}
\item Compute the center $Z(\mathcal{A}) = \{z \in \mathcal{A} : [z, a] = 0 \text{ for all } a\}$.
\item Construct the bar complex $\bar{B}^{(g)}(\mathcal{A}) = \bigoplus_n \Gamma(\overline{C}_n(X_g), \mathcal{A}^{\boxtimes n} \otimes \Omega^*_{\log})$.
\item Compute the total corrected differential $\Dg{g}$ on $\bar{B}^{(g)}(\mathcal{A})$
(Theorem~\ref{thm:genus-differential}), which includes fiberwise collision
residues, the Gauss--Manin connection, and modular corrections.
This differential satisfies $\Dg{g}^{\,2} = 0$ (Convention~\ref{conv:higher-genus-differentials}).
\item Compute the fiberwise collision differential $\dfib$ and verify
at chain level that $\dfib^{\,2} = \kappa(\mathcal{A}) \cdot \omega_g$;
this curvature measures the obstruction coefficient
$\kappa(\mathcal{A})$ and is interpreted coderivedly before passing to
the flat differential~$\Dg{g}$.
\item Take cohomology of the total differential:
$Q_g(\mathcal{A}) = H^*(\bar{B}^{(g)}(\mathcal{A}),\, \Dg{g})$.
\item Verify complementarity: $\dim Q_g(\mathcal{A}) + \dim Q_g(\mathcal{A}^!) = \dim H^*(\overline{\mathcal{M}}_g, \mathcal{Z}(\mathcal{A}))$.
\end{enumerate}
\end{remark}

\begin{example}[Algorithm applied to Heisenberg]
For $\mathcal{H}_\kappa$ at genus 1:
\begin{enumerate}
\item $Z(\mathcal{H}_\kappa) = \mathbb{C} \cdot \mathbb{1} \oplus \mathbb{C} \cdot
K$, where $K$ is the central element of the universal Heisenberg
(specializing to the level parameter~$\kappa$).
\item $\bar{B}^{(1)}(\mathcal{H}_\kappa) = \bigoplus_n \Gamma(\overline{C}_n(E_\tau),
\mathcal{H}_\kappa^{\boxtimes n} \otimes \Omega^*_{\log})$.
\item The total differential $d^{(1)}_{\mathrm{total}} = d_{\mathrm{residue}} +
d_{\mathrm{elliptic}} + d_{\mathrm{modular}}$ satisfies
$(d^{(1)}_{\mathrm{total}})^2 = 0$
(Theorem~\ref{thm:genus1-d-squared}).
\item The fiberwise collision differential
$d_{\mathrm{fib}} = \sum_{i<j} \operatorname{Res}_{D_{ij}} \cdot \eta_{ij}^{(1)}$ satisfies
$(d_{\mathrm{fib}})^2 = \kappa \cdot \omega_1 \neq 0$,
identifying the obstruction coefficient $\kappa(\mathcal{H}_\kappa) = \kappa$.
\item $Q_1(\mathcal{H}_\kappa) = H^*(\bar{B}^{(1)}(\mathcal{H}_\kappa),\,
d^{(1)}_{\mathrm{total}}) = \mathbb{C} \cdot \kappa$.
\item $\dim Q_1(\mathcal{H}_\kappa) + \dim Q_1(\mathcal{H}_\kappa^!) = 1 + 1 = 2 =
\dim H^*(\overline{\mathcal{M}}_{1,1})$.
\end{enumerate}
\end{example}


\begin{remark}[Connection to string theory]
\label{rem:string-theory-complementarity}
The complementarity theorem is the mathematical foundation for the
complementarity of Type A and Type B topological strings, the
exchange of quantum corrections under mirror symmetry, and the
modular constraints on genus expansions.
\end{remark}

\subsection{The complementarity landscape}
\label{subsec:complementarity-landscape}
\index{complementarity!landscape|textbf}
\index{Koszul conductor|textbf}
\index{anomaly ratio|textbf}

The complementarity sum $\kappa(\cA) + \kappa(\cA^!)$ is a
level-independent invariant for every standard family
(Theorem~\ref{thm:complementarity-root-datum}). We collect the
full data here.

\begin{proposition}[Complementarity landscape; \ClaimStatusConditional]
\label{prop:complementarity-landscape}
\index{complementarity!landscape}
Let $(\cA, \cA^!)$ be a chiral Koszul pair from the standard
landscape. The complementarity sum
$\kappa(\cA) + \kappa(\cA^!)$ is determined as follows.
\begin{enumerate}[label=\textup{(\roman*)}]
\item \emph{Free fields.}\;
 For the Heisenberg~$\cH_k$, the free fermion, the
 $\beta\gamma$~system at conformal weight~$\lambda$, and
 lattice vertex algebras~$V_\Lambda$:
 \begin{equation}\label{eq:free-field-antisymmetry}
 \kappa(\cA) + \kappa(\cA^!) \;=\; 0.
 \end{equation}
 In each case the Koszul dual inverts the sign of $\kappa$:
 $\kappa(\cH_k^!) = -k$,\;
 $\kappa(\beta\gamma_\lambda^!) = \kappa(bc_\lambda)
 = -(6\lambda^2 - 6\lambda + 1)$,\;
 $\kappa(V_\Lambda^!) = -\operatorname{rank}(\Lambda)$.
 The $\beta\gamma$ system is \emph{not} self-dual: its Koszul
 dual is the $bc$ system \textup{(}different statistics\textup{)}. The
 polynomial symmetry
 $\kappa_{\beta\gamma}(\lambda)
 = \kappa_{\beta\gamma}(1-\lambda) = 6\lambda^2 - 6\lambda + 1$
 is an internal symmetry of the $\beta\gamma$ formula, not a
 duality statement.

\item \emph{Affine Kac--Moody.}\;
 For $\widehat{\fg}_k$ with
 $\kappa(\widehat{\fg}_k)
 = \dim(\fg)(k + h^\vee)/(2h^\vee)$:
 \begin{equation}\label{eq:km-antisymmetry}
 \kappa(\widehat{\fg}_k)
 + \kappa(\widehat{\fg}_{k'}) \;=\; 0,
 \qquad k' = -k - 2h^\vee.
 \end{equation}
 The Koszul conductor is
 $K = c(k) + c(k') = 2\dim(\fg)$
 \textup{(}level-independent\textup{)}.

\item \emph{Principal $\cW$-algebras.}\;
 For $\cW_N = \cW(\mathfrak{sl}_N, f_{\mathrm{prin}})$
 at central charge~$c$, with
 $\kappa(\cW_N) = \varrho_N \cdot c$ where the anomaly ratio
 $\varrho_N = H_N - 1 = \sum_{j=2}^{N} 1/j$:
 \begin{equation}\label{eq:wn-complementarity-sum}
 \kappa(\cW_N^k) + \kappa(\cW_N^{k'})
 \;=\; \varrho_N \cdot K_N,
 \end{equation}
 where $K_N = c(k) + c(k') = 2(N{-}1)(2N^2{+}2N{+}1)
 = 4N^3 - 2N - 2$ is the \emph{Koszul conductor}
 \textup{(}Remark~\textup{\ref{rem:koszul-conductor-explicit})}.
 The self-dual point is
 $c_* = K_N/2 = (N{-}1)(2N^2{+}2N{+}1)$, where
 $\kappa = \kappa^! = \varrho_N K_N/2$.
\end{enumerate}
\end{proposition}

\begin{proof}
Part~(i): the Heisenberg level-inversion $k \mapsto -k$
(with $h^\vee = 0$) negates $\kappa$. For $\beta\gamma$:
$\kappa(\beta\gamma_\lambda) = 6\lambda^2 - 6\lambda + 1$ and
$\kappa(bc_\lambda) = -(6\lambda^2 - 6\lambda + 1)$ by definition
of the dual system; the sum vanishes identically in~$\lambda$.
The lattice case: $\kappa(V_\Lambda) = \operatorname{rank}(\Lambda)$,
$\kappa(V_\Lambda^!) = -\operatorname{rank}(\Lambda)$
(Theorem~\ref{thm:genus-universality}(ii)).

Part~(ii): by the Feigin--Frenkel involution $k \mapsto k' = -k - 2h^\vee$,
$\kappa(\widehat{\fg}_{k'})
= \dim(\fg)(k' + h^\vee)/(2h^\vee)
= \dim(\fg)(-k - h^\vee)/(2h^\vee)
= -\kappa(\widehat{\fg}_k)$.
The central charge sum
$c(k) + c(k') = 2\dim(\fg)$ is Proposition~\ref{prop:sl2-complementarity-all-genera}.

Part~(iii): by Theorem~\ref{thm:wn-obstruction},
$\kappa(\cW_N) = \varrho_N \cdot c$, so
$\kappa + \kappa' = \varrho_N(c + c') = \varrho_N K_N$.
The Koszul conductor $K_N = c(k) + c(k')$ is level-independent by
Theorem~\ref{thm:central-charge-complementarity}, and equals
$2(N{-}1)(2N^2{+}2N{+}1)$ by direct evaluation of the
Fateev--Lukyanov central charge formula under $k \mapsto -k - 2N$.
The self-dual point satisfies
$\kappa(c_*) = \kappa^!(c_*)$, i.e.,
$\varrho_N c_* = \varrho_N(K_N - c_*)$, giving $c_* = K_N/2$.
\end{proof}

\begin{table}[ht]
\centering
\caption{The complementarity landscape. The column
$\kappa + \kappa^!$ records the complementarity sum
$\kappa(\cA) + \kappa(\cA^!)$
(Proposition~\ref{prop:complementarity-landscape});
$K = c + c'$ is the Koszul conductor; $c_*$ is the self-dual
central charge. All entries are verified against
\texttt{compute/tests/test\_complementarity\_landscape.py}
(79~tests). The Virasoro row is $N = 2$ in the
$\cW_N$ series.}
\label{tab:complementarity-landscape}
\index{complementarity!landscape table}
\renewcommand{\arraystretch}{1.4}
\small
\begin{tabular}{@{}llllll@{}}
\toprule
Family $\cA$ & $\kappa(\cA)$ & $\kappa(\cA^!)$
 & $\kappa + \kappa^!$ & $K = c{+}c'$ & $c_*$ \\[2pt]
\midrule
\multicolumn{6}{@{}l}{\textit{Free fields:
 $\kappa + \kappa^! = 0$}} \\[2pt]
$\cH_k$
 & $k$ & $-k$ & $0$ & n/a & $k = 0$ \\
Free fermion
 & $1/4$ & $-1/4$ & $0$ & n/a & n/a \\
$V_\Lambda$ (rank $r$)
 & $r$ & $-r$ & $0$ & $2r$ & n/a \\
$\beta\gamma_\lambda$
 & $6\lambda^2{-}6\lambda{+}1$
 & $-(6\lambda^2{-}6\lambda{+}1)$
 & $0$ & n/a & n/a \\[4pt]
\midrule
\multicolumn{6}{@{}l}{\textit{Affine Kac--Moody:
 $\kappa + \kappa^! = 0$, $K = 2\dim(\fg)$}} \\[2pt]
$\widehat{\mathfrak{sl}}_2$
 & $3(k{+}2)/4$ & $-3(k{+}2)/4$ & $0$ & $6$ & n/a \\
$\widehat{\mathfrak{sl}}_3$
 & $4(k{+}3)/3$ & $-4(k{+}3)/3$ & $0$ & $16$ & n/a \\
$\widehat{G}_2$
 & $7(k{+}4)/4$ & $-7(k{+}4)/4$ & $0$ & $28$ & n/a \\
$\widehat{E}_8$
 & $62(k{+}30)/15$ & $-62(k{+}30)/15$ & $0$ & $496$ & n/a \\[4pt]
\midrule
\multicolumn{6}{@{}l}{\textit{$\cW_N$-algebras:
 $\kappa + \kappa^! = \varrho_N K_N \neq 0$
}} \\[2pt]
$\mathrm{Vir}_c$ ($= \cW_2$)
 & $c/2$ & $(26{-}c)/2$
 & $13$ & $26$ & $13$ \\
$\cW_3$
 & $5c/6$ & $5(100{-}c)/6$
 & $250/3$ & $100$ & $50$ \\
$\cW_4$
 & $13c/12$ & $13(246{-}c)/12$
 & $533/2$ & $246$ & $123$ \\
$\cW_5$
 & $77c/60$ & $77(488{-}c)/60$
 & $9394/15$ & $488$ & $244$ \\
$\cW_6$
 & $29c/20$ & $29(850{-}c)/20$
 & $2465/2$ & $850$ & $425$ \\
$\cW_7$
 & $223c/140$ & $223(1356{-}c)/140$
 & $75597/35$ & $1356$ & $678$ \\[2pt]
\bottomrule
\end{tabular}
\end{table}

\begin{remark}[Independent derivation via Euler characteristic and index theory]
\label{rem:complementarity-index-theory}
\index{complementarity!index-theoretic derivation}
\index{index theory!complementarity}
There is a second derivation of the scalar complementarity constant
that does not add the explicit $\kappa$-formulas term by term.
Write
\[
K_\kappa(\cA)
\;:=\;
\kappa(\cA) + \kappa(\cA^!).
\]
The bar-cobar package
\textup{(}Theorem~\textup{\ref{thm:chiral-koszul-duality}},
Theorem~\textup{\ref{thm:verdier-bar-cobar}}\textup{)}
places the Koszul pair inside the single ambient complementarity
complex $\mathbf{C}_g(\cA)$ of
Theorem~\textup{\ref{thm:quantum-complementarity-main}}, with Verdier
eigenspace splitting
\[
\mathbf{C}_g(\cA)
\;\simeq\;
\mathbf{Q}_g(\cA) \oplus \mathbf{Q}_g(\cA^!).
\]
On the proved uniform-weight lane, and at genus~$1$ for every standard
family, define the normalized scalar index
\[
\chi_g^{\mathrm{sh}}(\cA)
\;:=\;
\frac{F_g(\cA)}{\lambda_g^{\mathrm{FP}}}
\qquad \textup{(UNIFORM-WEIGHT)},
\]
where $\lambda_g^{\mathrm{FP}}$ is the universal tautological
integral of Theorem~\textup{\ref{thm:family-index}}. The modular
characteristic theorem together with the family-index theorem
\textup{(}Theorems~\textup{\ref{thm:modular-characteristic}}
and~\textup{\ref{thm:family-index}}\textup{)}
identify this normalized Euler characteristic with the scalar shadow
coefficient:
\[
\chi_g^{\mathrm{sh}}(\cA) \;=\; \kappa(\cA).
\]
Thus $\kappa(\cA)$ is the Euler characteristic of the shadow tower
after dividing by the universal Hodge factor.

Passing to the total Verdier-self-dual bar package gives
\[
\chi_g^{\mathrm{tot}}(\cA)
\;:=\;
\chi_g^{\mathrm{sh}}(\cA) + \chi_g^{\mathrm{sh}}(\cA^!)
\;=\;
\frac{F_g(\cA) + F_g(\cA^!)}{\lambda_g^{\mathrm{FP}}}
\;=\;
\kappa(\cA) + \kappa(\cA^!)
\qquad \textup{(UNIFORM-WEIGHT)}
\;=\;
K_\kappa(\cA).
\]
This is the Euler-characteristic route to complementarity:
$K_\kappa(\cA)$ is the normalized Euler characteristic of the total
bar package. It is also the cohomological index of the total
bar-cobar differential on the Koszul pair, because the index of a
differential is by definition the Euler characteristic of its
cohomology. Replacing $\cA$ by $\cA^!$ only exchanges the two Verdier
eigensummands, so the total index depends only on the underlying
$\chirAss$ duality orbit. Over a point the Hodge factor is absent and
the classical commutative Koszul pair has index~$0$; on a curve the
same construction acquires the universal Hodge term and can produce a
nonzero scalar constant.

The three standard checks are immediate. For the Virasoro family the
involution is $c \mapsto 26-c$, whose fixed point is $c = 13$
\textup{(}Remark~\textup{\ref{rem:self-dual-complementarity}}\textup{)}.
Since $\chi_g^{\mathrm{tot}}$ is constant along the duality orbit, it
may be evaluated at the fixed point:
\[
K_\kappa(\mathrm{Vir})
\;=\;
\chi_g^{\mathrm{tot}}(\mathrm{Vir}_{13})
\;=\;
2\,\chi_g^{\mathrm{sh}}(\mathrm{Vir}_{13})
\;=\;
2\,\kappa(\mathrm{Vir}_{13})
\;=\;
2 \cdot \frac{13}{2}
\;=\;
13.
\]
This recovers the Virasoro scalar complementarity constant from the
self-dual fixed point plus index theory, independently of the direct
identity $c/2 + (26-c)/2 = 13$.

For affine Kac--Moody, the Feigin--Frenkel involution
$k' = -k - 2h^\vee$ gives
\[
\kappa(V_k(\fg)) + \kappa(V_{k'}(\fg)) \;=\; 0,
\]
so the total bar-cobar index vanishes on the scalar lane:
$K_\kappa(\widehat{\fg}) = 0$.

For Bershadsky--Polyakov, Proposition~\ref{prop:bp-kappa} gives
$\kappa(\mathcal{B}^k) = c(k)/6$, while
Proposition~\ref{prop:bp-self-duality} gives
$c(k) + c(-k-6) = 196$. Therefore
\[
K_\kappa(\mathcal{B})
\;=\;
\kappa(\mathcal{B}^k) + \kappa(\mathcal{B}^{-k-6})
\;=\;
\frac{c(k) + c(-k-6)}{6}
\;=\;
\frac{196}{6}
\;=\;
\frac{98}{3}.
\]
Thus the same bar-cobar package carries two related invariants: the
unnormalized central-charge conductor $c + c'$ and the scalar
complementarity constant $K_\kappa = \kappa + \kappa'$. For
Bershadsky--Polyakov the former is $196$ and the latter is $98/3$.
\end{remark}

\begin{remark}[The Koszul conductor polynomial]
\label{rem:koszul-conductor-polynomial}
\index{Koszul conductor!polynomial identity}
The identity $K_N = 2(N{-}1)(2N^2{+}2N{+}1) = 4N^3 - 2N - 2$
shows that the Koszul conductor is a cubic polynomial in~$N$.
Its growth is controlled by the leading term $4N^3$, so
$K_N/\dim(\mathfrak{sl}_N) = K_N/(N^2-1) \to 4N$ as
$N \to \infty$. In contrast, the affine Kac--Moody Koszul
conductor $K^{\mathrm{aff}} = 2\dim(\fg) = 2(N^2 - 1)$ grows
quadratically. The excess $K_N - 2(N^2 - 1) = 4N^3 - 2N^2 - 2N
= 2N(2N^2 - N - 1)$ is the contribution of the DS ghost sector.
\end{remark}

\begin{remark}[Self-dual points and the complementarity sum]
\label{rem:self-dual-complementarity}
\index{complementarity!self-dual point}
\index{Virasoro algebra!chiral Koszul self-dual point}
At the chiral Koszul self-dual point $c = c_* = K_N/2$, the scalar
identity gives
$\kappa(\cW_N) = \kappa(\cW_N^!) = \varrho_N K_N/2$.
If a balanced anti-equivariant self-duality representative is
constructed on the genus-$g$ ambient complex, then the
complementarity splitting
$Q_g(\cA) \oplus Q_g(\cA^!)$ is a balanced decomposition
with $\dim Q_g(\cA) = \dim Q_g(\cA^!)$
(Theorem~\ref{thm:self-dual-halving}).
The Virasoro chiral Koszul self-dual point is $c = 13$ (not $c = 26$;
see Corollary~\ref{cor:virasoro-quantum-dim}).
The critical value $c = 26$, where $\kappa(\mathrm{Vir}_{26}^!)
= 0$, is the point where the \emph{dual} bar complex becomes
uncurved (Theorem~\ref{thm:critical-dimension}). These two
distinguished values $13$ and $26$ correspond to different
phenomena: $c = 13$ is self-duality of the Koszul pair, $c = 26$
is vanishing of the dual curvature.
\end{remark}

\begin{remark}[Level independence]
\label{rem:level-independence-landscape}
\index{complementarity!level independence}
In every row of Table~\ref{tab:complementarity-landscape}, the
complementarity sum $\kappa + \kappa^!$ is independent of
the level parameter~$k$. For affine Kac--Moody this is
immediate from the Feigin--Frenkel anti-symmetry
$\kappa(k) = -\kappa(k')$. For $\cW_N$-algebras it follows
because $\kappa = \varrho_N \cdot c(k)$ and
$c(k) + c(k') = K_N$ is level-independent
(Theorem~\ref{thm:central-charge-complementarity}). For free
fields the sum vanishes identically, with no level to vary.
The level independence is a topological invariant of the Koszul
pair: it depends only on the root datum (for
Lie-theoretic families) or the conformal-weight polynomial
(for $\beta\gamma/bc$).
\end{remark}

\section{Higher genus extension: descent and acyclicity}
\label{sec:higher-genus-descent}

\subsection{Beilinson--Drinfeld foundations: genus zero review}
\label{subsec:BD-genus-zero-review}

We review the Beilinson--Drinfeld construction at genus zero, identifying what generalizes.

\begin{theorem}[BD 3.4.12: genus zero acyclicity; \ClaimStatusProvedHere]\label{thm:BD-genus-zero}
\index{Beilinson--Drinfeld!higher genus}
For a smooth projective curve $X$ and chiral algebra $\mathcal{A}$, the Chevalley--Cousin complex $C(\mathcal{A})$ defined over the Ran space $R(X)$ is acyclic:
\[H^i(R(X), C(\mathcal{A})) =
\begin{cases}
\mathcal{A} & i = 0 \\
0 & i \neq 0
\end{cases}\]
\end{theorem}

\begin{proof}[Key steps from BD §3.4]
\emph{Step 1 (BD 3.4.10):} Embed $M(X) \hookrightarrow M(X^S)$ using the diagonal embedding $\Delta^{(S)}_*$. This is fully faithful and pseudo-tensor.

\emph{Step 2 (BD 3.4.11):} Construct the Chevalley--Cousin complex:
\[C(\mathcal{A})_{X^I} = \bigoplus_{T \in Q(I)} \Delta^{(I/T)}_* j^{(I/T)}_* j^{(I/T)*} \omega_{X^T}[|T|]\]
where $j^{(I/T)}: X^T \to X^I$ removes the diagonals.

\emph{Step 3 (BD 3.4.12):} Prove acyclicity via Cousin filtration. The key ingredients:
\begin{enumerate}
\item \emph{Descent compatibility:} The natural map $C(\mathcal{A}) \to \Delta^{(S)}_* \mathcal{A}$ is a quasi-isomorphism
\item \emph{Stratification:} Boundary divisors have normal crossings (Fulton--MacPherson)
\item \emph{Residue calculus:} The differential computes via iterated residues at collision divisors
\end{enumerate}
\end{proof}

\begin{remark}[Requirements for higher genus]\label{rem:what-to-preserve}
To extend to higher genus, we must ensure:
\begin{enumerate}
\item The factorization structure persists: $\mathcal{A}(U \sqcup V) \simeq \mathcal{A}(U) \otimes \mathcal{A}(V)$
\item Normal crossings are maintained in the boundary divisors
\item Descent data are compatible with moduli stack stratification
\item Quantum corrections (from $H^1(\Sigma_g, \C)$, the period data of the fiber curve) preserve acyclicity
\end{enumerate}
\end{remark}

\subsection{The universal curve and relative Ran space}
\label{subsec:universal-curve-ran}

At higher genus, we work over the moduli stack $\mathcal{M}_g$.

\begin{construction}[Relative Ran space]\label{const:relative-ran}
Let $\pi: \mathcal{C}_g \to \mathcal{M}_g$ be the universal curve of genus $g$. The \emph{relative Ran space} is:
\[R(\mathcal{C}_g/\mathcal{M}_g) := \colim_{n \geq 0} (\mathcal{C}_g)^{(n)}/\mathcal{M}_g\]
where $(\mathcal{C}_g)^{(n)} = \mathcal{C}_g^n \setminus \{\text{diagonals}\}$ is the configuration space of $n$ distinct points.

\emph{Fiber over a point.} For $[\Sigma_g] \in \mathcal{M}_g$, the fiber is:
\[R(\mathcal{C}_g/\mathcal{M}_g)|_{[\Sigma_g]} = R(\Sigma_g)\]
the ordinary Ran space of the Riemann surface $\Sigma_g$.
\end{construction}

\begin{proposition}[Factorization over moduli; \ClaimStatusProvedHere]\label{prop:factorization-over-moduli}
For disjoint open sets $U, V \subset \Sigma_g$ varying in families over $\mathcal{M}_g$:
\[\mathcal{A}(U \sqcup V) \simeq \mathcal{A}(U) \otimes_{\mathcal{O}_{\mathcal{M}_g}} \mathcal{A}(V)\]
The factorization is $\mathcal{O}_{\mathcal{M}_g}$-linear.
\end{proposition}

\begin{proof}
The chiral product $\mu\colon j_*j^*(\mathcal{A} \boxtimes \mathcal{A}) \to \Delta_*\mathcal{A}$ is a morphism of $\mathcal{D}$-modules defined locally on $X$ (Beilinson--Drinfeld~\cite{BD04}, \S3.4). For disjoint open sets $U, V \subset \Sigma_g$, the factorization isomorphism $\mathcal{A}(U \sqcup V) \simeq \mathcal{A}(U) \otimes \mathcal{A}(V)$ depends only on the restriction of $\mathcal{A}$ to $U \sqcup V$ and the factorization isomorphism of the $\mathcal{D}$-module structure. Since $U \cap V = \emptyset$, no global data (periods, modular parameter) enters: the isomorphism is determined by the local chiral product and the $\mathcal{D}$-module structure on each open set. The base $\mathcal{O}_{\mathcal{M}_g}$-linearity follows because the family $\pi\colon \mathcal{C}_{g} \to \mathcal{M}_g$ is a smooth family of curves, and the relative $\mathcal{D}$-module structure on $\mathcal{A}$ over $\mathcal{M}_g$ ensures all structure maps are $\mathcal{O}_{\mathcal{M}_g}$-linear.
\end{proof}

\subsection{Normal crossings: Deligne--Mumford + Fulton--MacPherson}
\label{subsec:normal-crossings-combined}

\begin{theorem}[Normal crossings persist at higher genus; \ClaimStatusProvedHere]\label{thm:normal-crossings-persist}
\index{normal crossing!boundary}
The fiber product:
\[\mathcal{Z}_{g,n} := \overline{\mathcal{M}}_{g,n} \times_{X^n} \overline{C}_n(X)\]
has boundary divisors in normal crossings.
\end{theorem}

\begin{proof}
\emph{Step 1: Deligne--Mumford normal crossings.}

By Deligne--Mumford~\cite{DeligneM69}, $\overline{\mathcal{M}}_{g,n}$ is a smooth Deligne--Mumford stack with boundary:
\[\partial \overline{\mathcal{M}}_{g,n} = \bigcup D_{g_1, g_2, S}\]
parametrizing stable curves with nodes. Each boundary divisor has equation $q = 0$ locally, where $q$ is the nodal parameter.

\emph{Step 2: Fulton--MacPherson normal crossings.}

The configuration space $\overline{C}_n(X)$ has boundary:
\[\partial \overline{C}_n(X) = \bigcup D_{I|J}\]
where $I \sqcup J = [n]$ parametrizes collisions. Each divisor has local equation $\epsilon_{ij} = |z_i - z_j| = 0$ (after blow-up).

\emph{Step 3: Fiber product preservation.}

The maps $\overline{\mathcal{M}}_{g,n} \to X^n$ and $\overline{C}_n(X) \to X^n$ are both:
\begin{itemize}
\item Proper
\item With normal crossing boundaries
\item Transverse to each other
\end{itemize}

By the Knudsen--Mumford construction~\cite{Knudsen83}, the fiber product of normal crossing divisors is normal crossing.

\emph{Step 4: Local coordinates near boundary.}

Near a point where both boundaries intersect, we have local coordinates:
\begin{align*}
&(q, \tau_1, \ldots, \tau_{3g-3+n}) && \text{for } \overline{\mathcal{M}}_{g,n} \\
&(\epsilon_{ij}, w_1, \ldots, w_{2n-2}) && \text{for } \overline{C}_n(X)
\end{align*}

The boundary has equation $q \cdot \epsilon_{ij} = 0$, which is normal crossing. \qedhere
\end{proof}

\begin{theorem}[Chevalley--Cousin acyclicity at higher genus; \ClaimStatusConditional]\label{thm:CC-acyclicity-higher-genus}
\label{thm:chevalley-cousin}
\label{thm:BD-extension-higher-genus}
Let $X$ be a smooth projective curve. The Chevalley--Cousin complex $C(\mathcal{A})$
defined over the moduli stack $\mathcal{M}_{g,n}$ remains acyclic, extending
Beilinson--Drinfeld's genus-zero result (BD \cite[Theorem 3.4.12]{BD04}).

\emph{Statement.} For each genus $g \geq 0$ and $n \geq 1$, the natural map
\[
\begin{aligned}
R\Gamma(R(X), C(\mathcal{A}))
&\to R\Gamma\bigl(\mathcal{M}_{g,n} \times_{X^n} C_n(X),\, C(\mathcal{A})\bigr)
\end{aligned}
\]
is a quasi-isomorphism. Here $C(\mathcal{A})$ is equipped with quantum
corrections parameterized by the period data
\[
t_g \in H^0\!\bigl(\mathcal{M}_g,\, R^1\pi_* Z(\mathcal{A})\bigr),
\]
where $\pi \colon \mathcal{C}_g \to \mathcal{M}_g$ is the universal curve and $Z(\mathcal{A})$ is the center. (The parameters arise from $H^1(\Sigma_g, \C)$, not from $H^1(\mathcal{M}_g)$, which vanishes for $g \geq 2$ by Harer's theorem; see Convention~\ref{conv:higher-genus-differentials}.)
\end{theorem}

\begin{proof}

We extend the BD genus-zero proof systematically and handle each higher-genus phenomenon in turn.

\subsection*{Higher-genus relative geometry}

\begin{center}
\begin{tabular}{|l|p{5cm}|p{5cm}|}
\hline
\textbf{Structure} & \textbf{Genus 0 (BD)} & \textbf{Genus $g \geq 1$} \\
\hline
Base space & $X$ (curve) & $\mathcal{C}_g \to \mathcal{M}_g$ (universal curve) \\
\hline
Ran space & $R(X) = \colim_n C_n(X)$ & $R(\mathcal{C}_g/\mathcal{M}_g)$ \\
\hline
Moduli & $\text{pt}$ (no moduli) & $\mathcal{M}_g$ (moduli stack, dim $3g-3$) \\
\hline
Quantum corrections & None & From $H^1(\Sigma_g) \neq 0$ (period integrals, $g \geq 1$) \\
\hline
Differential forms & Logarithmic, rational & Logarithmic + elliptic/abelian \\
\hline
Boundary & Normal crossings (FM) & Normal crossings (FM + DM) \\
\hline
\end{tabular}
\end{center}

\subsection*{Relative descent compatibility}

\emph{Step 1: Descent data along $R(X) \to X$.}

Following BD \cite[§3.4.10-3.4.11]{BD04}, embed $M(X)$ into the larger category
$M(X^S)$ with tensor structures $\otimes_*$ and $\otimes_{ch}$. The key is that
$\Delta^{(S)}_*: M(X)_{ch} \hookrightarrow M(X^S)_{ch}$ is a fully faithful
pseudo-tensor embedding.

Does this embedding preserve good properties when we replace $X$ with $\mathcal{C}_g \to \mathcal{M}_g$?

\begin{lemma}[Relative diagonal embedding; \ClaimStatusProvedHere]\label{lem:relative-diagonal}
The relative diagonal embedding:
\[\Delta^{(S)}_{/\mathcal{M}_g}: M(\mathcal{C}_g/\mathcal{M}_g) \hookrightarrow M((\mathcal{C}_g)^S/\mathcal{M}_g)\]
is fully faithful and pseudo-tensor, fiberwise over $\mathcal{M}_g$.
\end{lemma}

\begin{proof}
The embedding is defined fiberwise: for each $[\Sigma_g] \in \mathcal{M}_g$, we have:
\[\Delta^{(S)}|_{[\Sigma_g]}: M(\Sigma_g) \hookrightarrow M(\Sigma_g^S)\]
which is the BD embedding for the specific curve $\Sigma_g$.

\emph{Full faithfulness.} For $\mathcal{F}, \mathcal{G} \in M(\mathcal{C}_g/\mathcal{M}_g)$:
\begin{align*}
\Hom(\Delta^{(S)}_* \mathcal{F}, \Delta^{(S)}_* \mathcal{G})
&= \int_{[\Sigma_g] \in \mathcal{M}_g} \Hom_{\Sigma_g}(\Delta^{(S)}_* \mathcal{F}|_{[\Sigma_g]}, \Delta^{(S)}_* \mathcal{G}|_{[\Sigma_g]}) \\
&\simeq \int_{[\Sigma_g] \in \mathcal{M}_g} \Hom_{\Sigma_g}(\mathcal{F}|_{[\Sigma_g]}, \mathcal{G}|_{[\Sigma_g]}) && \text{(BD, fiberwise)} \\
&= \Hom(\mathcal{F}, \mathcal{G})
\end{align*}
The second isomorphism uses BD's full faithfulness on each fiber.

\emph{Pseudo-tensor.} The tensor structure $\otimes_{ch}$ on $M(\mathcal{C}_g/\mathcal{M}_g)$ is defined fiberwise, so preservation follows from BD fiberwise. \qedhere
\end{proof}

\subsection*{Stable-stratum compatibility}

\emph{Step 2: Compatibility with stratification by stable curves.}

\begin{definition}[Boundary strata]\label{def:boundary-strata}
The boundary of $\overline{\mathcal{M}}_{g,n}$ has components:
\begin{enumerate}
\item \emph{Separating nodes:} $D_{g_1,g_2,S}$ where $g_1 + g_2 = g$, $S \sqcup T = [n]$
\[D_{g_1,g_2,S} \simeq \overline{\mathcal{M}}_{g_1, |S|+1} \times \overline{\mathcal{M}}_{g_2, |T|+1}\]
Parametrizes curves that split into two components of genera $g_1, g_2$.

\item \emph{Non-separating nodes:} $D_{\text{irr}}$
\[D_{\text{irr}} \simeq \overline{\mathcal{M}}_{g-1, n+2}\]
Parametrizes curves with a self-node (attaching handle).
\end{enumerate}
\end{definition}

\begin{proposition}[Gluing formula at nodes; \ClaimStatusProvedHere]\label{prop:gluing-at-nodes}
For a stable curve $C$ with a node $p$ splitting it into $C_1 \cup_p C_2$:
\[\mathcal{A}(C) \simeq \mathcal{A}(C_1) \otimes_{\mathcal{A}(p)} \mathcal{A}(C_2)\]
where the tensor product is over the fiber algebra $\mathcal{A}(p)$ at the node.
\end{proposition}

\begin{proof}
\emph{Step 1: Formal neighborhood of node.}

Near a node $p$, we have local analytic coordinates $(u,v)$ with $uv = t$ where $t \to 0$ as we approach the boundary. The two branches are:
\begin{align*}
C_1^{\text{loc}} &= \{(u, v) : v = 0, u \neq 0\} \cup \{p\} \\
C_2^{\text{loc}} &= \{(u, v) : u = 0, v \neq 0\} \cup \{p\}
\end{align*}

\emph{Step 2: Chiral algebra factorizes.}

For disjoint opens $U_1 \subset C_1$, $U_2 \subset C_2$ with $U_1 \cap U_2 = \emptyset$:
\[\mathcal{A}(U_1 \sqcup U_2) = \mathcal{A}(U_1) \otimes \mathcal{A}(U_2)\]
by the factorization axiom.

\emph{Step 3: Taking limits.}

As $t \to 0$ (node formation), the two branches $C_1^{\text{loc}}$ and $C_2^{\text{loc}}$ come together at $p$. The factorization persists:
\[\mathcal{A}(C_1^{\text{loc}} \cup C_2^{\text{loc}}) = \lim_{t \to 0} \mathcal{A}(U_1(t) \sqcup U_2(t)) = \mathcal{A}(C_1) \otimes_{\mathcal{A}(p)} \mathcal{A}(C_2)\]
where the tensor product over $\mathcal{A}(p)$ accounts for the gluing. The limit passage is justified by the factorization property of chiral homology \cite[Theorem~3.4.9]{BD04}: for a family of curves $\mathcal{C} \to S$ degenerating to a nodal curve $C_0 = C_1 \cup_p C_2$, the chiral homology satisfies $H^{\mathrm{ch}}(C_0, \mathcal{A}) \simeq H^{\mathrm{ch}}(C_1, \mathcal{A}) \otimes_{H^{\mathrm{ch}}(\{p\}, \mathcal{A})} H^{\mathrm{ch}}(C_2, \mathcal{A})$ (the tensor product is over the stalk, which is $\mathcal{A}(p)$ for the unitalization). This is the bar complex incarnation of the gluing axiom for chiral homology. \qedhere
\end{proof}

\begin{lemma}[Boundary compatibility; \ClaimStatusProvedHere]\label{lem:boundary-compatible}
The restriction of $C(\mathcal{A})$ to each boundary stratum
$\mathcal{M}_{g_1,n_1+1} \times \mathcal{M}_{g_2,n_2+1}$ (with $g_1+g_2=g$, $n_1+n_2=n$)
is computed by the gluing formula:
\[C(\mathcal{A})|_{\text{boundary}} \simeq C(\mathcal{A})|_{\mathcal{M}_{g_1,n_1+1}}
\otimes_{A(p)} C(\mathcal{A})|_{\mathcal{M}_{g_2,n_2+1}}\]
where the tensor product is over the fiber $A(p)$ at the nodal point.
\end{lemma}

\begin{proof}[Proof of Lemma]
At a node $p$ in a stable curve, we have local coordinate patches $(U_1, z_1)$ and
$(U_2, z_2)$ with $z_1 \cdot z_2 = t$ where $t \to 0$ as we approach the boundary.

The factorization property of chiral algebras gives:
\[\mathcal{A}(U_1 \sqcup U_2) \simeq \mathcal{A}(U_1) \otimes_{A(p)} \mathcal{A}(U_2)\]

The Chevalley--Cousin complex respects this factorization:
\[C(\mathcal{A}(U_1 \sqcup U_2)) \simeq C(\mathcal{A}(U_1)) \otimes_{A(p)} C(\mathcal{A}(U_2))\]

As $t \to 0$ (approaching boundary), this tensor product structure persists in the limit.
\end{proof}

\begin{corollary}[Chevalley--Cousin at boundary; \ClaimStatusProvedHere]\label{cor:CC-at-boundary}
The Chevalley--Cousin complex respects boundary stratification:
\[C(\mathcal{A})|_{\partial \overline{\mathcal{M}}_{g,n}} = \bigoplus_{\text{strata } D} C(\mathcal{A})|_D\]
where each $C(\mathcal{A})|_D$ is computed via the gluing formula.
\end{corollary}
\begin{proof}
Each boundary stratum $D \subset \partial\overline{\mathcal{M}}_{g,n}$ corresponds to a stable curve with nodes. By the preceding theorem, the Chevalley--Cousin complex is compatible with the tensor product decomposition at each node. The direct sum decomposition follows from the normal crossings property of $\partial\overline{\mathcal{M}}_{g,n}$ (Theorem~\ref{thm:normal-crossings-persist}).
\end{proof}

\subsection*{Quantum-corrected differential}

\emph{Step 3: Quantum corrections and modular parameters.}

The genus-$g$ bar complex receives quantum corrections parametrized by the
period data of the genus-$g$ curve. For $g=1$: the single modular parameter $\tau \in \mathfrak{h}$
(upper half-plane). For $g \geq 2$: the $g$ $A$-periods $\oint_{A_i}\omega_j$ provide
$g$ independent parameters.

(Note: $H^1(\mathcal{M}_g, \mathbb{Q}) = 0$ for $g \geq 2$ by Harer's theorem. The
parameters $t_i$ below arise from $H^1(\Sigma_g, \mathbb{C})$, the cohomology of the
\emph{curve}, not of the \emph{moduli space}.)

These enter the total corrected differential
(Convention~\ref{conv:higher-genus-differentials}) as:
\[\Dg{g} = \dzero + \sum_{i=1}^g t_i \cdot d_i\]
where $t_i$ are the period parameters arising from $H^1(\Sigma_g)$ and $d_i$ are the
genus-$g$ correction terms coming from period integrals.

\begin{definition}[Quantum-corrected differential]\label{def:quantum-differential}
At genus $g$, the total corrected differential on the Chevalley--Cousin
complex (Convention~\ref{conv:higher-genus-differentials}) receives corrections:
\[\Dg{g} = \dzero + \sum_{k=1}^{g} t_k \cdot d_k\]
where:
\begin{itemize}
\item $\dzero$ is the genus-zero (classical) differential from BD
\item $t_k = \oint_{A_k} \omega_k$ are the $A$-period parameters arising from the Hodge decomposition of $H^1(\Sigma_g, \mathbb{C})$
\item $d_k$ are correction operators encoding quantum effects via integration over the $A$-cycle $\gamma_k$
\end{itemize}

\emph{Explicit form.} For genus $g=1$ (elliptic case), there is one $A$-period $t_1 = \tau$, and:
\[\Dg{1} = \dzero + \tau \cdot d_{\text{elliptic}}\]
where $\tau$ is the modulus of the torus and $d_{\text{elliptic}}$ involves elliptic functions.
\end{definition}

\begin{theorem}[Key property: \texorpdfstring{$\Dg{g}^{\,2} = 0$}{D(g) squared = 0}; \ClaimStatusConditional]\label{thm:quantum-diff-squares-zero}
\textup{[Regime: curved-central
\textup{(}i.e., curvature $\kappa$ is a central scalar;
the theorem proves that $\Dg{g}$ is \emph{flat}\/\textup{)};
all genera
\textup{(}Convention~\textup{\ref{conv:regime-tags})}.]}

The total corrected differential $\Dg{g}$ \textup{(}Convention~\textup{\ref{conv:higher-genus-differentials}(ii))} satisfies $\Dg{g}^{\,2} = 0$.
\end{theorem}

\begin{proof}
\emph{Step 1: Expansion of $\Dg{g}^{\,2}$.}

\begin{align*}
\Dg{g}^{\,2} &= \left(\dzero + \sum_k t_k d_k\right)^2 \\
&= \dzero^2 + \sum_k t_k (\dzero d_k + d_k \dzero) + \sum_{k,l} t_k t_l d_k d_l
\end{align*}

\emph{Step 2: Classical term vanishes.}

$\dzero^2 = 0$ by the genus-zero cobar nilpotency (Theorem~\ref{thm:cobar-d-squared-zero}).

\emph{Step 3: Mixed terms vanish.}

The correction operators $d_k$ are constructed from period integrals over the $A$-cycles $\gamma_k \subset \Sigma_g$. We claim:
\[\dzero d_k + d_k \dzero = 0\]

The residue operator $\dzero = \sum_D \operatorname{Res}_D$ computes residues at collision divisors $D_{ij} = \{z_i = z_j\} \subset \overline{C}_n(\Sigma_g)$, while $d_k = \oint_{\gamma_k}$ integrates a single variable $z_i$ over the cycle $\gamma_k$ on $\Sigma_g$. Both compositions vanish independently:
\[
\dzero \circ d_k = 0 \quad \text{and} \quad d_k \circ \dzero = 0.
\]
For $\dzero \circ d_k = 0$: after integrating $z_i$ over $\gamma_k$, the result is a form on a configuration space with one fewer variable. The collision singularity $(z_i - z_j)^{-1}$ is replaced by $\oint_{\gamma_k} (z_i - z_j)^{-1} dz_i$, which is either $0$ (if $z_j \notin \gamma_k$) or $2\pi i$ (if $z_j \in \gamma_k$), in both cases producing a regular function of $z_j$ with no residue to extract. For $d_k \circ \dzero = 0$: extracting the residue $\operatorname{Res}_{D_{ij}}$ at $z_i = z_j$ eliminates the variable $z_i$ (replacing it by the OPE coefficient at $z_j$); the subsequent period integral $\oint_{\gamma_k} dz_i$ has nothing to act on, since $z_i$ no longer appears.

\emph{Pole-order separation.} The full OPE contains higher-order poles $(z_i - z_j)^{-m}$ for $m \geq 2$, but these contribute to the correction operators $d_k$ with $k \geq 2$ in the expansion $\Dg{g} = \dzero + \sum_k t_k d_k$, not to $\dzero$ or $d_1$. The relation $\dzero d_k + d_k \dzero = 0$ established above involves only the residue operator $\dzero$ and the period integral $d_k$; the simple-pole analysis suffices because $\dzero$ extracts first-order residues by definition (Definition~\ref{def:bar-differential-complete}).

\emph{Step 4: Quantum terms vanish; the Fay trisecant identity
is the higher-genus Borcherds identity.}

Three propagators control the bar differential at genus~$g$,
producing three chain-level models with different algebraic types
(Convention~\ref{conv:higher-genus-differentials}):
\begin{center}
\small
\begin{tabular}{lllll}
\textbf{Model} & \textbf{Propagator} & \textbf{Arnold 3-form} & \textbf{Differential} & \textbf{Category} \\[2pt]
Classical &
 $(z_i - z_j)^{-1}\,dz_i$ &
 $\mathcal{A}_3^{(0)} = 0$ &
 $\dzero^2 = 0$ &
 derived \\[2pt]
Corrected hol.\ &
 $\partial_{z_i}\!\log E(z_i,z_j)$ &
 $\mathcal{A}_3^{\mathrm{hol}} = 0$ &
 $\Dg{g}^{\,2} = 0$ &
 derived \\[2pt]
Curved geom.\ &
 $\eta_{ij}^{(g)}$ (Arakelov) &
 $\mathcal{A}_3^{(g)} = 2\pi i\,\omega_{\mathrm{Ar}}$ &
 $\dfib^{\,2} = \kappa\cdot\omega_g$ &
 coderived
\end{tabular}
\end{center}
The first line is the classical Arnold relation
(Theorem~\ref{thm:arnold-three}).
The third line is the quantum-corrected Arnold relation
(Theorem~\ref{thm:quantum-arnold-relations}):
the non-holomorphic correction in the Arakelov propagator
breaks the Arnold relation by exactly
$2\pi i\,\omega_{\mathrm{Ar}}$, producing the curvature
$\dfib^{\,2} = \kappa(\cA)\cdot\omega_g$.
The middle line is the key:

\emph{The holomorphic propagator satisfies the exact Arnold
relation at every genus.} This was established
genus-by-genus in the proof of
Theorem~\ref{thm:quantum-arnold-relations}:
at genus~$1$, the holomorphic cyclic sum
$\zeta(z_{12})\zeta(z_{23}) + \mathrm{cyc}
= \frac{1}{2}[\wp(z_{12}) + \wp(z_{23}) + \wp(z_{31})]$
is symmetric, hence killed by wedge with antisymmetric
form factors; at genus~$\geq 2$, the same vanishing
follows from the Fay trisecant identity for the Siegel
$\wp$-matrix (equation~\eqref{eq:fay-g2} and its
generalization).

The total corrected differential $\Dg{g}$ is the bar
differential built from the holomorphic propagator
$h_{ij} = \partial_{z_i}\!\log E(z_i, z_j)$
(Definition~\ref{def:higher-genus-log-forms}).
Since $\mathcal{A}_3^{\mathrm{hol}} = 0$, the nine-term
verification of $\Dg{g}^{\,2} = 0$ proceeds exactly as at
genus~$0$ (Theorem~\ref{thm:cobar-d-squared-zero}),
with the Fay identity replacing the classical Arnold
relation at each three-point interaction.
The Fay trisecant identity is the higher-genus
Borcherds identity.

\emph{Component-level consequence.}
From $\Dg{g}^{\,2} = 0$ together with the
independently-proved $\dzero^2 = 0$ (Step~2) and
$\dzero\, d_k + d_k\,\dzero = 0$ (Step~3):
\[
 \sum_{k,l} t_k\,t_l\, d_k d_l
 = \Dg{g}^{\,2} - \dzero^2
 - \sum_k t_k\bigl(\dzero\,d_k + d_k\,\dzero\bigr)
 = 0.
\]
The parameters $t_k$ vary independently as
$[\Sigma_g]$ moves in the Siegel upper half-space
$\mathfrak{H}_g$, so the vanishing holds coefficient-by-coefficient:
\[
 d_k\,d_l + d_l\,d_k = 0
 \quad\text{for all } 1 \le k,l \le g.
\]
This is the operator-level manifestation of the Lagrangian
isotropicity $\gamma_k \cdot \gamma_l = 0$
of the $A$-cycle subspace in $H_1(\Sigma_g, \mathbb{Z})$:
the Riemann bilinear relations connect iterated period
integrals of the Cauchy kernel to the symplectic
intersection pairing, and restricting to a Lagrangian
subspace forces all anticommutators to vanish.

\emph{(Note: The two vanishings $\Dg{g}^{\,2} = 0$ and
$\dfib^{\,2} \neq 0$ are not in tension; they are the two
faces of the chiral Riemann--Hilbert correspondence
(Remark~\textup{\ref{rem:curvature-riemann-hilbert}}).
The holomorphic propagator defines a flat connection on
the universal cover of $\Sigma_g$. The Arakelov
propagator defines the single-valued but curved
projection to $\Sigma_g$ itself. The total corrected
differential $\Dg{g}$ lives in the flat gauge; the
fiberwise differential $\dfib$ lives in the curved gauge.
The passage $\dfib \leadsto \Dg{g}$ is the chiral
exponential map.)} \qedhere
\end{proof}

\begin{remark}[Curved twisting morphisms at higher genus]
\label{rem:curved-twisting-higher-genus}
\index{twisting morphism!curved|textbf}
\index{Maurer--Cartan equation!curved}
At genus~$0$, the universal twisting morphism
$\tau\colon \bar{B}(\cA) \to \cA$ satisfies the flat
Maurer--Cartan equation $d\tau + \tau \star \tau = 0$
(Corollary~\ref{cor:three-bijections},
Remark~\ref{rem:MC-is-Stokes}).
At genus~$g \geq 1$, the propagator acquires monodromy and the
twisting morphism becomes \emph{curved}:
\begin{equation}\label{eq:curved-MC-higher-genus}
d\tau_g + \tau_g \star \tau_g = \kappa(\cA) \cdot \omega_g.
\end{equation}
The curvature $\kappa(\cA) \cdot \omega_g$ is the Arakelov
$(1,1)$-form weighted by the modular characteristic
(Theorem~\ref{thm:quantum-arnold-relations}): it measures
the failure of the single-valued propagator to satisfy the
Arnold relation at genus~$g$.

The three propagators in the table of Step~4 above correspond
to three twisting morphisms: $\tau_0$ (classical, flat MC),
$\tau_g^{\mathrm{hol}}$ (holomorphic, flat MC via the Fay
identity), and $\tau_g^{\mathrm{Ar}}$ (Arakelov, curved MC
with curvature~$\kappa \cdot \omega_g$).
The total corrected differential $\Dg{g}$ corresponds to
$\tau_g^{\mathrm{hol}}$; the fiberwise differential $\dfib$
corresponds to $\tau_g^{\mathrm{Ar}}$.
The passage between them is the chiral exponential map.

The three incarnations descend through the H/M/S hierarchy:
\begin{itemize}
\item \emph{H-level.} The shadow obstruction tower
 $\Theta_\cA^{\leq r}$
 (Definition~\ref{def:shadow-postnikov-tower})
 assembles the modular data at each finite order;
 its all-degree limit
 $\Theta_\cA \in
 \mathrm{MC}(\mathfrak{g}_\cA^{\mathrm{mod}})$
 gives a single modular twisting cochain
 (Theorem~\ref{thm:mc2-bar-intrinsic};
 Remark~\ref{rem:theta-modular-twisting}).
\item \emph{M-level.} At each genus, the three twisting
 morphisms in the table above are explicit dg maps,
 with the curved MC equation~\eqref{eq:curved-MC-higher-genus}
 encoding the Arakelov defect.
\item \emph{S-level.} The scalar shadow is
 $\kappa(\cA) \cdot \lambda_g \in H^*(\overline{\mathcal{M}}_g)$
 \textup{(}UNIFORM-WEIGHT\textup{)}
 (Theorem~\ref{thm:genus-universality}):
 one number per genus,
 determined by the genus-$1$ curvature alone.
 For multi-weight algebras at $g \geq 2$, the scalar formula
 receives a cross-channel correction $\delta F_g^{\mathrm{cross}}$
 (Theorem~\ref{thm:multi-weight-genus-expansion}).
\end{itemize}
\end{remark}

\begin{lemma}[Quantum corrections preserve acyclicity; \ClaimStatusConditional]\label{lem:quantum-preserves-acyclicity}
The total corrected differential $\Dg{g}$ satisfies $\Dg{g}^{\,2} = 0$ and preserves the
acyclicity of the Chevalley--Cousin complex.
\end{lemma}

\begin{proof}
We have shown $\Dg{g}^{\,2} = 0$ in Theorem~\ref{thm:quantum-diff-squares-zero}. It remains to show acyclicity.

\emph{Step 1: Spectral sequence with quantum corrections.}

Filter $C(\mathcal{A})$ by the Cousin filtration $F_p$. The $E_1$ page is:
\[E_1^{p,q} = H^{p+q}(\mathcal{M}_g \times C_p(X), \text{gr}_p C(\mathcal{A}))\]

Quantum corrections affect only the differentials $d_r: E_r^{p,q} \to E_r^{p+r, q-r+1}$ for $r \geq 1$.

\emph{Step 2: Classical acyclicity implies quantum acyclicity.}

The quantum corrections $t_k d_k$ preserve the Cousin filtration and act as derivations with respect to the chiral product. Write $\Dg{g} = \dzero + \epsilon$ where $\epsilon = \sum_{k \geq 1} t_k d_k$. Since $\epsilon$ increases filtration degree, it contributes only to higher pages of the spectral sequence: the $E_1$ page is determined by $\dzero$ alone. By classical BD acyclicity \cite[Theorem~3.4.12]{BD04}, the $E_1$ page is concentrated in a single degree, hence $E_2 = E_\infty$. The quantum corrections modify $d_r$ for $r \geq 2$, but since $E_2^{p,q} = 0$ for $q \neq 0$, these higher differentials vanish for degree reasons. Therefore acyclicity is preserved.

\emph{Step 3: Explicit verification for low genus.}

\emph{Genus 1.} The correction involves the elliptic Weierstrass $\wp$ function. By theta function identities:
\[H^i(C(\mathcal{A}), d_1) = H^i(C(\mathcal{A}), d_0) \otimes_\mathbb{C} \mathcal{O}(\mathfrak{h})\]
where $\mathcal{O}(\mathfrak{h})$ denotes holomorphic functions of $\tau \in \mathfrak{h}$ (the coefficients are quasi-modular forms in $\mathbb{C}[E_2, E_4, E_6]$, not polynomials in $\tau$). This is acyclic since $H^i(C(\mathcal{A}), d_0) = 0$ for $i > 0$.

\emph{Genus 2.} The correction involves Riemann theta functions for the genus-$2$ Jacobian. The same tensor product decomposition applies: $H^i(C(\mathcal{A}), d_2) = H^i(C(\mathcal{A}), d_0) \otimes_\mathbb{C} \mathcal{O}(\mathcal{A}_2)$ where $\mathcal{O}(\mathcal{A}_2)$ denotes holomorphic functions on the Siegel upper half-space of degree~$2$. Acyclicity follows from the genus-$0$ result as in genus~$1$.

\emph{General genus.} By descent from the Torelli space $\mathcal{T}_g \to \mathcal{M}_g$ (which is a covering), acyclicity for $\mathcal{T}_g$ implies acyclicity for $\mathcal{M}_g$.
\end{proof}

\subsection*{Acyclicity at all genera}

\emph{Step 4: Acyclicity via Filtration.}

We prove acyclicity by induction on the Cousin filtration, now accounting for quantum corrections, incorporating Lemmas \ref{lem:relative-diagonal}, \ref{lem:boundary-compatible}, \ref{lem:quantum-preserves-acyclicity}.

\begin{remark}[Higher-genus acyclicity package]\label{rem:summary-higher-genus}
The extension of BD's Chevalley--Cousin acyclicity to all genera uses:
\begin{enumerate}
\item \emph{Relative Ran space:} Working over the universal curve $\mathcal{C}_g \to \mathcal{M}_g$
\item \emph{Boundary compatibility:} Gluing formula at nodes, normal crossings preserved
\item \emph{Quantum corrections:} $\Dg{g} = \dzero + \sum t_k d_k$ with $\Dg{g}^{\,2} = 0$
\item \emph{Acyclicity preservation:} Spectral sequence argument
\end{enumerate}

These four inputs give the relative genus-$g$ Chevalley--Cousin
acyclicity statement.
\end{remark}

\begin{lemma}[Graded piece acyclicity; \ClaimStatusProvedHere]\label{lem:graded-acyclic}
For each $n \geq 0$ and $g \geq 0$:
\[H^i(\mathcal{M}_g \times R(X)_n^o, \text{gr}_n C(\mathcal{A})) =
\begin{cases}
\mathcal{A} & i=0, n=0 \\
0 & \text{otherwise}
\end{cases}\]
where $R(X)_n^o = X^n / \Sigma_n$ is the configuration space of $n$ unordered points.
\end{lemma}

\begin{proof}[Proof of Lemma]
The Leray spectral sequence for $\mathcal{M}_g \times X^n \to \mathcal{M}_g$ gives:
\[E_2^{p,q} = H^p(\mathcal{M}_g) \otimes H^q(X^n, \mathcal{A}^{\boxtimes n})
\Rightarrow H^{p+q}(\mathcal{M}_g \times X^n, \mathcal{A}^{\boxtimes n})\]

For $g=0$: $\mathcal{M}_0 = \text{pt}$, recovering BD's genus-zero result.

For $g=1$: $H^*(\overline{\mathcal{M}}_{1,1}) = \mathbb{Q}[\lambda_1]/(\lambda_1^2)$ where $\lambda_1 = c_1(\mathbb{E})$ is the first Chern class of the Hodge bundle. The quantum corrections enter through this class.

For $g \geq 2$: $\mathcal{M}_g$ has dimension $3g-3$, and its cohomology is generated
by the Hodge classes $\lambda_i \in H^{2i}(\mathcal{M}_g)$.

The key is that $\mathcal{A}^{\boxtimes n}$ is a $D$-module, so by BD \cite[Lemma 4.2.10]{BD04},
its cohomology vanishes in degrees $> n$. Combined with the structure of $H^*(\mathcal{M}_g)$,
this forces the higher cohomology to vanish except in the stated cases.
\end{proof}

\emph{Step 5: Conclusion.}

By Lemmas \ref{lem:boundary-compatible}, \ref{lem:quantum-preserves-acyclicity}, and
\ref{lem:graded-acyclic}, the Chevalley--Cousin complex $C(\mathcal{A})$ over
$\mathcal{M}_{g,n}$ is acyclic for all $g \geq 0$.

Therefore, the descent from $R(X)$ to $X$ extends to all genera with quantum corrections.
\end{proof}

\begin{remark}[Physical interpretation]\label{rem:physical-higher-genus-descent}
The theorem states that configuration space integrals computing correlation functions on $\Sigma_g$ remain well-defined at all genera, provided the quantum corrections (central charges, anomalies) parametrized by $H^*(\overline{\mathcal{M}}_g)$ are included.
\end{remark}

\section[Verdier duality and Ayala-Francis compatibility]{Verdier duality and Ayala--Francis compatibility}
\label{sec:verdier-ayala-francis}

\subsection{Three levels of duality}
\label{subsec:three-dualities}

\begin{definition}[Three duality structures]\label{def:three-dualities}
We work with three dualities simultaneously.

\emph{Verdier duality (geometric).}
For $X$ a smooth variety of dimension $d$, Verdier duality is
$\mathbb{D}_X\colon D^b_c(X) \to D^b_c(X)^{\mathrm{op}}$,
$\mathbb{D}_X(\mathcal{F}) = R\mathcal{H}om(\mathcal{F}, \omega_X[d])$.
It satisfies $\mathbb{D}_X^2 \simeq \operatorname{id}$, Poincar\'e duality
\[
R\Gamma_c(X, \mathbb{D}_X \mathcal{F})
\simeq R\Hom(R\Gamma_c(X, \mathcal{F}), \mathbb{C})^\vee,
\]
and compatibility with proper pushforward
$\mathbb{D}_Y \circ f_* \simeq f_! \circ \mathbb{D}_X$.

\emph{Ayala--Francis duality (topological).}
For an $E_n$-algebra $A$ and manifold $M$,
$\int_M A \simeq \mathbb{D}_M(\int_{-M} A^\vee)$,
where $A^\vee$ is the $E_n$-coalgebra Koszul dual and $-M$ denotes the opposite orientation.
This satisfies oriented involution, gluing
$\int_{M_1 \cup M_2} A \simeq \int_{M_1} A \otimes_{\int_{\partial}} \int_{M_2} A$,
and Poincar\'e--Koszul duality (bar and cobar are dual under integration).

\emph{Linear duality (algebraic).}
For vector spaces, $V^\vee = \Hom(V, \mathbb{C})$.

We show these three dualities are compatible via the de Rham functor.
\end{definition}

\subsection{The de Rham functor: bridge between geometry and topology}
\label{subsec:de-rham-functor}

\begin{definition}[de Rham functor]\label{def:de-rham-functor}
The de Rham functor:
\[\text{DR}: D^b(D\text{-mod}(X)) \to D^b(\text{Vect}_\mathbb{C})\]
is defined by taking global sections followed by de Rham cohomology:
\[\text{DR}(\mathcal{M}) = R\Gamma(X, \Omega^\bullet_X \otimes_{\mathcal{D}_X} \mathcal{M})\]

For a right $\mathcal{D}_X$-module $\mathcal{M}$, this computes
the de Rham cohomology with coefficients in $\mathcal{M}$.
\end{definition}

\begin{proposition}[DR preserves duality structures; \ClaimStatusProvedHere]\label{prop:DR-preserves-duality}
The de Rham functor is compatible with duality in the following sense:
\[\text{DR}(\mathbb{D}_X \mathcal{M}) \simeq \text{DR}(\mathcal{M})^\vee[-d]\]
where $d = \dim X$ and $(-)^\vee$ is linear duality.
\end{proposition}

\begin{proof}
\emph{Step 1: Verdier duality on $\mathcal{D}$-modules.}

For a $\mathcal{D}_X$-module $\mathcal{M}$:
\[\mathbb{D}_X(\mathcal{M}) = R\mathcal{H}om_{\mathcal{D}_X}(\mathcal{M}, \mathcal{D}_X \otimes_{\mathcal{O}_X} \omega_X[d])\]

\emph{Step 2: Apply de Rham.}

\begin{align*}
\text{DR}(\mathbb{D}_X \mathcal{M})
&= R\Gamma(X, \Omega^\bullet_X \otimes_{\mathcal{D}_X} R\mathcal{H}om_{\mathcal{D}_X}(\mathcal{M}, \mathcal{D}_X \otimes \omega_X[d])) \\
&\simeq R\Gamma(X, R\mathcal{H}om_{\mathcal{D}_X}(\mathcal{M}, \Omega^\bullet_X \otimes \omega_X[d])) && \text{(tensor-hom adjunction)} \\
&\simeq R\Hom_{\mathcal{D}_X}(\mathcal{M}, R\Gamma(X, \Omega^\bullet_X \otimes \omega_X[d])) && \text{(global sections)} \\
&\simeq R\Hom(\text{DR}(\mathcal{M}), \mathbb{C})[-d] && \text{(Serre duality)}
\end{align*}

The last step uses Serre duality: $R\Gamma(X, \omega_X[d]) \simeq \mathbb{C}$ for proper smooth $X$.

Therefore: $\text{DR}(\mathbb{D}_X \mathcal{M}) \simeq \text{DR}(\mathcal{M})^\vee[-d]$. \qedhere
\end{proof}

\begin{theorem}[Geometric-topological duality compatibility; \ClaimStatusProvedHere]\label{thm:verdier-AF-compat}
The Verdier duality functor on configuration spaces:
\[\mathbb{D}_{\text{Conf}_n(X)}: D^b(\text{Conf}_n(X)) \to D^b(\text{Conf}_n(X))^{op}\]
is compatible with the Ayala--Francis factorization homology duality via the de Rham functor:
\[\text{DR}: D\text{-mod}(X) \to \text{Vect}_\mathbb{C}\]

\emph{Precise Statement.} The following diagram commutes up to canonical isomorphism:
\[
\begin{tikzcd}
D^b(D\text{-mod}(\text{Conf}_n(X))) \arrow[r, "\mathbb{D}"] \arrow[d, "\text{DR}"]
& D^b(D\text{-mod}(\text{Conf}_n(X)))^{op} \arrow[d, "\text{DR}"] \\
D^b(\text{Vect}_\mathbb{C}) \arrow[r, "\text{AF-dual}"] & D^b(\text{Vect}_\mathbb{C})^{op}
\end{tikzcd}
\]
\end{theorem}

\begin{proof}

\subsection*{Part 1: setup and notation}

For any $\mathcal{D}$-module $\mathcal{M}$ on $\mathrm{Conf}_n(X)$, we prove:
\[\text{DR}(\mathbb{D}_{\text{Conf}_n(X)}(\mathcal{M})) \simeq \text{AF-dual}(\text{DR}(\mathcal{M}))\]
where the left side uses Verdier duality and the right side uses Ayala--Francis (topological) duality.

Both sides compute a form of ``dual'' of $\text{DR}(\mathcal{M})$; we must show they give the same result.

\subsection*{Part 2: Verdier duality on configuration spaces}

\begin{lemma}[Verdier dual of chiral algebra; \ClaimStatusProvedHere]\label{lem:verdier-dual-chiral}
For a chiral algebra $\mathcal{A}$ on $X$, consider the $\mathcal{D}$-module:
\[\mathcal{M}_n = \mathcal{A}^{\boxtimes n} \otimes \Omega^k_{\log}(\text{Conf}_n(X))\]
on the configuration space $\text{Conf}_n(X)$. Its Verdier dual is:
\[\mathbb{D}_{\text{Conf}_n(X)}(\mathcal{M}_n) \simeq (\mathcal{A}^\vee)^{\boxtimes n} \otimes \Omega^{2n-k}_c(\text{Conf}_n(X))\]
where $\mathcal{A}^\vee$ is the linear dual of $\mathcal{A}$ and $\Omega^{2n-k}_c$ are compactly supported forms
($\dim_{\mathbb{R}}\operatorname{Conf}_n(X) = 2n$, so
$k + (2n{-}k) = 2n$).
\end{lemma}

\begin{proof}[Proof of Lemma]
\emph{Step 1: Dimension of configuration space.}

For a smooth curve $X$ of complex dimension~$1$, the configuration space $\operatorname{Conf}_n(X) = X^n \setminus \{\text{diagonals}\}$ has complex dimension~$n$ (real dimension~$2n$).

\emph{Step 2: Dualizing complex.}

The dualizing complex of $\text{Conf}_n(X)$ is:
\[\omega_{\text{Conf}_n(X)} = \omega_X^{\boxtimes n}[n]\]
(product of dualizing complexes, shifted).

\emph{Step 3: Compute Verdier dual.}

\begin{align*}
\mathbb{D}(\mathcal{M}_n)
&= R\mathcal{H}om(\mathcal{M}_n, \omega_{\text{Conf}_n(X)}) \\
&= R\mathcal{H}om(\mathcal{A}^{\boxtimes n} \otimes \Omega^k_{\log}, \omega_X^{\boxtimes n}[n]) \\
&\simeq (\mathcal{A}^\vee)^{\boxtimes n} \otimes R\mathcal{H}om(\Omega^k_{\log}, \omega_X^{\boxtimes n}[n]) && \text{(tensor-hom)} \\
&\simeq (\mathcal{A}^\vee)^{\boxtimes n} \otimes \Omega^{2n-k}_c[n] && \text{(forms with compact support)}
\end{align*}

The last step uses the pairing between logarithmic forms and compactly supported forms. \qedhere
\end{proof}

\subsection*{Part 3: Ayala--Francis duality on factorization homology}

\begin{lemma}[AF duality for chiral algebras; \ClaimStatusProvedHere]\label{lem:AF-dual-chiral}
The Ayala--Francis duality for a chiral algebra $\mathcal{A}$ is:
\[\int_{\text{Conf}_n(X)} \mathcal{A} \simeq \mathbb{D}_{\text{top}}\left(\int_{-\text{Conf}_n(X)} \bar{B}(\mathcal{A})\right)\]
where $\bar{B}(\mathcal{A})$ is the bar coalgebra computing the dual
coalgebra $\mathcal{A}^i = H^*(\barB(\mathcal{A}))$, and
$\mathbb{D}_{\text{top}}$ is topological (linear) duality.
\end{lemma}

\begin{proof}[Proof of Lemma]
This is Ayala--Francis Theorem 4.5~\cite{AF15}.

\emph{Step 1: Factorization homology as colimit.}
\[\int_{\text{Conf}_n(X)} \mathcal{A} = \colim_{U_1 \sqcup \cdots \sqcup U_n \subset X} \mathcal{A}(U_1) \otimes \cdots \otimes \mathcal{A}(U_n)\]

\emph{Step 2: Bar construction as limit.}
\[\int_{-\text{Conf}_n(X)} \bar{B}(\mathcal{A}) = \lim_{U_1 \sqcup \cdots \sqcup U_n \subset X} \bar{B}(\mathcal{A})(U_1) \otimes \cdots \otimes \bar{B}(\mathcal{A})(U_n)\]

\emph{Step 3: Duality interchanges colim and lim.}
\[\mathbb{D}_{\text{top}}(\lim) \simeq \colim(\mathbb{D}_{\text{top}})\]

Therefore: $\mathbb{D}_{\text{top}}\left(\int_{-\text{Conf}_n(X)} \bar{B}(\mathcal{A})\right) \simeq \int_{\text{Conf}_n(X)} \mathcal{A}$. \qedhere
\end{proof}

\subsection*{Part 4: the de Rham functor intertwines the dualities}

Now we show that DR makes the two dualities compatible.

\begin{proposition}[Key compatibility; \ClaimStatusProvedHere]\label{prop:key-compat-DR}
The following diagram commutes:
\[
\begin{tikzcd}
\mathcal{M}_n \arrow[r, "\mathbb{D}_X"] \arrow[d, "\text{DR}"]
& \mathbb{D}_X(\mathcal{M}_n) \arrow[d, "\text{DR}"] \\
\int_{\text{Conf}_n(X)} \mathcal{A} \arrow[r, "\text{AF-dual}"]
& \mathbb{D}_{\text{top}}\left(\int_{-\text{Conf}_n(X)} \bar{B}(\mathcal{A})\right)
\end{tikzcd}
\]
\end{proposition}

\begin{proof}[Proof of Proposition]
\emph{Step 1: Apply DR to Verdier dual (left-then-down).}

From Lemma~\ref{lem:verdier-dual-chiral}:
\[\mathbb{D}_X(\mathcal{M}_n) \simeq (\mathcal{A}^\vee)^{\boxtimes n} \otimes \Omega^{2n-k}_c\]

Applying DR:
\begin{align*}
\text{DR}(\mathbb{D}_X(\mathcal{M}_n))
&= R\Gamma(\text{Conf}_n(X), \Omega^\bullet \otimes_{\mathcal{D}} ((\mathcal{A}^\vee)^{\boxtimes n} \otimes \Omega^{2n-k}_c)) \\
&\simeq R\Gamma_c(\text{Conf}_n(X), \mathcal{A}^\vee)^{\boxtimes n} && \text{(Poincaré duality)} \\
&\simeq \left(R\Gamma(\text{Conf}_n(X), \mathcal{A})^{\boxtimes n}\right)^\vee && \text{(linear duality)}
\end{align*}

\emph{Step 2: Apply AF-dual to DR (down-then-right).}

From the definition of factorization homology:
\[\text{DR}(\mathcal{M}_n) \simeq \int_{\text{Conf}_n(X)} \mathcal{A}\]

Applying AF-dual (Lemma~\ref{lem:AF-dual-chiral}):
\[\text{AF-dual}\left(\int_{\text{Conf}_n(X)} \mathcal{A}\right) \simeq \mathbb{D}_{\text{top}}\left(\int_{-\text{Conf}_n(X)} \bar{B}(\mathcal{A})\right)\]

\emph{Step 3: Compare the two paths.}

We need to show:
\[\left(R\Gamma(\text{Conf}_n(X), \mathcal{A})^{\boxtimes n}\right)^\vee \simeq \mathbb{D}_{\text{top}}\left(\int_{-\text{Conf}_n(X)} \bar{B}(\mathcal{A})\right)\]

By the Verdier intertwining (Theorem~\ref{thm:bar-cobar-verdier}):
\[\mathbb{D}_{\operatorname{Ran}}\bar{B}(\mathcal{A}) \simeq \cA^!_\infty
\qquad\text{(factorization \emph{algebra}, not coalgebra)}\]

Therefore:
\begin{align*}
\mathbb{D}_{\text{top}}\left(\int_{-\text{Conf}_n(X)} \bar{B}(\mathcal{A})\right)
&\simeq \mathbb{D}_{\text{top}}\left(\int_{-\text{Conf}_n(X)} \mathcal{A}^\vee\right) \\
&\simeq \left(\int_{\text{Conf}_n(X)} \mathcal{A}\right)^\vee && \text{(AF duality)} \\
&\simeq \left(R\Gamma(\text{Conf}_n(X), \mathcal{A})^{\boxtimes n}\right)^\vee && \text{(definition of factorization homology)}
\end{align*}

The two paths agree. \qedhere
\end{proof}

\subsection*{Part 5: full theorem conclusion}

Combining Proposition~\ref{prop:key-compat-DR} with Lemmas \ref{lem:verdier-dual-chiral} and \ref{lem:AF-dual-chiral}, we have proven that the diagram in the theorem statement commutes.

In particular, Verdier duality on $\mathcal{D}$-modules, Ayala--Francis duality on factorization algebras, and linear duality on vector spaces are all compatible via the de Rham functor.
\end{proof}

\begin{corollary}[Bar complex computes factorization cohomology; \ClaimStatusProvedHere]\label{cor:bar-is-fh}
The geometric bar complex $\bar{B}^{\text{geom}}(\mathcal{A})$ computes factorization homology:
\[\text{DR}(\bar{B}^{\text{geom}}(\mathcal{A})) \simeq \int_X \mathcal{A}\]
\end{corollary}

\begin{proof}
By the compatibility theorem just proved, the de Rham functor intertwines Verdier duality on $\mathcal{D}$-modules with Ayala--Francis duality on factorization algebras. The bar complex is defined as $\bar{B}^{\mathrm{geom}}(\mathcal{A}) = \bigoplus_n \Gamma(\overline{C}_{n+1}(X), \mathcal{A}^{\boxtimes(n+1)} \otimes \Omega^n_{\log})$, which under the de Rham functor becomes the factorization homology $\int_X \mathcal{A}$: both compute the derived tensor product $\mathcal{A} \otimes^{\mathbb{L}}_{\mathrm{Disk}(X)} \mathrm{pt}$ (Ayala--Francis~\cite{AF15}, Theorem~3.24; for the identification with our bar complex, see Theorem~\ref{thm:fact-homology-quantum}).
\end{proof}

\begin{remark}\label{rem:why-verdier-AF-matters}
This compatibility theorem ensures the construction is not ad hoc. It shows:

\begin{enumerate}
\item \emph{Geometric bar-cobar} (via configuration space integrals and Verdier duality)
\item \emph{Topological factorization homology} (via $E_\infty$ operads and Ayala--Francis)
\item \emph{Algebraic Koszul duality} (via bar-cobar adjunction)
\end{enumerate}

are all manifestations of the same underlying structure. The de Rham functor bridges geometry and topology, making the three perspectives equivalent.
\end{remark}

\subsection{Verification}

\emph{Setup: Three Levels of Duality}

We must reconcile three different notions of duality:

\begin{enumerate}
\item \emph{Verdier duality (geometric):} For $X$ smooth, $\mathbb{D}_X: D^b_c(X)
\to D^b_c(X)$ sends sheaf $\mathcal{F}$ to $\mathbb{D}_X(\mathcal{F}) =
R\mathcal{H}om(\mathcal{F}, \omega_X[\dim X])$.

\item \emph{Ayala--Francis duality (topological):} For $E_n$-algebras $A$,
factorization homology $\int_M A$ has a dual $\int_M A^{\vee}$ where $A^{\vee}$
is the $E_n$-coalgebra dual.

\item \emph{Linear duality (algebraic):} For vector spaces $V$, $V^* =
\text{Hom}(V, \mathbb{C})$.
\end{enumerate}

\emph{Step 1: De Rham Functor as Bridge}

The de Rham functor is defined by:
\[\text{DR}(\mathcal{M}) = R\Gamma(X, \Omega^{\bullet}_X \otimes_{\mathcal{D}_X} \mathcal{M})\]
for $\mathcal{M} \in D\text{-mod}(X)$.

\begin{lemma}[De Rham and Verdier duality; \ClaimStatusProvedHere]\label{lem:DR-verdier-compat}
For $\mathcal{M} \in D^b_c(D\text{-mod}(X))$ with $X$ smooth:
\[\text{DR}(\mathbb{D}_X(\mathcal{M})) \simeq \text{DR}(\mathcal{M})^* [\dim X]\]
where $(-)^*$ is linear duality.
\end{lemma}

\begin{proof}[Proof of Lemma]
This is a classical result in $D$-module theory (Kashiwara-Schapira \cite{KS90}).
The key steps are:

1. By definition, $\mathbb{D}_X(\mathcal{M}) = R\mathcal{H}om_{\mathcal{D}_X}(\mathcal{M},
\omega_X[\dim X])$.

2. The de Rham complex of $\mathbb{D}_X(\mathcal{M})$ is:
\[\text{DR}(\mathbb{D}_X(\mathcal{M})) = R\Gamma(X, \Omega^{\bullet}_X
\otimes_{\mathcal{D}_X} R\mathcal{H}om_{\mathcal{D}_X}(\mathcal{M}, \omega_X[\dim X]))\]

3. By adjunction:
\[\Omega^{\bullet}_X \otimes_{\mathcal{D}_X} R\mathcal{H}om_{\mathcal{D}_X}(\mathcal{M},
\omega_X[\dim X]) \simeq R\mathcal{H}om(\Omega^{\bullet}_X \otimes_{\mathcal{D}_X}
\mathcal{M}, \omega_X[\dim X])\]

4. Using Serre duality:
\[R\Gamma(X, R\mathcal{H}om(\text{DR}(\mathcal{M}), \omega_X[\dim X]))
\simeq R\Gamma(X, \text{DR}(\mathcal{M}))^* [\dim X]\]
\end{proof}

\emph{Step 2: Configuration Spaces and Ran Space}

For configuration spaces, we must be more careful. The Ran space $\text{Ran}(X)$ is:
\[\text{Ran}(X) = \text{colim}_{n} X^{(n)}\]
where $X^{(n)} = X^n / \Sigma_n$ is the $n$-fold symmetric product.

Beilinson--Drinfeld \cite{BD04} show that chiral algebras are factorization algebras on
$\text{Ran}(X)$. Ayala--Francis \cite{AF15} work with factorization algebras on manifolds.

The connection is through the Riemann--Hilbert correspondence:
\[\text{RH}: D\text{-mod}(X) \xrightarrow{\sim} \text{Local systems}(X^{an})\]

\begin{lemma}[Ran space duality; \ClaimStatusProvedHere]\label{lem:ran-duality-AF}
Verdier duality on $D\text{-mod}(\text{Ran}(X))$ corresponds under $\text{RH}$ to
Ayala--Francis duality on factorization algebras on $X^{an}$.
\end{lemma}

\begin{proof}[Proof of Lemma]
The Riemann--Hilbert correspondence extends to the Ran space by taking colimits:
\[\text{RH}: D\text{-mod}(\text{Ran}(X)) \xrightarrow{\sim}
\text{Fact}(X^{an}, \text{Vect}_\mathbb{C})\]

For $\mathcal{A} \in \text{ChirAlg}(X)$, view it as a $D$-module on $\text{Ran}(X)$. Then:
\[\text{RH}(\mathcal{A}) = \text{forget structure}(\mathcal{A})\]
is its underlying local system.

Ayala--Francis duality for $\text{RH}(\mathcal{A})$ is defined by:
\[\int_M \text{RH}(\mathcal{A})^{\vee} := \left(\int_M \text{RH}(\mathcal{A})\right)^*\]

On the $D$-module side, this is:
\[\mathbb{D}_{\text{Ran}(X)}(\mathcal{A}) = R\mathcal{H}om_{\mathcal{D}}(\mathcal{A},
\omega_{\text{Ran}(X)})\]

By Lemma~\ref{lem:DR-verdier-compat}, applying $\text{DR}$ to both sides gives the same result.
\end{proof}

\emph{Step 3: Configuration Space Level}

Now specialize to $\text{Conf}_n(X) = X^n \setminus \Delta$. We have:
\[\bar{B}^n(\mathcal{A}) = \int_{\text{Conf}_n(X)} \mathcal{A}^{\boxtimes n}\]

\begin{lemma}[Bar as factorization homology; \ClaimStatusProvedHere]\label{lem:bar-as-fact-hom-AF}
The bar construction computes factorization homology:
\[\bar{B}(\mathcal{A}) \simeq \int_{X} \mathcal{A}\]
in the sense of Ayala--Francis \cite{AF15}.
\end{lemma}

\begin{proof}[Proof of Lemma]
By Ayala--Francis \cite[Theorem 4.19]{AF15}, factorization homology is computed by:
\[\int_X \mathcal{A} = \text{colim}_n \left(\int_{\text{Conf}_n(X)}
\mathcal{A}^{\boxtimes n}\right)^{\Sigma_n}\]

This is the bar construction:
\[\bar{B}^n(\mathcal{A}) = \Gamma(\overline{C}_n(X), j_* j^* \mathcal{A}^{\boxtimes n}
\otimes \Omega^n_{\log})\]

The logarithmic forms $\Omega^n_{\log}$ provide the integration measure, and taking
$\Sigma_n$-coinvariants gives the symmetric quotient.
\end{proof}

\emph{Step 4: Homotopy Koszul Dual via Verdier Duality}

The homotopy Koszul dual $\cA^!_\infty$ is characterized by
Verdier duality on the bar coalgebra
(Convention~\ref{conv:bar-coalgebra-identity}):
\[\cA^!_\infty \simeq \mathbb{D}_{\text{Ran}(X)}(\bar{B}(\mathcal{A}))\]
The strict Koszul dual $\cA^!$ is recovered on cohomology
by Theorem~\ref{thm:chiral-koszul-duality}.

\begin{lemma}[Algebra structure from Verdier dual; \ClaimStatusProvedHere]\label{lem:coalgebra-verdier-AF}
Under $\text{DR}$, the algebra structure on $\cA^!_\infty$ comes from
applying Verdier duality to the coalgebra
$\bar{B}(\mathcal{A})$: Verdier duality converts the
factorization coalgebra $\bar{B}(\cA)$ into the factorization
algebra $\cA^!_\infty$.
\end{lemma}

\begin{proof}[Proof of Lemma]
The multiplication $\mu: \mathcal{A} \otimes \mathcal{A} \to \mathcal{A}$ dualizes to:
\[\mathbb{D}(\mu): \mathbb{D}(\mathcal{A}) \to \mathbb{D}(\mathcal{A} \otimes \mathcal{A})
\simeq \mathbb{D}(\mathcal{A}) \otimes \mathbb{D}(\mathcal{A})\]

Applying $\text{DR}$:
\[\text{DR}(\mathbb{D}(\mu)): \text{DR}(\mathcal{A})^* \to \text{DR}(\mathcal{A})^*
\otimes \text{DR}(\mathcal{A})^*\]

This is the coproduct structure on the coalgebra.

In factorization terms, this says:
\[\Delta: \int_X \mathcal{A}^! \to \int_{X \sqcup X} \mathcal{A}^!
= \left(\int_X \mathcal{A}^!\right) \otimes \left(\int_X \mathcal{A}^!\right)\]

which is the Ayala--Francis coalgebra structure.
\end{proof}

\emph{Step 5: Full Compatibility}

\begin{lemma}[Diagram commutes; \ClaimStatusProvedHere]\label{lem:diagram-commutes-AF}
The diagram in the theorem statement commutes up to natural isomorphism.
\end{lemma}

\begin{proof}[Proof of Lemma]
By Lemmas \ref{lem:DR-verdier-compat}, \ref{lem:ran-duality-AF}, \ref{lem:bar-as-fact-hom-AF}, and
\ref{lem:coalgebra-verdier-AF}, we have:
\[\text{DR}(\mathbb{D}(\mathcal{A})) \simeq \text{DR}(\mathcal{A})^*
\simeq \text{AF-dual}(\text{DR}(\mathcal{A}))\]

The naturality in $\mathcal{A}$ ensures this is a natural isomorphism of functors.
\end{proof}

\begin{remark}[Compatibility with chiral Koszul duality]\label{rem:importance-koszul-AF}
The geometric construction (Verdier duality on configuration spaces) and the topological construction (Ayala--Francis duality on factorization algebras) produce equivalent results.

Without this compatibility, we could not be sure that the ``dual coalgebra'' we construct
geometrically is the same as the ``Koszul dual'' in the abstract algebraic sense.

The theorem provides the bridge: geometric duality via $D$-modules corresponds to
topological duality via factorization homology, both giving the same Koszul dual algebra.
\end{remark}

With Verdier duality and Ayala--Francis compatibility established, the bar-cobar adjunction at higher genus can be inverted on the Koszul locus.

%================================================================
% HIGHER GENUS QUASI-ISOMORPHISM
%================================================================

\section{Bar-cobar quasi-isomorphism at higher genus}
\label{sec:bar-cobar-qi-higher-genus}

\begin{lemma}[Open-stratum quasi-isomorphism; \ClaimStatusConditional]\label{lem:higher-genus-open-stratum-qi}
Let $U_g \subset \overline{\mathcal{M}}_g$ be the open stratum of smooth genus-$g$
curves. The restriction
\[j_g^*\psi_g: j_g^*\Omega_g(\bar{B}_g(\mathcal{A})) \longrightarrow j_g^*\mathcal{A}_g\]
is a quasi-isomorphism.
\end{lemma}

\begin{proof}
Fix a geometric point $s \in U_g$, corresponding to a smooth
genus-$g$ curve~$X_s$. The fiber of $j_g^*\psi_g$ at~$s$ is
the bar-cobar counit
$\psi_s\colon \Omega(\bar{B}(\cA|_{X_s})) \to \cA|_{X_s}$
on the \emph{fixed} curve~$X_s$.
By axiom~\ref{MK:modular}
(Definition~\ref{def:modular-koszul-chiral}), the PBW spectral
sequence of $(\bar{B}^{(g)}(\cA), \dfib)$ has concentrated
$E_\infty$ page at genus~$g$. This concentration is verified
independently of bar-cobar inversion: for Heisenberg by
Theorem~\ref{thm:heisenberg-higher-genus}, for affine
Kac--Moody by Theorem~\ref{thm:pbw-allgenera-km}, for Virasoro
by Theorem~\ref{thm:pbw-allgenera-virasoro}, and for principal
$\mathcal{W}$-algebras by
Theorem~\ref{thm:pbw-allgenera-principal-w}.

Given this concentration, the fundamental theorem of chiral
twisting morphisms
(Theorem~\ref{thm:fundamental-twisting-morphisms})
applies to the bar complex on~$X_s$: the bar-cobar
spectral sequence collapses at~$E_2$, and the counit~$\psi_s$
is a quasi-isomorphism.

The cone $C = \operatorname{Cone}(j_g^*\psi_g)$ is therefore a complex of
constructible sheaves on~$U_g$ whose stalk cohomology vanishes at every
geometric point (by the fiberwise quasi-isomorphism just established).
Constructibility of~$C$ is inherited from that of the fiber
cohomology sheaves on~$U_g$
(established in the proof of Theorem~\ref{thm:ss-quantum},
Step~2).
The cohomology sheaves $\mathcal{H}^q(C)$ are therefore
constructible sheaves on the smooth connected stack~$U_g$
whose stalks $\mathcal{H}^q(C)_s = H^q(C_s) = 0$ vanish at
every geometric point~$s$.
A constructible sheaf on a connected base whose stalks are all
zero is the zero sheaf
(see~\cite[Expos\'e~IX, Th\'eor\`eme~2.14]{SGA4});
hence $\mathcal{H}^q(C) = 0$ for all~$q$, which means $C$ is
acyclic and $j_g^*\psi_g$ is a quasi-isomorphism on~$U_g$.
\end{proof}

\begin{lemma}[Boundary-stratum compatibility of \texorpdfstring{$\psi_g$}{psi-g}; \ClaimStatusConditional]\label{lem:higher-genus-boundary-qi}
For every boundary stratum $D_\Gamma \subset \partial\overline{\mathcal{M}}_g$
indexed by a stable graph $\Gamma$, the restricted map
\[\psi_g|_{D_\Gamma}: \Omega_g(\bar{B}_g(\mathcal{A}))|_{D_\Gamma}
\longrightarrow \mathcal{A}_g|_{D_\Gamma}\]
is a quasi-isomorphism, provided the higher-genus inversion statement is known for
all genera $< g$.
\end{lemma}

\begin{proof}
By Proposition~\ref{prop:gluing-at-nodes} and
Lemma~\ref{lem:boundary-compatible}, restriction of the bar data to $D_\Gamma$
factorizes over nodal gluings into tensor products of lower-genus pieces.
Functoriality of bar-cobar inversion
(Theorem~\ref{thm:bar-cobar-inversion-functorial}) identifies
$\psi_g|_{D_\Gamma}$ with the corresponding tensor product of lower-genus inversion maps.

Every vertex genus in $\Gamma$ is strictly smaller than $g$: in the separating case,
$g=g_1+g_2$ with $g_1,g_2 \geq 0$ and each $g_i < g$ (note: one component may have genus~$0$, which is the base case of the induction); in the non-separating case, the single vertex has genus $g-1 < g$. In either case, the induction hypothesis applies to each vertex. Tensor products of quasi-isomorphisms
are quasi-isomorphisms by the K\"unneth formula for factorization
algebras at nodes: the factorization tensor product at a node is
computed fiberwise over the normalization, reducing to the plain
tensor product of chain complexes over the ground field, which
preserves quasi-isomorphisms by K\"unneth. This yields the claim.
\end{proof}

\begin{lemma}[Extension across boundary; \ClaimStatusProvedHere]\label{lem:extension-across-boundary-qi}
Let $X$ be a stack with open-closed decomposition $X = U \sqcup Z$ (with
$j:U\hookrightarrow X$, $i:Z\hookrightarrow X$). For a morphism of complexes
$f:K\to L$ on $X$, if $j^*f$ and $i^*f$ are quasi-isomorphisms, then $f$ is a
quasi-isomorphism.
\end{lemma}

\begin{proof}
Let $C=\operatorname{Cone}(f)$. It suffices to prove $C$ is acyclic.
From $j^*f$ and $i^*f$ quasi-isomorphisms, $j^*C$ and $i^*C$ are acyclic.
Apply the localization triangle
\[j_!j^*C \longrightarrow C \longrightarrow i_*i^*C \xrightarrow{+1}.\]
Both outer terms are acyclic, hence $C$ is acyclic. Therefore $f$ is a
quasi-isomorphism.
\end{proof}

\begin{theorem}[Higher genus inversion; \ClaimStatusConditional]\label{thm:higher-genus-inversion}
\index{bar-cobar inversion!higher genus}
\textup{[}Regime: curved-central on the Koszul locus
\textup{(}Convention~\textup{\ref{conv:regime-tags})].}

\smallskip\noindent
The genus-$0$ inversion of~\S\ref{sec:frame-inversion} extends to all
genera on the Koszul locus. At the scalar level, quantum corrections
are controlled by $\kappa(\cA) \cdot \lambda_g$ on the
uniform-weight lane
\textup{(}Theorem~\textup{\ref{thm:modular-characteristic})}; for
multi-weight algebras at $g \geq 2$, the scalar formula receives the
cross-channel correction of
Theorem~\textup{\ref{thm:multi-weight-genus-expansion}}.

Let $\cA$ be a modular pre-Koszul chiral algebra
\textup{(}Definition~\textup{\ref{def:modular-koszul-chiral}},
axioms \textup{MK1--MK3)}.
\begin{enumerate}[label=\textup{(\alph*)}]
\item \emph{Koszul locus.}
For each genus $g \geq 0$, the bar-cobar counit is a
quasi-isomorphism:
\[\psi_g \colon \Omega_g(\bar{B}_g(\cA))
\xrightarrow{\;\sim\;} \cA_g,\]
where $\cA_g$ denotes the genus-$g$ component.
The spectral sequence $E_1^{p,q}(g) \Rightarrow
H^{p+q}(\Omega_g \bar{B}_g(\cA))$ collapses at $E_2$.
\end{enumerate}
\end{theorem}

\begin{proof}
Induct on $g$.

\emph{Base case ($g=0$).}
The genus-$0$ bar-cobar inversion is proved independently by
Theorem~\ref{thm:bar-nilpotency-complete} ($d^2 = 0$) and
Theorem~\ref{thm:chiral-koszul-duality} (Koszul quasi-isomorphism).
The induction below uses only these genus-$0$ inputs and the
lower-genus induction hypothesis.

\emph{Inductive step.}
Fix $g>0$ and assume the claim for all genera $<g$. Set
$X=\overline{\mathcal{M}}_g$, with $j:U_g\hookrightarrow X$ and
$i:\partial X\hookrightarrow X$.
\begin{enumerate}[label=(\arabic*)]
\item By Lemma~\ref{lem:higher-genus-open-stratum-qi}, $j^*\psi_g$ is a
quasi-isomorphism.
\item By Lemma~\ref{lem:higher-genus-boundary-qi} and the induction hypothesis
on the lower-genus vertices of each stable graph $\Gamma$ indexing a boundary
stratum $D_\Gamma$, $\psi_g|_{D_\Gamma}$ is a quasi-isomorphism for every
$D_\Gamma$. The passage from ``quasi-isomorphism on each~$D_\Gamma$''
to ``quasi-isomorphism on all of~$\partial\overline{\mathcal{M}}_g$''
follows by a Mayer--Vietoris argument on the normal crossing boundary:
the boundary strata form a simple normal crossing divisor, so the
cone of~$i^*\psi_g$ is computed by iterated extensions along the
pairwise intersections $D_\Gamma \cap D_{\Gamma'}$, each of which
is a deeper boundary stratum where the same inductive hypothesis
applies. Hence $i^*\psi_g$ is a quasi-isomorphism.
\item Apply Lemma~\ref{lem:extension-across-boundary-qi} to
$X=U_g\sqcup\partial X$ and $f=\psi_g$. Since $j^*f$ and $i^*f$ are
quasi-isomorphisms by (1)--(2), $f$ is a quasi-isomorphism.
\end{enumerate}
Therefore $\psi_g$ is a quasi-isomorphism on $\overline{\mathcal{M}}_g$.
The asserted $E_2$-collapse is proved separately in
Lemma~\ref{lem:e2-collapse-higher-genus}.
\end{proof}

\begin{remark}[Off-locus coderived continuation kept separate]
\index{bar-cobar inversion!off Koszul locus}
Outside the Koszul locus \textup{(}admissible affine levels,
minimal-model central charges\textup{)}, the filtered curved object
$\Omega_X \bar{B}_X(\cA)$ still exists and its failure to invert in the
ordinary derived category is measured by the shadow obstruction tower
$\Theta_{\cA}^{\leq r}$
\textup{(}Definition~\textup{\ref{def:shadow-postnikov-tower},}
equation~\eqref{eq:universal-MC}\textup{)}. The present file does not
prove that the counit becomes an equivalence in the provisional
coderived category
$\mathsf{D}^{\mathrm{co}}_{\mathrm{prov}}(X)$; that off-locus upgrade is
left to the separate coderived-model formalism
\textup{(}Definition~\textup{\ref{def:provisional-coderived}}\textup{)}
and is intentionally not folded into
Theorem~\ref{thm:higher-genus-inversion}, so no circular dependence on
Proposition~\ref{prop:coderived-adequacy}(a) remains.
\end{remark}

\begin{remark}[Geometric substrate (Volume~II)]
\label{rem:theorem-b-lagrangian}
\index{Lagrangian self-intersection!Theorem B}
Volume~II compares Koszul inversion with the recovery
of a Lagrangian from a clean self-intersection datum: the derived
self-intersection groupoid
$\mathfrak{S} = \mathcal{L} \times_{\mathcal{M}} \mathcal{L}$
is \'etale precisely on the Koszul locus, and an \'etale groupoid
determines its base up to equivalence.
\end{remark}

\begin{remark}[Structural explanation: the Getzler--Kapranov involution]
\label{rem:inversion-from-ft-squared}
\index{bar-cobar inversion!from FT squared}
The above inductive proof verifies bar-cobar inversion stratum by
stratum. The \emph{structural reason} it holds is the
Getzler--Kapranov involution
$\mathrm{FT}^2 \simeq \mathrm{id}$
(Theorem~\ref{thm:feynman-involution}): since
$\barB = \mathrm{FT}(\cA)$ and
$\Omega(\barB) = \mathrm{FT}(\mathrm{FT}(\cA)) \simeq \cA$
on the homotopy category of modular operad algebras satisfying
Poincar\'e duality
(Corollary~\ref{cor:feynman-duality-qch}), bar-cobar inversion
at all genera is a formal consequence of the modular operad axioms.
The inductive proof above is the explicit verification that the
abstract homotopy equivalence $\mathrm{FT}^2 \simeq \mathrm{id}$
descends to a concrete quasi-isomorphism on each stratum of
$\overline{\mathcal{M}}_g$.
\end{remark}

\begin{lemma}[\texorpdfstring{$E_2$}{E2} collapse at higher genus;
\ClaimStatusProvedHere]\label{lem:e2-collapse-higher-genus}
\index{spectral sequence!E2 collapse@$E_2$ collapse!higher genus}
For a Koszul chiral algebra~$\cA$ on a genus-$g$ curve, the PBW
spectral sequence
$E_1^{p,q}(g) \Rightarrow H^{p+q}(\Omega_g \bar{B}_g(\cA))$
degenerates at~$E_2$.
\end{lemma}

\begin{proof}
The PBW filtration is a filtration by chiral algebras, and the
Koszulness hypothesis ensures that the associated graded is a
Koszul complex in the classical sense. The $d_r$ differential
for $r \geq 2$ maps between PBW-graded components separated
by~$r$ steps, and the Koszul concentration (all bar cohomology
is in bar-degree~$0$ on the associated graded) forces these maps
to have zero source or zero target. This is
Corollary~\ref{cor:bar-cohomology-koszul-dual}, Step~4, applied
fiberwise over~$\overline{\mathcal{M}}_g$: on each geometric
fiber (a fixed smooth or stable curve), the collision differential
is genus-$0$ type, and the genus-$g$ quantum corrections live in
higher Leray degrees, contributing only to $d_r$ for $r \geq 2$;
since~$E_2$ is already concentrated on the diagonal by Koszul
concentration, these higher differentials vanish.
\end{proof}

\begin{remark}[Homotopy-native formulation]\label{rem:homotopy-native-b}
The counit $\psi_g\colon\Omega_g(\bar{B}_g(\cA))\xrightarrow{\sim}\cA_g$
is a QI of $\Ainf$-chiral algebras.
The genus induction is the shadow of factorization-homology descent
(excision, Ayala--Francis~\cite{AF15}); $E_2$-collapse on the Koszul
locus corresponds to formality of the filtered $\Ainf$-algebra.
Off the Koszul locus, the statement lives in
$\cat{D}^{\mathrm{co}}_\infty(\cat{Fact}_g(X))$.
\end{remark}

\begin{remark}[Elementary model presentation;
Convention~\ref{conv:proof-architecture}]
\label{rem:theorem-b-model}
\index{bar-cobar inversion!model presentation}
\emph{Step~B} (M-level): induction on genus via open-stratum QI
(Lemma~\ref{lem:higher-genus-open-stratum-qi}), boundary QI
(Lemma~\ref{lem:higher-genus-boundary-qi}), and extension
(Lemma~\ref{lem:extension-across-boundary-qi}).
\emph{B$\to$A}: lifts to H-level by factorization-homology excision
(Remark~\ref{rem:homotopy-native-b}) and model independence.
\emph{Step~C}: PBW spectral sequence collapses at $E_2$
(Lemma~\ref{lem:e2-collapse-higher-genus}), yielding Riordan/Motzkin
generating functions.
\end{remark}

\begin{remark}[Scope boundary at admissible levels]
\label{rem:admissible-scope-boundary}
\index{admissible level!scope boundary}
\index{bar-cobar inversion!failure at admissible levels}
Theorem~B requires Koszulity (MK1). For simple admissible affine
quotients $L_k(\fg)$ at levels $k=-h^\vee+p/q$, that hypothesis is
not currently verified: null vectors obstruct the PBW/Shapovalov
route, and the bar-cobar counit is therefore not promoted here as a
derived equivalence. The obstruction class lives in the higher bar
data encoded in $\Theta_\cA$.
The replacement is the coderived factorization co-contra comparison on
the bar-coalgebra surface
\textup{(}Theorem~\ref{thm:factorization-positselski}\textup{)},
which is available without Koszulness.
\end{remark}

\subsection{Factorization homology descent: the homotopy-native proof}
\label{subsec:fh-descent}

The proof of Theorem~B proceeds by induction on genus at M-level.
This induction is the chain-level shadow of a
single structural statement at H-level: \emph{factorization
homology descent}. This subsection makes the passage from
M-level to H-level explicit and identifies the precise point
where the homotopy-theoretic perspective is necessary.

\begin{proposition}[Pants decomposition as excision;
\ClaimStatusProvedHere]
\label{prop:pants-excision}
\index{factorization homology!excision}
\index{pants decomposition}
Let $\Sigma_g$ be a closed oriented surface of genus~$g \geq 2$,
and let $\Sigma_g = P_1 \cup_{\gamma_1} P_2 \cup_{\gamma_2}
\cdots \cup_{\gamma_{2g-3}} P_{2g-2}$ be a pants decomposition,
where each $P_i$ is a pair of pants (sphere with $3$ boundary
circles) and each $\gamma_j$ is a separating or non-separating
simple closed curve. Then for any $\Eone$-chiral algebra $\cA$ on
$\Sigma_g$, the factorization homology satisfies:
\begin{equation}\label{eq:fh-excision}
\int_{\Sigma_g} \cA
\;\simeq\;
\int_{P_1} \cA
\;\underset{\int_{\gamma_1}\cA}{\otimes^{\mathbb{L}}}\;
\int_{P_2} \cA
\;\underset{\int_{\gamma_2}\cA}{\otimes^{\mathbb{L}}}\;
\cdots
\;\underset{\int_{\gamma_{2g-3}}\cA}{\otimes^{\mathbb{L}}}\;
\int_{P_{2g-2}} \cA\,,
\end{equation}
where each derived tensor product is over the $\Eone$-algebra
$\int_{\gamma_j}\cA$ and the equivalence is in the
$\infty$-category of chain complexes.
\end{proposition}

\begin{proof}
This is Ayala--Francis excision \cite{AF15} applied to the
collar-gluing decomposition $\Sigma_g = P_1 \cup_{\gamma_1 \times
\mathbb{R}} P_2 \cup \cdots$ (each gluing region is a collar
$\gamma_j \times (-\epsilon, \epsilon)$). Since $\cA$ is an
$\Eone$-chiral algebra, hence a factorization algebra on $\Sigma_g$,
\cite[Theorem~3.24]{AF15} gives the equivalence
\[
\textstyle\int_{\Sigma_g}\cA \;\simeq\;
\textstyle\int_{P_1}\cA
\;\underset{\int_{\gamma_1 \times \mathbb{R}}\cA}{\otimes^{\mathbb{L}}}\;
\textstyle\int_{P_2}\cA
\;\underset{}{\otimes^{\mathbb{L}}}\; \cdots
\]
The collar identification
$\int_{\gamma_j \times \mathbb{R}}\cA \simeq \int_{\gamma_j}\cA$
follows from $\mathbb{R}$-invariance of the factorization structure
\cite[Proposition~3.10]{AF15}.
\end{proof}

\begin{proposition}[Genus induction is iterated excision;
\ClaimStatusConditional]
\label{prop:genus-induction-excision}
\index{factorization homology!genus induction}
The genus-$g$ quasi-isomorphism $\psi_g$ of
Theorem~\textup{B} corresponds, under
Proposition~\textup{\ref{prop:pants-excision}}, to iterated
application of the Mayer--Vietoris sequence for factorization
homology:
\begin{enumerate}[label=(\roman*)]
\item \textbf{Base case} ($g = 0$):
$\int_{\mathbb{P}^1}\cA \simeq \Omega(\Bbar^{(0)}(\cA))
\simeq \cA$ is bar-cobar inversion on the Koszul locus
(Theorem~\textup{A}).

\item \textbf{Inductive step} ($g \to g+1$):
Attaching a handle is the collar-gluing
$\Sigma_{g+1} = \Sigma_g \cup_{\gamma \times \mathbb{R}}
(\text{handle})$. The Mayer--Vietoris sequence
\[
\cdots \to
H_n\bigl(\textstyle\int_\gamma\cA\bigr) \to
H_n\bigl(\textstyle\int_{\Sigma_g}\cA\bigr) \oplus
H_n\bigl(\textstyle\int_{\mathrm{handle}}\cA\bigr) \to
H_n\bigl(\textstyle\int_{\Sigma_{g+1}}\cA\bigr)
\to \cdots
\]
is the H-level content of the extension lemma
(Lemma~\textup{\ref{lem:extension-across-boundary-qi}}).

\item \textbf{Handle contribution}: The factorization homology
of a handle (annulus with two boundary circles identified) is
controlled by the Hochschild homology of the $\Eone$-algebra
$\int_\gamma\cA$, by the Hochschild--Kostant--Rosenberg
theorem for factorization homology
\cite[Theorem~5.5.3.10]{HA}.
The genus-$g$ curvature
$\mcurv \cdot \omega_g$ arises as the Euler class of the handle's
tangent bundle, pulled back to the Hochschild complex
via the Chern character.
\end{enumerate}
\end{proposition}

\begin{proof}
The correspondence between the chain-level proof of Theorem~B and
factorization homology descent proceeds step by step.

For (i), the genus-$0$ factorization homology $\int_{\mathbb{P}^1}\cA$
is computed by the Beilinson--Drinfeld chiral homology
\cite[Theorem~4.3.1]{BD04}, which identifies with the cobar
construction $\Omega(\Bbar^{(0)}(\cA))$: the factorization structure
maps correspond to the cobar differential $d_{\mathrm{cobar}}$,
and factorization along collisions gives the bar differential
$d_{\mathrm{bar}}$.

For (ii), the open stratum quasi-isomorphism
(Lemma~\ref{lem:higher-genus-open-stratum-qi}) is the restriction to
$\Sigma_g \smallsetminus \gamma \hookrightarrow \Sigma_{g+1}$,
and the boundary quasi-isomorphism
(Lemma~\ref{lem:higher-genus-boundary-qi}) is the restriction to
a tubular neighborhood of $\gamma$. The extension lemma assembles
these via the long exact Mayer--Vietoris sequence; the induction
hypothesis guarantees that $\psi_g$ is already a quasi-isomorphism.

For (iii), the Hochschild complex of $\int_\gamma\cA$ computes
the factorization homology of the annulus $\gamma \times [0,1]$
by \cite[Theorem~5.5.3.10]{HA}. The identification of the handle
(where the two boundary circles are identified) introduces a trace
map $\mathrm{tr} \colon HH_*(\int_\gamma\cA) \to
\int_{\text{handle}}\cA$. The curvature arises because this trace
map, when lifted to chains, satisfies $d_{\mathrm{tr}}^2 =
\mcurv \cdot \lambda_1$ (the Hodge class), which is the content of
Theorem~\ref{thm:genus-universality} viewed from the factorization
homology perspective.
\end{proof}

\begin{remark}[Why the homotopy perspective is necessary]
\label{rem:why-homotopy-necessary}
\index{factorization homology!necessity of homotopy perspective}
The H-level statement is strictly stronger than the M-level
quasi-isomorphism in three ways:
(i)~excision~\eqref{eq:fh-excision} is an $\Eone$-algebra
equivalence, producing $\Ainf$-structure automatically;
(ii)~Mayer--Vietoris is functorial in the pants decomposition;
(iii)~the coderived category
$\cat{D}^{\mathrm{co}}_\infty(\cat{Fact}_g(X))$ governing
the off-Koszul regime has no chain-level analogue.
\end{remark}

\begin{proposition}[\texorpdfstring{$E_2$}{E2}-collapse as formality;
\ClaimStatusConditional]
\label{prop:e2-collapse-formality}
\index{formality!spectral sequence collapse}
On the Koszul locus, the following conditions on a chiral algebra
$\cA$ are equivalent:
\begin{enumerate}[label=(\roman*)]
\item The bar spectral sequence
$E_1^{p,q}(g) \Rightarrow H^{p+q}(\Omega_g\Bbar_g(\cA))$
collapses at~$E_2$ for all $g \geq 0$.
\item The filtered $\Ainf$-algebra $\Omega_g(\Bbar_g(\cA))$ is
\emph{formal}: quasi-isomorphic to its cohomology with the
transferred $\Ainf$-structure
(Theorem~\textup{\ref{thm:chiral-htt}}).
\item The transferred higher operations satisfy
$\tilde{m}_n^{(g)} = 0$ for $n \geq 3$ and all~$g$
(Koszulness at all genera).
\end{enumerate}
In the factorization homology picture, formality means that
$\int_{\Sigma_g}\cA$ is determined by $\int_{\mathbb{P}^1}\cA$
and the handle attachments, with no higher coherence obstructions.
\end{proposition}

\begin{proof}
(i) $\Rightarrow$ (ii): The spectral sequence arises from the bar
filtration on $\Omega_g(\Bbar_g(\cA))$. Collapse at~$E_2$ means
that $E_2 \cong E_\infty \cong H^*(\Omega_g(\Bbar_g(\cA)))$,
and the filtration on cohomology is split. This splitting produces
the SDR $(\Omega_g(\Bbar_g(\cA)),\, H^*,\, p,\, \iota,\, h)$
of Theorem~\ref{thm:bar-cobar-htt}, and the transferred structure
on $H^*$ has $\tilde{m}_n = 0$ for $n \geq 3$ by the argument of
Step~3 in that proof (Koszul-line concentration forces vanishing).
Hence $\Omega_g(\Bbar_g(\cA))$ is formal.

(ii) $\Rightarrow$ (iii): If $\Omega_g(\Bbar_g(\cA))$ is formal,
the transferred operations $\tilde{m}_n$ are the operations on the
formal model, which has $\tilde{m}_n = 0$ for $n \geq 3$ by
definition.

(iii) $\Rightarrow$ (i): If $\tilde{m}_n = 0$ for $n \geq 3$,
the minimal model of $\Omega_g(\Bbar_g(\cA))$ has only binary
operations. The spectral sequence of the bar filtration on a
formal $\Ainf$-algebra necessarily collapses at~$E_2$:
the $E_2$-differentials are the only non-trivial ones
(they encode $m_2$), and all higher differentials vanish because
$\tilde{m}_n = 0$ for $n \geq 3$
(each $E_r$-differential for $r \geq 3$ is controlled by
$\tilde{m}_r$; cf.\ \cite[Theorem~3.4.1]{LV12}).
\end{proof}

\begin{remark}[The factorization algebra as the natural object]
\label{rem:factorization-natural}
\index{factorization algebra!naturality}
The factorization algebra $\int_{\Sigma_g}\cA$ is the H-level object
of which $\Omega_g(\Bbar_g(\cA))$ is a presentation (M-level) and
the Hilbert series is a shadow (S-level).
Different presentations (SDR choices, pants decompositions) produce
canonically equivalent H-level data.
The passage M$\to$H requires homotopy invariance of the bar-cobar
adjunction, addressed in Chapter~\ref{ch:en-koszul-duality}.
\end{remark}

%================================================================
% KOSZUL DUALITY ACROSS GENERA
% (Absorbed from koszul_across_genera.tex)
%================================================================


%================================================================
% AMBIENT COMPLEMENTARITY AND THE LAGRANGIAN UPGRADE
%================================================================

\section{Ambient complementarity and the Lagrangian upgrade}
\label{sec:ambient-complementarity-lagrangian}

The preceding sections prove that complementarity decomposes the
genus-$g$ cohomology into two dual packages:
$Q_g(\cA) \oplus Q_g(\cA^!) \cong H^*(\overline{\mathcal{M}}_g, Z(\cA))$.
This is the \emph{linear shadow} of a deeper geometric fact. The
correct theorem is not an additive splitting of vector spaces. The
correct theorem is that there is \emph{one} ambient deformation problem,
equipped with a $(-1)$-shifted symplectic structure, and the two dual
packages are realized inside it as \emph{complementary Lagrangians}.
Once that geometry is installed, the scalar characteristic
$\kappa(\cA)$ and the spectral package $\Delta_{\cA}$ become the
quadratic face of a larger nonlinear theory whose higher jets are the
cubic, quartic, and all-degree shadow tensors computed in
\S\ref{sec:heisenberg-shadow-gaussianity},
\S\ref{sec:affine-cubic-shadow},
\S\ref{sec:betagamma-quartic-birth}, and
\S\ref{sec:mixed-cubic-quartic-shadows}.

\begin{remark}[Why additive complementarity is not enough:
the clutching obstruction]
\label{rem:why-additive-not-enough}
Consider a naive additive splitting
$H^*(\overline{\mathcal{M}}_g, Z(\cA))
= V_{\cA} \oplus V_{\cA^!}$
of vector spaces. Such a splitting is \emph{unstable under
genus gluing}: if $(x,y) \in V_{\cA} \times V_{\cA^!}$ at
genera $g_1, g_2$, the separating clutching map $\xi_*$ on
$\Sigma_{g_1} \cup_{\mathrm{node}} \Sigma_{g_2}$ does
\emph{not} respect the direct-sum decomposition unless the
deformation spaces are related by a Lagrangian condition:
\[
\xi_*^{-1}(V_{\cA} \oplus V_{\cA^!})
\;\neq\;
V_{\cA} \oplus V_{\cA^!}
\qquad\text{generically.}
\]
The ambient theorem repairs this by imposing a single deformation
problem $(L_{\cA} \oplus K_{\cA} \oplus L_{\cA^!})$ with a
$(-1)$-shifted symplectic form and isotropicity of each side.
Then $\xi_*$ respects the Lagrangian splitting because
isotropicity is \emph{equivariant} under the modular operad
action, whereas arbitrary direct sums are not.
\end{remark}

\subsection{The ambient complementarity datum}
\label{subsec:ambient-complementarity-datum}

\begin{definition}[Ambient complementarity datum]
\label{def:ambient-complementarity-datum}
Let $\cA$ be a modular Koszul chiral algebra on a smooth projective
curve, and let $\cA^!$ denote its Koszul dual on the locus where the
duality package is defined. An \emph{ambient complementarity datum}
consists of:
\begin{enumerate}[label=\textup{(\roman*)}]
\item two cyclic deformation complexes
\[
L_{\cA} := \operatorname{Def}_{\mathrm{cyc}}(\cA),
\qquad
L_{\cA^!} := \operatorname{Def}_{\mathrm{cyc}}(\cA^!),
\]
viewed as filtered complete cyclic $L_\infty$-algebras;
\item a completed kernel object
\[
K_{\cA} := (\cA^! \widehat\otimes \cA)[1]
\]
with differential twisted by the universal twisting kernel;
\item a Maurer--Cartan element
\[
\tau_{\cA} \in \operatorname{MC}(\cA^! \widehat\otimes \cA)
\]
representing the canonical bar--cobar twisting kernel;
\item linearizations of the kernel in the two deformation directions,
\[
\nabla_{\cA} \colon L_{\cA} \to K_{\cA},
\qquad
\nabla_{\cA^!} \colon L_{\cA^!} \to K_{\cA};
\]
\item a degree $-1$ cyclic pairing on the total complex
$L_{\cA} \oplus K_{\cA} \oplus L_{\cA^!}$
for which the kernel differential is skew-adjoint.
\end{enumerate}
\end{definition}

\begin{definition}[Ambient complementarity tangent complex]
\label{def:ambient-complementarity-tangent-complex}
Given an ambient complementarity datum, the
\emph{ambient complementarity tangent complex} is
\begin{equation}
\label{eq:ambient-tangent-complex}
T_{\mathrm{comp}}(\cA)
:=
\operatorname{fib}\Bigl(
\nabla_{\cA} - \nabla_{\cA^!} \colon
L_{\cA} \oplus L_{\cA^!} \longrightarrow K_{\cA}
\Bigr).
\end{equation}
A tangent vector in $T_{\mathrm{comp}}(\cA)$ is a first-order
deformation of $\cA$, a first-order deformation of $\cA^!$, and a
first-order deformation of the universal kernel, constrained by the
linearized Maurer--Cartan equation.
\end{definition}

\begin{theorem}[Ambient complementarity in tangent form]
\label{thm:ambient-complementarity-tangent}
\ClaimStatusProvedHere
Assume an ambient complementarity datum exists for $\cA$. Then:
\begin{enumerate}[label=\textup{(\roman*)}]
\item the complex $T_{\mathrm{comp}}(\cA)$ carries a canonical
nondegenerate pairing of degree $-1$;
\item the one-sided tangent complexes
\[
T_{\cA} := \operatorname{fib}(\nabla_{\cA} \colon L_{\cA} \to K_{\cA}),
\qquad
T_{\cA^!} := \operatorname{fib}(\nabla_{\cA^!} \colon L_{\cA^!} \to K_{\cA})
\]
are isotropic inside $T_{\mathrm{comp}}(\cA)$;
\item if the pairing on $L_{\cA} \oplus K_{\cA} \oplus L_{\cA^!}$
is perfect and bar--cobar duality identifies the normal complex of
one side with the shifted dual tangent complex of the other, then
$T_{\cA}$ and $T_{\cA^!}$ are complementary Lagrangians in
$T_{\mathrm{comp}}(\cA)$.
\end{enumerate}
\end{theorem}

\begin{proof}
The degree $-1$ pairing on $T_{\mathrm{comp}}(\cA)$ is induced from the
cyclic pairing on the direct sum by passage to the homotopy fiber.
Because the total differential is skew-adjoint, the pairing descends to
cohomology and is compatible with the differential.

For the one-sided tangent complex $T_{\cA}$, the pullback of the ambient
pairing reduces to the quadratic term in the differentiated
Maurer--Cartan equation for the universal kernel. That quadratic term
vanishes because the Maurer--Cartan equation is exactly the isotropy
condition for the graph of the one-sided variation. The same argument
applies to $T_{\cA^!}$.

If the cyclic pairing is perfect and bar--cobar duality identifies the
normal complex to one side with the shifted dual tangent complex of the
other, maximal isotropicity follows. This is the derived form of the
principle that the two sides are opposite polarizations of a single
symplectic deformation problem.
\end{proof}

\begin{theorem}[Ambient complementarity as shifted symplectic formal
moduli problem]
\label{thm:ambient-complementarity-fmp}
\ClaimStatusConditional
Under the perfectness and nondegeneracy hypotheses of
Theorem~\textup{\ref{thm:ambient-complementarity-tangent}(iii)},
assume that $T_{\mathrm{comp}}(\cA)$ is represented by a complete
filtered dg~Lie algebra or cyclic $L_\infty$-algebra with a closed
invariant non-degenerate pairing of degree~$-1$, and that the
one-sided complexes are filtered $L_\infty$ subalgebras with
Lagrangian tangent maps. Then $T_{\mathrm{comp}}(\cA)$ integrates to a
formal moduli problem $\mathcal{M}_{\mathrm{comp}}(\cA)$ carrying
a canonical $(-1)$-shifted symplectic structure, and the one-sided
deformation problems define Lagrangian maps
\[
\mathcal{M}_{\cA}
\longrightarrow
\mathcal{M}_{\mathrm{comp}}(\cA)
\longleftarrow
\mathcal{M}_{\cA^!}.
\]
The genus-$g$ complementarity complex is the linear shadow of this
picture.
\end{theorem}

\begin{proof}
By hypothesis, the modular cyclic deformation complex
$\Defcyc^{\mathrm{mod}}(\cA)$
(Definition~\ref{def:modular-cyclic-deformation-complex})
is a complete filtered dg~Lie algebra: the product topology over
$(g,n)$ with the genus filtration is complete, and the cyclic
structure from the non-degenerate invariant form on~$\cA$ makes it
a cyclic dg~Lie algebra (hence a strict cyclic $L_\infty$-algebra).

By the Lurie--Pridham correspondence
\cite{LurieDAGX, Pridham17},
every complete filtered dg~Lie algebra $\fg$ integrates to a
formal moduli problem $\mathcal{M}_\fg$
whose tangent complex is $\fg[1]$.
Applied to $\fg = \Defcyc^{\mathrm{mod}}(\cA)$, this gives
the formal moduli problem $\mathcal{M}_{\mathrm{comp}}(\cA)$.

The cyclic pairing on $\Defcyc^{\mathrm{mod}}(\cA)$ is a
non-degenerate invariant pairing of degree~$-1$. By the
PTVV/Kontsevich--Pridham principle \cite{PTVV13,Pridham17}, an
invariant non-degenerate pairing of degree~$n$ on a dg~Lie
algebra $\fg$ induces an $n$-shifted symplectic structure
on $\mathcal{M}_\fg$. The degree~$-1$ pairing on
$\Defcyc^{\mathrm{mod}}(\cA)$ thus gives a
$(-1)$-shifted symplectic structure on
$\mathcal{M}_{\mathrm{comp}}(\cA)$.

For the Lagrangian structure: under the filtered subalgebra and
perfectness hypotheses,
the one-sided tangent complexes $T_\cA$ and $T_{\cA^!}$ are
complementary Lagrangians in $T_{\mathrm{comp}}(\cA)$
(Theorem~\ref{thm:ambient-complementarity-tangent}(iii)).
By the PTVV Lagrangian correspondence, isotropic sub-dg-Lie
algebras whose tangent complexes are Lagrangian
integrate to Lagrangian maps of formal moduli problems.
\end{proof}

\begin{remark}[PTVV scope]
The comparison with \cite{PTVV13} has two distinct layers.
At the ambient level,
Theorem~\ref{thm:ambient-complementarity-fmp}
uses the degree-$(-1)$ cyclic pairing on
$T_{\mathrm{comp}}(\cA)$ to produce a
$(-1)$-shifted symplectic formal moduli problem.
At fixed genus,
Proposition~\ref{prop:ptvv-lagrangian}
uses the linear derived stack attached to
$C_g = R\Gamma(\overline{\mathcal{M}}_g,\mathcal{Z}(\cA))$:
the Verdier pairing gives the constant
$(-(3g{-}3))$-shifted symplectic form of
\cite[Example~1.4]{PTVV13}, and the perfect subcomplexes whose
cohomology groups are denoted
$Q_g(\cA)$ and $Q_g(\cA^!)$
satisfy the Lagrangian conditions of
\cite[Definition~2.8, Theorem~2.9]{PTVV13}.

This matches the PTVV framework, but it does not identify
complementarity with a derived-intersection computation, and it does
not derive the scalar conductor identity
$\kappa(\cA)+\kappa(\cA^!)=K$ from shifted symplectic geometry alone.
The conductor formula still enters only after the Verdier
eigenspace decomposition is projected to the scalar lane and the
family-specific computations of Theorem~D are used.
\end{remark}

\begin{theorem}[Global Lagrangian complementarity across genera:
clutching-coherent section of the Verdier-symplectic bundle on
$\overline{\mathcal{M}}_{g,n}$; clauses~\textup{(i)--(iii)}~\ClaimStatusConditional;
clause~\textup{(iv)}~\ClaimStatusConditional]
\label{thm:lagrangian-complementarity-global-upgrade}
\phantomsection\label{V1-thm:lagrangian-complementarity-global-upgrade}%
\index{complementarity!global Lagrangian}%
\index{shifted symplectic!Verdier, $-(3g-3)$}%
Let $\cA$ be a modular Koszul chiral algebra on a smooth projective
curve satisfying the perfectness hypotheses of
Theorem~\textup{\ref{thm:ambient-complementarity-tangent}(iii)}, and let
$\cA^!$ denote its Koszul dual on the locus where the duality package
is defined. Restrict to the uniform-weight scalar lane
(multi-weight cases require the cross-channel $\delta F_g^{\mathrm{cross}}$
correction of Theorem~D). Then:
\begin{enumerate}[label=\textup{(\roman*)}]
\item \emph{Verdier-symplectic bundle.}
For each $(g,n)$ with $2g-2+n>0$, the Verdier self-duality of
$\mathcal{Z}(\cA)$ on $\overline{\mathcal{M}}_{g,n}$ equips the sheaf
\[
 \mathbf{T}^{\mathrm{mod}}_{g,n}
 := R\Gamma\bigl(\overline{\mathcal{M}}_{g,n},\, \mathcal{Z}(\cA) \oplus
 \mathcal{Z}(\cA^!)\bigr)
\]
with a canonical $(-(3g-3+n))$-shifted nondegenerate skew pairing,
extending \textup{\cite[Example~1.4]{PTVV13}} to the boundary of
$\overline{\mathcal{M}}_{g,n}$. This is the \emph{Verdier-symplectic
structure on the modular tangent bundle} and is distinct from the
$(-1)$-shifted ambient structure of
Theorem~\textup{\ref{thm:ambient-complementarity-fmp}}.
\item \emph{Genus-wise Lagrangian section.}
The Koszul pair $(\cA,\cA^!)$ defines, at each $(g,n)$, a pair of
perfect subcomplexes $\mathbf{L}_{g,n}(\cA)$ and
$\mathbf{L}_{g,n}(\cA^!)$ of
$\mathbf{T}^{\mathrm{mod}}_{g,n}$ that are complementary Lagrangians
for the Verdier pairing in~\textup{(i)}, in the sense of
Proposition~\textup{\ref{prop:ptvv-lagrangian}} applied fiberwise
over boundary strata.
\item \emph{Clutching coherence.}
For every separating clutching map
\[
 \xi: \overline{\mathcal{M}}_{g_1,n_1+1} \times
 \overline{\mathcal{M}}_{g_2,n_2+1}
 \longrightarrow \overline{\mathcal{M}}_{g_1+g_2,\, n_1+n_2},
\]
the section commutes with $\xi$ up to $R$-matrix-twisted K\"unneth:
\[
 \xi^*\mathbf{L}_{g_1+g_2}(\cA)
 \;\cong\; \bigl(\mathbf{L}_{g_1}(\cA) \boxtimes \mathbf{L}_{g_2}(\cA)\bigr)^{R\text{-twist}},
\]
with $R(z)$ the spectral $R$-matrix of Theorem~H. An analogous
identity holds for non-separating clutching via the annular bar
model $B^{\mathrm{ann}}$.
\item \emph{Global assembly \textup{(}\ClaimStatusConditional\textup{)}.}
The assembled data $(\mathbf{L}_{\bullet}(\cA),\mathbf{L}_{\bullet}(\cA^!))$
is a \emph{global Lagrangian section} of the Verdier-symplectic bundle
on the moduli operad $\overline{\mathcal{M}}_{\bullet,\bullet}$.
Conditional on the modular-bootstrap-to-curved-Dunn bridge
$\Phi$ in cohomological degrees $\leq 2$, the genus-wise
Lagrangian sections cohere to a single global Lagrangian subobject,
i.e.\ the assignment $(g,n)\mapsto\mathbf{L}_{g,n}$ factors through
a Lagrangian correspondence of modular operads.
\emph{Conditional.}~Clause~\textup{(iv)} Global assembly depends on
the Vol.~II modular-bootstrap-to-curved-Dunn comparison theorem
Proposition~\ref{V2-prop:modular-bootstrap-to-curved-dunn-bridge}
\textup{(}\texttt{curved\_dunn\_higher\_genus.tex:138}\textup{)};
its proof is \emph{not} inscribed in Vol~I.
\end{enumerate}
\end{theorem}

\begin{proof}
\textbf{(i)} Verdier duality on each $\overline{\mathcal{M}}_{g,n}$ gives
$\mathbb{D}\mathcal{Z}(\cA)\cong \mathcal{Z}(\cA)[2(3g-3+n)]$ under the
cyclic perfectness hypothesis (invariant form on $\cA$). The Serre pairing
$H^k\otimes H^{2(3g-3+n)-k}\to k$ shifts to degree $-(3g-3+n)$ after the
canonical shift, exactly PTVV Example~1.4 extended to the DM boundary
via Saito's log-perverse extension. Nondegeneracy on boundary strata
follows from strata-wise applicability of
Theorem~\textup{\ref{thm:quantum-complementarity-main}}.

\textbf{(ii)} At smooth $(g,n)$, the Lagrangian condition is
Proposition~\textup{\ref{prop:ptvv-lagrangian}}. At boundary strata, the
clutching formula in~(iii) plus fiberwise Verdier duality preserves
isotropy; closure under the bar differential gives the subcomplex.

\textbf{(iii)} Separating clutching factorizes the Verdier pairing as
the Serre pairing on each factor plus the residual pairing at the node.
The $R$-matrix enters through the chiral Hochschild pairing at the
node (annular bar $B^{\mathrm{ann}}$ with monodromy $R(z)$), exactly the
$R$-twisted K\"unneth of the Theorem~A$^{\infty,2}$ descent. Non-separating
clutching reduces to the $S^1$-sewing of $B^{\mathrm{ann}}$ as in the
higher-genus foundations theorem~\textup{\ref{V2-thm:modular-bar}}.

\textbf{(iv)} Genus-wise Lagrangians assemble if and only if the
obstruction cocycle in the curved-Dunn twisting-cochain
complex is exact. Under $\Phi$, the modular-bootstrap
vanishing $H^2_{\mathrm{boot}}(\overline{\mathcal{M}}_{g,n})=0$
\textup{(}Proposition~\textup{\ref{prop:modular-bootstrap-H2-vanishing}}%
\textup{)} maps to the curved-Dunn obstruction class; its image
vanishes in cohomological degree~$\leq 2$, which suffices for
assembly of a Lagrangian section (isotropicity + closure are degree~$\leq 2$
conditions). The bridge map $\Phi$ itself, together with its
cohomological iso on $H^2$ needed to transport vanishing, is
the Vol.~II comparison theorem
Proposition~\ref{V2-prop:modular-bootstrap-to-curved-dunn-bridge};
we treat it as an imported input here, which is why clause~(iv)
carries the \ClaimStatusConditional\ tag in contrast to the
self-contained clauses (i)--(iii). Multi-weight cases require the
cross-channel $\delta F_g^{\mathrm{cross}}$ correction of Theorem~D to
enter the isotropy check.
\end{proof}

\begin{remark}[Verdier and Calabi--Yau shifted structures]
\label{rem:fm116-scope-global-upgrade}
Theorem~\textup{\ref{thm:lagrangian-complementarity-global-upgrade}}
is the $(-(3g-3+n))$-shifted \emph{Verdier-symplectic} upgrade;
it does \emph{not} claim PTVV-for-$\mathrm{CY}_d$ on the moduli stack.
The distinction is mathematical: Verdier-on-$\overline{\mathcal{M}}_{g,n}$
is an intrinsic pairing built from Serre duality and Verdier self-dual
$\mathcal{Z}(\cA)$; PTVV-for-$\mathrm{CY}_d$ is the external pairing built
from a Calabi--Yau target and lives on mapping stacks of the target.
These two structures are related only through the chiral-to-CY
functor of Vol~III, which is \emph{not} invoked here.
\end{remark}

\begin{remark}[Global form of Theorem~C]
\label{rem:platonic-upgrade-lagrangian-global}
Theorem~\textup{\ref{thm:quantum-complementarity-main}} (Theorem~C)
delivers a genus-wise Lagrangian decomposition; Theorem~\textup{\ref{thm:ambient-complementarity-fmp}}
delivers a $(-1)$-shifted ambient picture. The global form combines
both: the $(-(3g-3+n))$-shifted Verdier-symplectic modular tangent
bundle admits a \emph{single} global Lagrangian section whose
specializations at fixed genus recover Theorem~C and whose shadow
at the cyclic deformation side recovers Theorem~\textup{\ref{thm:ambient-complementarity-fmp}}.
The bridge between the two shifted structures is the
PTVV-Kontsevich--Pridham cyclic-to-Verdier principle: cyclicity of the
pairing on $\Defcyc^{\mathrm{mod}}(\cA)$ pulls back from the Verdier
pairing on $\mathcal{Z}(\cA)$ under the Chern--Weil trace of
\S\ref{sec:ambient-complementarity-lagrangian}.
\end{remark}

\subsection{The complementarity potential: direct definition}
\label{subsec:cotangent-normal-form}

The complementarity potential is the generating function of one
Lagrangian in the cotangent chart centered at the other. We define
it directly from the cyclic deformation data and then prove its
existence via the cotangent normal form.

\begin{definition}[Complementarity potential: direct construction]
\label{def:complementarity-potential}
Let $V_{\cA} \subset H^*(\operatorname{Def}_{\mathrm{cyc}}(\cA))$
be a finite-dimensional cyclic slice with nondegenerate Hessian
$H_{\cA}$ and propagator $P_{\cA} = H_{\cA}^{-1}$. The
\emph{complementarity potential} is the formal function
\begin{equation}
\label{eq:complementarity-potential-direct}
S_{\cA}(x)
\;:=\;
\sum_{r \geq 2} \frac{1}{r!}\,
\operatorname{Sh}_r(\cA)(x, \ldots, x),
\end{equation}
where $\operatorname{Sh}_r(\cA)(x_1,\ldots,x_r)
:= \sum_{\mathrm{cyc}} \langle x_1,\,
\ell_{r-1}^{\mathrm{tr}}(x_2,\ldots,x_r)\rangle$
is the degree-$r$ shadow jet, and $\ell_n^{\mathrm{tr}}$ are the
transferred cyclic $\Ainf$-operations on the minimal model.
The first three terms are the \emph{shadow jets}:
\begin{equation}
\label{eq:shadow-jets-direct}
S_{\cA}(x)
\;=\;
\underbrace{\tfrac12 H_{\cA}(x,x)}_{\text{Hessian}}
\;+\;
\underbrace{\tfrac{1}{6}\mathfrak{C}_{\cA}(x,x,x)}_{\text{cubic shadow}}
\;+\;
\underbrace{\tfrac{1}{24}\mathfrak{Q}_{\cA}(x,x,x,x)}_{\text{quartic shadow}}
\;+\; O(x^5).
\end{equation}
The dual Lagrangian $\mathcal{M}_{\cA^!}$ is the graph of $dS_{\cA}$
in the cotangent chart centered at $\mathcal{M}_{\cA}$.
\end{definition}

\begin{remark}[What the potential computes]
At genus~$0$: $S_{\cA}$ is the classical complementarity action.
Its critical locus $\operatorname{Crit}(S_{\cA})$ governs
simultaneous self-dual deformations of $\cA$ and $\cA^!$.
When $S_{\cA}$ is exactly quadratic (i.e., only $H_{\cA}$
survives), complementarity is ``fake,'' a linear splitting.
The first nonzero higher jet of $S_{\cA}$ is the first genuine
nonlinear modular invariant. For the five standard families:
\begin{center}
\begin{tabular}{@{}lll@{}}
\toprule
Family & Leading nonlinear jet & Potential type \\
\midrule
Heisenberg & none ($S = \frac12 H$) & exactly quadratic \\
Affine $\widehat{\mathfrak{sl}}_2$ &
$\frac16\kappa(a,[a,a])$ & polynomial (cubic) \\
$\beta\gamma$ & $\frac{1}{24}\operatorname{cyc}(m_3)$ &
polynomial (quartic) \\
Virasoro & $\frac16(2x^3) + \frac{1}{24}Q_0 x^4$ &
\textbf{non-polynomial} \\
$\mathcal{W}_N$ ($N \geq 3$) &
$\frac16\mathfrak{C}^{\mathrm{grav}} + \frac{1}{24}
\mathfrak{Q}^{\mathrm{contact}}$ &
\textbf{non-polynomial} \\
\bottomrule
\end{tabular}
\end{center}
\end{remark}

Each entry in the table can now be made explicit.

\begin{example}[Complementarity potential: Heisenberg]
\label{ex:complementarity-potential-heisenberg}
For the rank-$1$ Heisenberg $\mathcal{H}_\kappa$ with
$\operatorname{OPE}:J(z)J(w) \sim \kappa/(z{-}w)^2$, the cyclic
slice is $V_{\mathcal{H}} = \mathbb{C}\langle J\rangle$ with
Hessian $H_{\mathcal{H}} = \kappa\,x^2$ (the double-pole coefficient).
The OPE has \emph{no simple pole}: $\ell_2^{\mathrm{tr}}(J,J)=0$.
Therefore all higher shadows vanish, and
\[
S_{\mathcal{H}}(x) = \tfrac12\kappa\,x^2.
\]
The dual Lagrangian is $\mathcal{M}_{\mathcal{H}^!}
= \operatorname{graph}(dS)
= \operatorname{graph}(\kappa\,x\,dx)$, which is \emph{linear}.
Complementarity is fake
(Proposition~\ref{prop:fake-complementarity-criterion}).
\end{example}

\begin{example}[Complementarity potential: affine
$\widehat{\mathfrak{sl}}_2$]
\label{ex:complementarity-potential-affine}
For $\widehat{\mathfrak{sl}}_{2,k}$ with
$J^a(z)J^b(w) \sim (k{+}2)\kappa^{ab}/(z{-}w)^2
+ f^{ab}{}_c J^c(w)/(z{-}w)$,
the double pole gives $H_{\mathrm{aff}} = (k{+}2)\kappa$
and the simple pole gives the Lie bracket
$\ell_2^{\mathrm{tr}}(y,z) = [y,z]$.
The quartic obstruction vanishes by the Jacobi identity.
The complementarity potential is:
\begin{align}
S_{\mathrm{aff}}(a)
&= \tfrac12(k{+}2)\,\kappa(a,a)
+ \tfrac16\kappa(a,[a,a]).
\label{eq:complementarity-potential-affine}
\end{align}
At the \emph{critical level} $k = -2$ (where $h^\vee = 2$ for
$\mathfrak{sl}_2$), the Hessian vanishes and the potential
becomes purely cubic:
$S_{\mathrm{crit}}(a) = \tfrac16\kappa(a,[a,a])$.
The critical locus
$\operatorname{Crit}(S_{\mathrm{crit}}) = \{a : [a,a]=0\}$
is the Maurer--Cartan locus for $\mathfrak{sl}_2$. This is related to
the critical Feigin--Frenkel/oper story, but the manuscript does not
identify this critical locus itself with the oper space; the proved
bar-cohomological statement is Theorem~\ref{thm:oper-bar}.
\end{example}

\begin{example}[Complementarity potential: $\beta\gamma$]
\label{ex:complementarity-potential-betagamma}
On the weight-changing line $V_{\beta\gamma} = \mathbb{C}\langle\eta
\rangle$ with $\langle\eta,\eta\rangle = c'(\lambda)$
(the derivative of the central charge $c(\lambda)
= 1 - 12\lambda(1{-}\lambda) + 12\lambda^2$):
the transferred operations $\ell_n^{\mathrm{tr}}(\eta^n) = 0$ for
all $n \geq 2$ by rank-one abelian rigidity
(Theorem~\ref{thm:betagamma-rank-one-rigidity}). Therefore
\[
S_{\beta\gamma}(x) = \tfrac12 c'(\lambda)\,x^2.
\]
The potential is exactly quadratic on the weight line. Off this
line (on the full weight/contact slice), $\mathfrak{Q}_{\beta\gamma}$
contributes a quartic term from $m_3$.
\end{example}

\begin{example}[Complementarity potential: Virasoro]
\label{ex:complementarity-potential-virasoro}
For $\operatorname{Vir}_c$ on the primary line
$V_{\mathrm{Vir}} = \mathbb{C}\langle T\rangle$ with
$T(z)T(w) \sim (c/2)/(z{-}w)^4 + 2T(w)/(z{-}w)^2 + \partial
T(w)/(z{-}w)$:
\begin{align}
H_{\mathrm{Vir}} &= \tfrac{c}{2}\,x^2,
&
\mathfrak{C}_{\mathrm{Vir}} &= 2x^3,
&
\mathfrak{Q}^{\mathrm{contact}}_{\mathrm{Vir}}
&= \frac{10}{c(5c{+}22)}\,x^4.
\end{align}
The complementarity potential through quartic order is:
\begin{equation}
\label{eq:complementarity-potential-virasoro}
S_{\mathrm{Vir}}(x)
= \frac{c}{4}\,x^2
+ \frac{1}{3}\,x^3
+ \frac{10}{24\,c(5c{+}22)}\,x^4
+ \operatorname{Sh}_5 x^5 + \cdots
\end{equation}
The quintic $\operatorname{Sh}_5 \neq 0$ generically
(Theorem~\ref{thm:w-virasoro-quintic-forced}): the potential is
\textbf{non-polynomial} and the shadow obstruction tower is infinite. This
is the mechanism by which the Virasoro modular package carries
strictly more information than the $\hat{A}$-genus.
\end{example}

\begin{example}[Legendre duality: affine case worked]
\label{ex:legendre-duality-affine}
The Legendre dual of the affine potential
\eqref{eq:complementarity-potential-affine} is computed by
extremizing $\langle p, a\rangle - S_{\mathrm{aff}}(a)$:
\[
p = \frac{\partial S}{\partial a}
= (k{+}2)\kappa(a,\cdot) + \tfrac12\kappa([a,a],\cdot),
\qquad
a = \frac{p}{(k{+}2)\kappa} + O(p^2).
\]
Substituting into the Legendre formula and collecting through
cubic order:
\[
S_{\mathrm{aff}}^\vee(p)
= \frac{1}{2(k{+}2)}\kappa^{-1}(p,p)
- \frac{1}{6(k{+}2)^3}\kappa^{-1}(p,[\kappa^{-1}p,\kappa^{-1}p])
+ O(p^4).
\]
Since $S_{\mathrm{aff}}^\vee = S_{\widehat{\mathfrak{sl}}_2^!}$
by Proposition~\ref{prop:legendre-duality-potentials}, the cubic
shadow of the Koszul dual is:
\[
\mathfrak{C}_{\mathrm{aff}^!}(p,p,p)
= -\frac{1}{(k{+}2)^3}\kappa([\kappa^{-1}p,\kappa^{-1}p],
\kappa^{-1}p).
\]
As stated in Proposition~\ref{prop:legendre-duality-cubic}, the
cubic on the dual side is the Hessian-transported negative of the
cubic on the primal side.
\end{example}

The complementarity potential arises from shifted
symplectic geometry.

\begin{theorem}[Shifted cotangent normal form]
\label{thm:shifted-cotangent-normal-form}
\ClaimStatusConditional
Let $\mathcal{M}$ be a pointed $(-1)$-shifted symplectic formal moduli
problem, and let $i_+\colon\mathcal{L}_+\to\mathcal{M}$,
$i_-\colon\mathcal{L}_-\to\mathcal{M}$ be pointed Lagrangian maps.
Assume the induced tangent map
$T_0\mathcal{L}_+ \oplus T_0\mathcal{L}_-
\xrightarrow{\sim} T_0\mathcal{M}$
is a quasi-isomorphism. Assume that the shifted symplectic
neighborhood theorem applies to $i_+$ in the pointed formal category
under consideration, and that the formal Poincare lemma holds for
degree-zero one-forms on $\mathcal{L}_+$. Then there exists a pointed
formal
symplectomorphism
\[
(\mathcal{M}, i_+)
\simeq
(T^*[-1]\mathcal{L}_+, 0\text{-section})
\]
under which $\mathcal{L}_-$ becomes the graph of a closed
degree-zero one-form $\alpha_+ \in \Omega^1(\mathcal{L}_+)^0$.
In the pointed formal neighborhood this one-form is exact:
$\alpha_+ = dS_+$
for a degree-zero formal function $S_+$, unique up to constant.
\end{theorem}

\begin{proof}
By the shifted symplectic neighborhood hypothesis, the normal complex of $\mathcal{L}_+$
in $\mathcal{M}$ is canonically identified with
$T^*[-1]\mathcal{L}_+$. The symplectic neighborhood theorem
produces the pointed equivalence. Under this identification
$\mathcal{L}_+$ becomes the zero section; since $\mathcal{L}_-$
is complementary, it is the graph of a one-form. That one-form is
closed because the graph of a one-form in a cotangent bundle is
Lagrangian if and only if the one-form is closed. The formal
Poincare hypothesis makes this closed one-form exact.
\end{proof}

\begin{definition}[Complementarity potential: existence form]
\label{def:complementarity-potential-existence}
Assume the ambient complementarity formal moduli problem exists.
The formal function
\[
S_{\cA} \in \widehat{\mathcal{O}}(\mathcal{M}_{\cA})^0
\]
produced by Theorem~\ref{thm:shifted-cotangent-normal-form} for
$\mathcal{L}_+ = \mathcal{M}_{\cA}$,
$\mathcal{L}_- = \mathcal{M}_{\cA^!}$
is called the \emph{complementarity potential} of the Koszul pair
$(\cA, \cA^!)$.
\end{definition}

\begin{proposition}[Legendre duality of the two potentials]
\label{prop:legendre-duality-potentials}
\ClaimStatusProvedHere
Let $S_{\cA}$ and $S_{\cA^!}$ be the complementarity potentials
constructed from the two cotangent charts. Then $S_{\cA}$ and
$S_{\cA^!}$ are formal Legendre transforms of one another.
\end{proposition}

\begin{proof}
The two functions generate the same formal Lagrangian subspace of
the same shifted symplectic ambient moduli problem, viewed from
opposite cotangent charts. The coordinate change between the charts
is the formal Legendre transform.
\end{proof}

\begin{proposition}[Legendre duality of cubic tensors]
\label{prop:legendre-duality-cubic}
\ClaimStatusProvedHere
Let $H_{\cA}$ be the complementarity Hessian and
$\mathfrak{C}_{\cA}$ the cubic shadow. Then the cubic shadows
of $\cA$ and $\cA^!$ are related by
\[
\mathfrak{C}_{\cA^!}(p,p,p)
= -\mathfrak{C}_{\cA}(H_{\cA}^{-1}p,\, H_{\cA}^{-1}p,\,
H_{\cA}^{-1}p).
\]
More generally, at quartic order the Legendre transform gives
$S_{\cA}^\vee(p)
= \frac12 H_{\cA}^{-1}(p,p)
- \frac{1}{6}\mathfrak{C}_{\cA}
(H_{\cA}^{-1}p, H_{\cA}^{-1}p, H_{\cA}^{-1}p) + O(p^4)$.
\end{proposition}

\begin{proof}
Write $S_{\cA}(x) = \frac12 H_{\cA}(x,x)
+ \frac16 \mathfrak{C}_{\cA}(x,x,x) + O(x^4)$.
The critical point of $\langle p,x\rangle - S_{\cA}(x)$ is
$x = H_{\cA}^{-1}p + O(p^2)$. Substituting into the
Legendre formula and collecting through cubic order yields the
identity.
\end{proof}

\subsection{Derived critical locus and fake complementarity}
\label{subsec:derived-critical-locus}

\begin{theorem}[Derived critical locus of self-dual deformations]
\label{thm:derived-critical-locus}
\ClaimStatusProvedHere
Assume the ambient complementarity formal moduli problem exists.
Then there is a canonical equivalence
\[
\mathcal{M}_{\cA}
\times^h_{\mathcal{M}_{\mathrm{comp}}(\cA)}
\mathcal{M}_{\cA^!}
\;\simeq\;
\operatorname{Crit}(S_{\cA}).
\]
Simultaneous deformations of $\cA$ and $\cA^!$ preserving
complementarity are governed by the derived critical locus of the
complementarity potential.
\end{theorem}

\begin{proof}
In the cotangent normal form one Lagrangian is the zero section and
the other is the graph of $dS_{\cA}$. Their derived intersection
is the derived zero locus of $dS_{\cA}$, which is the derived
critical locus of $S_{\cA}$.
\end{proof}

\begin{proposition}[Criterion for fake complementarity]
\label{prop:fake-complementarity-criterion}
\ClaimStatusProvedHere
In the cotangent chart centered at $\mathcal{M}_{\cA}$, the
following are equivalent:
\begin{enumerate}[label=\textup{(\roman*)}]
\item the dual Lagrangian $\mathcal{M}_{\cA^!}$ is linear;
\item the complementarity potential $S_{\cA}$ is quadratic;
\item the derived critical locus $\operatorname{Crit}(S_{\cA})$
is controlled entirely by its Hessian complex, with no higher
nonlinear operations.
\end{enumerate}
The first obstruction to fake complementarity is the first nonzero
higher jet of $S_{\cA}$ beyond its quadratic part.
\end{proposition}

\begin{proof}
In the cotangent chart, $\mathcal{M}_{\cA^!}$ is the graph of
$dS_{\cA}$. The graph is linear iff $dS_{\cA}$ is linear,
iff $S_{\cA}$ is quadratic. When $S_{\cA}$ is quadratic, the
critical locus is cut out by a linear differential and is
governed by the Hessian alone.
\end{proof}

\begin{remark}[Shadow jets as the quantum $L_\infty$ structure
made visible]
\label{rem:shadow-jets-linfty}
\index{quantum $L_\infty$!shadow jets as}
The complementarity potential
$S_{\cA}(x) = \sum_{r \geq 2}(1/r!)\operatorname{Sh}_r(x^r)$
encodes the full quantum $L_\infty$ structure of
Theorem~\ref{thm:modular-quantum-linfty} in a single generating
function. The precise dictionary is:
\begin{center}
\begin{tabular}{@{}lll@{}}
\toprule
Shadow jet & $L_\infty$ operation & Geometric meaning \\
\midrule
$H_{\cA} = \operatorname{Sh}_2$ &
$\ell_2^{(0)}$ restricted to cyclic slice &
Hessian of complementarity \\
$\mathfrak{C}_{\cA} = \operatorname{Sh}_3$ &
$\ell_3^{(0)} = \ell_2^{\mathrm{tr}}$ (transferred) &
cubic bending of dual Lagrangian \\
$\mathfrak{Q}_{\cA} = \operatorname{Sh}_4$ &
$\ell_4^{(0)} = \ell_3^{\mathrm{tr}}$ (transferred) &
quartic contact correction \\
$\Lambda_P(\mathfrak{Q}_{\cA})$ &
$\ell_1^{(1)}$ applied to $\operatorname{Sh}_4$ &
genus-$1$ Hessian correction \\
$\operatorname{Sh}_r$ for $r \geq 5$ &
$\ell_r^{(0)}$ (higher homotopy) &
higher nonlinear modular data \\
\bottomrule
\end{tabular}
\end{center}
The vanishing or nonvanishing of each shadow jet
$\operatorname{Sh}_r$ is therefore a question about the transferred
$L_\infty$ operations $\ell_r^{\mathrm{tr}}$. The archetype
trichotomy
(Theorem~\ref{thm:w-archetype-trichotomy}) classifies families
by which transferred operations survive:
\begin{itemize}
\item Gaussian: all $\ell_n^{\mathrm{tr}} = 0$ for $n \geq 2$
(Heisenberg);
\item Lie/tree: $\ell_2^{\mathrm{tr}} \neq 0$,
$\ell_n^{\mathrm{tr}} = 0$ for $n \geq 3$ (affine);
\item Contact/quartic: $\ell_2^{\mathrm{tr}} = 0$,
$\ell_3^{\mathrm{tr}} \neq 0$ ($\beta\gamma$);
\item Mixed: $\ell_2^{\mathrm{tr}} \neq 0$ \emph{and}
$\ell_3^{\mathrm{tr}} \neq 0$ (Virasoro, $\mathcal{W}_N$).
\end{itemize}
The shadow obstruction tower is the homotopy structure made visible.
\end{remark}

\begin{remark}[The three archetypes from the Lagrangian perspective]
\label{rem:three-archetypes-lagrangian}
The fake-complementarity criterion reframes the archetype trichotomy
(Theorem~\ref{thm:nms-archetype-trichotomy}) in geometric language.
The Heisenberg is exactly fake: $S_{\mathcal{H}}$ is quadratic.
Affine $\widehat{\mathfrak{sl}}_2$ is cubic at leading order:
$S_{\mathrm{aff}} = \tfrac12 H + \tfrac16\kappa(a,[a,a])$,
so the dual Lagrangian bends cubically. The $\beta\gamma$ system is
quartic on the weight/contact slice. Virasoro is the first family
where the dual Lagrangian has both cubic and quartic bending,
and the complementarity potential is genuinely non-polynomial; its
shadow obstruction tower is infinite (Theorem~\ref{thm:nms-finite-termination}).
\end{remark}

\begin{remark}[Shadow obstruction tower as period correction]
\label{rem:shadow-tower-period-correction}
\index{period correction!shadow tower}
The shadow hierarchy
$\kappa \to \Delta \to \mathfrak{C}
\to \mathfrak{Q} \to \Theta$
admits a concrete interpretation as the Taylor expansion of the
\emph{full period correction} that restores nilpotency of the
genus-$g$ bar differential. At genus~$g$, the fiberwise
differential satisfies $d_{\mathrm{fib}}^2 = \kappa \cdot \omega_g$.
The total corrected differential
$\mathbb{D}_g = d_0 + \sum_{n \geq 1} t_g^{(n)} d_n$
absorbs this curvature:
\begin{itemize}
\item $t_g^{(1)}$: determined by $\kappa$ alone (the $m_2$
 contribution; the Faber--Pandharipande term);
\item $t_g^{(2)}$: involves $m_3$ paired with genus-$g$ propagators,
 giving the cubic shadow $\mathfrak{C}_{\cA}$;
\item $t_g^{(3)}$: involves $m_4$ and $m_3 \circ m_3$ terms,
 giving the quartic shadow $\mathfrak{Q}_{\cA}$;
\item $t_g^{(\infty)}$: the full series is $\Theta_{\cA}$,
 the universal Maurer--Cartan element.
\end{itemize}
The subleading corrections for $n \geq 2$ are \emph{not} determined
by $\kappa$: they depend on the higher $\Ainf$-operations interacting
with the Hodge data of $\overline{\mathcal{M}}_g$. This is the
mechanism by which algebra-dependent modular information enters
beyond the $\hat{A}$-genus.
\end{remark}

\section{Finite-rank dual comparison and the metaplectic anomaly}
\label{sec:holographic-complementarity-metaplectic}
\index{dual comparison!finite-rank|textbf}
\index{metaplectic correction!Gaussian|textbf}

In the finite-rank paired regime, complementarity admits a
linear comparison model. The decomposition
$\mathbf{C}_g(\cA) \simeq \mathbf{Q}_g(\cA) \oplus
\mathbf{Q}_g(\cA^!)$ of Theorem~\ref{thm:quantum-complementarity-main}
is an idempotent splitting of the genus-$g$ ambient center complex.
Under the non-degenerate Verdier pairing hypotheses it becomes a
pair of dual Lagrangian polarizations: one branch is the
$+1$-eigenspace of the represented Verdier involution, the other is
the $-1$-eigenspace, and the two branches exhaust the ambient
cohomology of moduli space. This is a typed linear comparison. No
bulk factorization algebra is constructed here.

The linear model is the following. The genus-$g$ ambient complex
$\mathbf{C}_g(\cA) = R\Gamma(\overline{\mathcal{M}}_g,
\mathcal{Z}(\cA))$ is a shifted-symplectic linear space in the sense of
Proposition~\ref{prop:ptvv-lagrangian}, with a
$(-(3g{-}3))$-shifted symplectic form.
The two complementary branches $\mathbf{Q}_g(\cA)$ and
$\mathbf{Q}_g(\cA^!)$ are Lagrangian subspaces of this phase space,
and the finite-rank passage from one polarization to the other is the
formal Fourier transform. Quantizing the linearized phase space
produces a Heisenberg algebra, a Fock representation, a Weyl algebra,
and a Fourier kernel that intertwines creation with annihilation.

When one goes beyond the linear theory (when two Riemann surfaces
are sewn together and the intermediate phase space is integrated
out), the composition of Gaussian kernels ceases to be strictly
multiplicative. It picks up a determinant factor, a square-root
correction $\det(1 - C_1 B_2)^{-1/2}$ that measures the failure of
naive composition. This determinant anomaly satisfies a $2$-cocycle
law and defines a central extension of the sewing semigroup: the
metaplectic extension. A choice of metaplectic splitting
trivializes the cocycle and restores strict composition, but the
anomaly itself is an invariant: it is the first genuinely nonlinear
determinant correction, detecting algebra-dependent structure
invisible to the scalar curvature $\kappa$.

\medskip
\noindent\emph{Finite-rank hypothesis.}\;
Every theorem in this section is derived from the proved core of
Theorems~A--D together with finite-rank linear algebra on perfect
dg models. The only additional assumption is that we work on
finite-rank parity-even models when carrying out explicit Gaussian
calculations; the graded extension follows by the standard Koszul
sign rule. No conjectural input is required until the closing
conjecture on bulk enhancement.

\subsection{The protected dual transform}
\label{subsec:holo-comp-protected-transform}

\begin{definition}[Protected dual transform]
\label{def:holo-comp-protected-transform}
\index{dual transform!protected|textbf}
Let $X$ be a smooth projective curve and let $\cA$ be a modular
Koszul chiral algebra on~$X$. The \emph{protected dual transform} is
\begin{equation}\label{eq:holo-comp-protected-transform}
\mathbb{H}_X(\cA)
:=
\mathbb{D}_{\operatorname{Ran}} \, \barB_X(\cA)
\in D^{\mathrm{co}}\!\bigl(\operatorname{Fact}(X)\bigr).
\end{equation}
For each $g \geq 0$, the \emph{dual cofiber} is
\begin{equation}\label{eq:holo-comp-cofiber}
\mathbf{H}^{\mathrm{hol}}_g(\cA)
:=
\cofib\!\bigl(\mathbf{Q}_g(\cA)
\longrightarrow \mathbf{C}_g(\cA)\bigr).
\end{equation}
\end{definition}

\begin{theorem}[Protected dual transform;
\ClaimStatusConditional]
\label{thm:holo-comp-bulk-reconstruction}
\index{dual transform!protected|textbf}
Let $\cA$ be a modular Koszul chiral algebra on~$X$. Assume the
Verdier bar--cobar comparison and higher-genus inversion are available
for $\cA$ and $\cA^!$ on the Koszul locus, and form the cofiber in the
strict flat ambient representative. Then:
\begin{enumerate}[label=\textup{(\roman*)}]
\item \emph{Dual transform.}\;
 $\mathbb{H}_X(\cA) \simeq \cA^!$.
\item \emph{Anti-involution.}\;
 $\mathbb{H}_X^2(\cA) \simeq \cA$.
\item \emph{Cofiber identification.}\;
 For every $g \geq 0$,
 $\mathbf{H}^{\mathrm{hol}}_g(\cA) \simeq \mathbf{Q}_g(\cA^!)$.
\end{enumerate}
In particular, the protected transform $\mathbb{H}_X$ identifies
canonically with the Koszul-dual branch and is a contravariant
involution in the stated coderived category. No bulk factorization
algebra is constructed in this theorem.
\end{theorem}

\begin{proof}
By Verdier intertwining
(Theorem~\ref{thm:bar-cobar-isomorphism-main}) and higher-genus
inversion (Theorem~\ref{thm:higher-genus-inversion}),
\[
\mathbb{H}_X(\cA)
=
\mathbb{D}_{\operatorname{Ran}} \, \barB_X(\cA)
\simeq
\cA^!.
\]
Applying the same argument to $\cA^!$ and using
$(\cA^!)^! \simeq \cA$ gives the anti-involution. For the cofiber,
the involution-splitting (Lemma~\ref{lem:involution-splitting})
gives $\mathbf{C}_g(\cA) \simeq \mathbf{Q}_g(\cA) \oplus
\mathbf{Q}_g(\cA^!)$; the canonical inclusion of the first summand
has cofiber the second.
\end{proof}

\subsection{Shifted cotangent realization and spectral reciprocity}
\label{subsec:holo-comp-cotangent-spectral}

Fix a genus $g \geq 1$ and write
$d_g := 3g - 3$,
$V_g(\cA) := \mathbf{Q}_g(\cA)$,
$W_g(\cA) := \mathbf{Q}_g(\cA^!)[d_g]$.
By Theorem~\ref{thm:quantum-complementarity-main},
$W_g(\cA) \simeq V_g(\cA)^\vee$.

\begin{theorem}[Shifted cotangent realization;
\ClaimStatusConditional]
\label{thm:holo-comp-cotangent-realization}
\index{cotangent realization!shifted|textbf}
Let $\cA$ be a modular Koszul chiral algebra. Then:
\begin{enumerate}[label=\textup{(\roman*)}]
\item \emph{Cotangent decomposition.}\;
 $\mathbf{C}_g(\cA) \simeq V_g(\cA) \oplus V_g(\cA)^\vee[-d_g]$.
\item \emph{Transposition.}\;
 $\mathbf{H}^{\mathrm{hol}}_g(\cA) \simeq V_g(\cA)^\vee[-d_g]$.
\item \emph{Protected self-duality.}\;
 If $\mathbb{H}_X(\cA) \simeq \cA$, then
 $V_g(\cA) \simeq V_g(\cA)^\vee[-d_g]$.
\end{enumerate}
\end{theorem}

\begin{proof}
By complementarity,
$\mathbf{Q}_g(\cA^!) \simeq W_g(\cA)[-d_g] \simeq
V_g(\cA)^\vee[-d_g]$. Substituting into the direct-sum
decomposition proves (i); (ii) follows from
Theorem~\ref{thm:holo-comp-bulk-reconstruction}(iii);
(iii) follows because protected self-duality identifies $\cA$ with
$\cA^!$.
\end{proof}

The cotangent realization has immediate spectral consequences.

\begin{corollary}[Spectral reciprocity and palindromicity;
\ClaimStatusProvedHere]
\label{cor:holo-comp-spectral-reciprocity}
\index{spectral reciprocity!palindromicity|textbf}
Assume $V_g(\cA)$ has finite-dimensional cohomology and set
$P_g(\cA; u) := \sum_i \dim H^i(V_g(\cA))\, u^i$. Then:
\begin{enumerate}[label=\textup{(\roman*)}]
\item \emph{Spectral reciprocity.}\;
 $P_g(\cA^!; u) = u^{d_g}\, P_g(\cA; u^{-1})$.
\item \emph{Protected self-dual palindromicity.}\;
 If $\mathbb{H}_X(\cA) \simeq \cA$, then
 $P_g(\cA; u) = u^{d_g}\, P_g(\cA; u^{-1})$.
\item \emph{Even-genus Euler vanishing.}\;
 If $\mathbb{H}_X(\cA) \simeq \cA$ and $g$ is even, then
 $\chi(V_g(\cA)) = 0$.
\item \emph{Determinant parity.}\;
 \begin{equation}\label{eq:holo-comp-det-parity}
 \Det\!\bigl(\mathbf{C}_g(\cA)\bigr)
 \cong
 \begin{cases}
 \C, & g \text{ odd}, \\[4pt]
 \Det\!\bigl(V_g(\cA)\bigr)^{\otimes 2}, & g \text{ even}.
 \end{cases}
 \end{equation}
\end{enumerate}
\end{corollary}

\begin{proof}
The transposition
$\mathbf{Q}_g(\cA^!) \simeq V_g(\cA)^\vee[-d_g]$ gives (i) on
cohomology. Protected self-duality forces palindromicity.
Evaluating at $u = -1$ and using the parity of $d_g = 3g - 3$
(odd when $g$ is even) gives (iii). For (iv), duality gives
$\Det(V^\vee) \cong \Det(V)^{-1}$, and a shift by $[-d_g]$
inverts the determinant line iff $d_g$ is odd. The direct-sum
formula then yields~\eqref{eq:holo-comp-det-parity}.
\end{proof}

\begin{remark}[The palindromic constraint on standard families]
\label{rem:holo-comp-palindromic-examples}
The Virasoro algebra is chiral Koszul self-dual at $c = 13$
($\mathrm{Vir}_c^! = \mathrm{Vir}_{26-c}$). At the self-dual
point the Poincar\'e polynomial of $V_g(\mathrm{Vir}_{13})$ is
forced to be palindromic when the protected self-duality hypothesis
holds. For the Heisenberg algebra, which is not protected self-dual,
the reciprocity law instead relates
$V_g(\mathcal{H})$ to $V_g(\mathcal{H}^!)$ via the shift by $d_g$;
the two eigenspaces can carry different spectral data.
\end{remark}

\subsection{The Heisenberg--Fock--Fourier--Weyl quantization package}
\label{subsec:holo-comp-quantization-package}

The shifted cotangent realization promotes the genus-$g$ ambient
complex from a mere direct sum to a linearized phase space. The
natural next step is to quantize it.

\begin{definition}[Genus-$g$ Heisenberg algebra and Fock object]
\label{def:holo-comp-heisenberg-fock}
\index{Heisenberg algebra!genus-$g$|textbf}
\index{Fock object!genus-$g$|textbf}
Write $V := V_g(\cA)$, $W := W_g(\cA) \simeq V^\vee$.
\begin{enumerate}[label=\textup{(\roman*)}]
\item The \emph{genus-$g$ Heisenberg dg Lie algebra} is
 $\mathfrak{h}_g(\cA) := V \oplus W \oplus \C\mathbf{1}$
 with $\mathbf{1}$ central, $V$ and $W$ abelian, and bracket
 $[w, v] = w(v)\,\mathbf{1}$.
\item The \emph{primal Fock object} is
 $\mathfrak{F}^\partial_g(\cA) :=
 \bigoplus_{n \geq 0} \operatorname{Sym}^n(V)$
 and the \emph{dual Fock object} is
 $\mathfrak{F}^{\mathrm{bulk}}_g(\cA) :=
 \bigoplus_{n \geq 0} \operatorname{Sym}^n(W)$.
\item The \emph{genus-$g$ Fourier kernel} is
 $\mathbf{K}_g^{\mathrm{hol}}(\cA) :=
 \exp(\operatorname{coev}_g(\cA))$
 in the completed tensor product
 $\widehat{\mathfrak{F}}{}^\partial_g(\cA)
 \,\widehat{\otimes}\,
 \widehat{\mathfrak{F}}{}^{\mathrm{bulk}}_g(\cA)$,
 where
 $\operatorname{coev}_g(\cA) \in V \otimes W$ corresponds to the
 identity of~$V$.
\end{enumerate}
\end{definition}

\begin{theorem}[Fourier intertwining;
\ClaimStatusProvedHere]
\label{thm:holo-comp-fourier-transport}
\index{Fourier kernel!intertwining|textbf}
Let $v \in V$, $w \in W \simeq V^\vee$. Then the Fourier kernel
intertwines creation and annihilation:
\begin{equation}\label{eq:holo-comp-fourier-intertwining}
(\iota_w \otimes 1)\,\mathbf{K}_g^{\mathrm{hol}}(\cA)
=
(1 \otimes m_w)\,\mathbf{K}_g^{\mathrm{hol}}(\cA),
\qquad
(m_v \otimes 1)\,\mathbf{K}_g^{\mathrm{hol}}(\cA)
=
(1 \otimes \iota_v)\,\mathbf{K}_g^{\mathrm{hol}}(\cA),
\end{equation}
where $m_v$ is multiplication and $\iota_w$ is contraction.
\end{theorem}

\begin{proof}
The kernel $\mathbf{K}_g^{\mathrm{hol}}(\cA) =
\sum_{n \geq 0} (1/n!)\,\operatorname{coev}_g(\cA)^{\odot n}$
is a formal exponential of the coevaluation tensor. Contraction
on the left by $w$ removes one tensor factor and returns
multiplication by $w$ on the right; multiplication by $v$ on the
left is dually contraction by $v$ on the right.
\end{proof}

\begin{definition}[Genus-$g$ Weyl algebra and Wick product]
\label{def:holo-comp-weyl-wick}
\index{Weyl algebra!genus-$g$|textbf}
\index{Wick product!genus-$g$|textbf}
The completed \emph{Weyl object} is
$\widehat{\mathscr{W}}_g(\cA) :=
\prod_{n \geq 0} \operatorname{Sym}^n(V \oplus W)[[\hbar]]$.
The evaluation pairing $W \otimes V \to \C$ defines a continuous
biderivation $\Gamma_g$ on the tensor square, and the
\emph{Wick star product} is
\begin{equation}\label{eq:holo-comp-wick-star}
f \star_{g,\hbar} h
:=
\mu \circ e^{\hbar \Gamma_g}(f \otimes h).
\end{equation}
\end{definition}

\begin{theorem}[Weyl associativity, PBW, and linear sewing;
\ClaimStatusProvedHere]
\label{thm:holo-comp-weyl-sewing}
\index{PBW theorem!Weyl algebra|textbf}
\index{sewing!linear Fourier kernel|textbf}
For each $g \geq 1$:
\begin{enumerate}[label=\textup{(\roman*)}]
\item $\star_{g,\hbar}$ is associative and unital.
\item $\operatorname{gr}(\widehat{\mathscr{W}}_g(\cA),
 \star_{g,\hbar}) \cong \operatorname{Sym}(V \oplus W)[[\hbar]]$
 (PBW).
\item The generator relations are
 $w \star_{g,\hbar} v = wv + \hbar\, w(v)$,
 while $v \star v' = vv'$ and $w \star w' = ww'$.
\item The linear Fourier kernels compose exactly:
 $\langle \mathbf{K}_{12}^{\mathrm{lin}},
 \mathbf{K}_{23}^{\mathrm{lin}}\rangle_{(2)}
 = \mathbf{K}_{13}^{\mathrm{lin}}$.
\end{enumerate}
\end{theorem}

\begin{proof}
Associativity: pairwise contraction operators commute, so both
parenthesizations of a triple product equal
$\mu_{123} \circ \exp(\hbar(\Gamma_{12} + \Gamma_{13} +
\Gamma_{23}))$. PBW: each Wick contraction lowers symmetric degree
by $2$. The generator relations follow because the only nontrivial
contraction pairs a left $W$-generator with a right $V$-generator.
The linear sewing law is the $B = C = 0$ specialization of the
Gaussian composition theorem below.
\end{proof}

\subsection{Gaussian composition, Schur complement, and the
determinant anomaly}
\label{subsec:holo-comp-gaussian-metaplectic}

The linear sewing law is exact but flat: it sees only the
coevaluation tensor and misses the nonlinear part of the comparison.
The first enrichment is to allow quadratic data. Once the kernels
carry quadratic terms, composition ceases to be anomaly-free and
acquires a determinant correction, the metaplectic anomaly.

\begin{convention}[Parity-even models]
\label{conv:holo-comp-parity-even}
For the rest of this subsection we work on a fixed finite-rank
perfect dg model for $V = V_g(\cA)$ and identify $W$ with $V^\vee$.
Formulas are stated in the parity-even sector; the graded extension
follows by the Koszul sign rule.
\end{convention}

\begin{definition}[Gaussian triple and Gaussian kernel]
\label{def:holo-comp-gaussian-triple}
\index{Gaussian kernel!finite-rank|textbf}
Let $V$ be an $n$-dimensional complex vector space with coordinates
$x = (x_1, \ldots, x_n)^t$ and dual coordinates
$\eta = (\eta_1, \ldots, \eta_n)^t$ on $W = V^\vee$. A
\emph{Gaussian triple} is a triple $(B, A, C)$ with
$B \in \operatorname{Hom}(V, W)$,
$A \in \operatorname{End}(V)$,
$C \in \operatorname{Hom}(W, V)$,
where $B$ and $C$ are symmetric. The associated
\emph{Gaussian kernel} is
\begin{equation}\label{eq:holo-comp-gaussian-kernel}
\mathbf{K}(B, A, C)(x, \eta)
:=
\exp\!\Bigl(
\tfrac{1}{2}\, x^t B x + x^t A^t \eta
+ \tfrac{1}{2}\, \eta^t C \eta
\Bigr).
\end{equation}
The linear Fourier kernel of the preceding subsection is the
special case $B = C = 0$, $A = \mathrm{id}$. The quadratic data
$(B, C)$ encode the first departure from linearity: they measure
how the two Lagrangian branches of complementarity bend away from
the linear decomposition.
\end{definition}

\begin{theorem}[Gaussian composition via Schur complement;
\ClaimStatusProvedHere]
\label{thm:holo-comp-gaussian-composition}
\index{Schur complement!Gaussian composition|textbf}
\index{determinant anomaly!Gaussian sewing|textbf}
Let $\mathbf{K}_1 = \mathbf{K}(B_1, A_1, C_1)(x, \eta)$ and
$\mathbf{K}_2 = \mathbf{K}(B_2, A_2, C_2)(\xi, y)$ be Gaussian
kernels and assume $1 - C_1 B_2$ is invertible. Then their middle
contraction is again Gaussian:
\begin{equation}\label{eq:holo-comp-gaussian-composition}
\langle \mathbf{K}_1, \mathbf{K}_2 \rangle_{(2)}
=
\det(1 - C_1 B_2)^{-1/2}\,
\mathbf{K}(B_{12}, A_{12}, C_{12})(x, y),
\end{equation}
where the composed triple $(B_{12}, A_{12}, C_{12})$ is given by
the Schur complements
\begin{align}
B_{12} &:= B_1 + A_1^t B_2(1 - C_1 B_2)^{-1} A_1,
\label{eq:holo-comp-schur-B} \\
A_{12}^t &:= A_1^t(1 - B_2 C_1)^{-1} A_2^t,
\label{eq:holo-comp-schur-A} \\
C_{12} &:= C_2 + A_2 C_1(1 - B_2 C_1)^{-1} A_2^t.
\label{eq:holo-comp-schur-C}
\end{align}
In particular, the linear sewing law of
Theorem~\textup{\ref{thm:holo-comp-weyl-sewing}(iv)} is recovered
by setting $B_1 = C_1 = B_2 = C_2 = 0$, $A_1 = A_2 = \mathrm{id}$.
\end{theorem}

\begin{proof}
Apply the formal Gaussian differential identity: for symmetric
matrices $S$, $T$ and vectors $a$, $b$ with $1 - ST$ invertible,
\[
\left.
\exp\!\Bigl(\tfrac{1}{2}\,\partial_z^t T \partial_z
+ b^t \partial_z\Bigr)
\exp\!\Bigl(a^t z + \tfrac{1}{2}\, z^t S z\Bigr)
\right|_{z=0}
=
\det(1 - ST)^{-1/2}
\exp\!\Bigl(\tfrac{1}{2}\, b^t T(1{-}ST)^{-1} b
+ a^t(1{-}ST)^{-1} b
+ \tfrac{1}{2}\, a^t S(1{-}TS)^{-1} a\Bigr).
\]
This identity is proved by the Gaussian heat-equation ansatz: the
one-parameter family
$F(u,z) := \exp(u\,\partial_z^t T \partial_z/2)
\exp(a^t z + z^t S z / 2)$
satisfies a heat equation whose solution via a Gaussian ansatz
reduces to the Riccati system $S'(u) = S(u) T S(u)$,
$a'(u) = S(u) T a(u)$,
$c'(u) = \tfrac{1}{2} a(u)^t T a(u)
+ \tfrac{1}{2}\operatorname{Tr}(T S(u))$,
with explicit solutions
$S(u) = S(1 - uTS)^{-1}$,
$a(u) = (1 - uST)^{-1} a$,
$c(u) = \tfrac{1}{2} a^t T(1 - uST)^{-1} a
- \tfrac{1}{2}\log\det(1 - uST)$.
Setting $S = C_1$, $T = B_2$, $a = A_1 x$, $b = A_2^t y$, and
evaluating at $u = 1$ yields exactly the
composition~\eqref{eq:holo-comp-gaussian-composition} with
Schur-complement
coefficients~\eqref{eq:holo-comp-schur-B}--\eqref{eq:holo-comp-schur-C}.
\end{proof}

\begin{remark}[Physical interpretation of the Schur complement]
\label{rem:holo-comp-schur-interpretation}
\index{Schur complement!physical interpretation}
The formulas~\eqref{eq:holo-comp-schur-B}--\eqref{eq:holo-comp-schur-C}
have a natural interpretation: they are the Schur complements of the
middle quadratic block in a $3 \times 3$ block matrix. The
determinant factor $\det(1 - C_1 B_2)^{-1/2}$ is the precise
algebraic cost of integrating out the intermediate finite-rank phase
space. At the level of Gaussian integration, this is a one-loop
determinant:
the intermediate $\eta$-variables have been traced over, and the
resulting effective kernel on the remaining $(x, y)$-variables
carries a fluctuation determinant.
\end{remark}

\subsection{The metaplectic cocycle and its strictification}
\label{subsec:holo-comp-metaplectic-cocycle}

\begin{theorem}[Metaplectic $2$-cocycle and strictification;
\ClaimStatusProvedHere]
\label{thm:holo-comp-metaplectic-cocycle}
\index{cocycle!metaplectic|textbf}
\index{metaplectic correction!strictification|textbf}
Let $G_i = (B_i, A_i, C_i)$ ($i = 1, 2, 3$) be Gaussian triples
with all pairwise compositions defined. Write
$\alpha_g(C, B) := \det(1 - CB)^{-1/2}$. Then:
\begin{enumerate}[label=\textup{(\roman*)}]
\item \emph{$2$-cocycle law.}\;
 \begin{equation}\label{eq:holo-comp-cocycle-law}
 \alpha_g(C_1, B_2)\,\alpha_g(C_{12}, B_3)
 =
 \alpha_g(C_2, B_3)\,\alpha_g(C_1, B_{23}),
 \end{equation}
 where $(B_{12}, A_{12}, C_{12}) = G_1 \sharp G_2$ and
 $(B_{23}, A_{23}, C_{23}) = G_2 \sharp G_3$ are the Schur-composed
 triples.
\item \emph{Central extension.}\;
 The law
 $(\lambda_1, G_1) \cdot (\lambda_2, G_2)
 := (\lambda_1 \lambda_2 \alpha_g(C_1, B_2),\, G_1 \sharp G_2)$
 defines an associative central extension of the Gaussian semigroup
 by $\C^\times$.
\item \emph{Strictification.}\;
 If a subsemigroup admits a $1$-cochain $\mu_g$ satisfying
 $\alpha_g(C_1, B_2) = \mu_g(G_1)\,\mu_g(G_2)\,
 \mu_g(G_1 \sharp G_2)^{-1}$,
 then the corrected kernels
 $\widetilde{\mathbf{K}}(G) := \mu_g(G)^{-1}\,\mathbf{K}(G)$
 compose strictly:
 $\langle \widetilde{\mathbf{K}}(G_1),
 \widetilde{\mathbf{K}}(G_2) \rangle
 = \widetilde{\mathbf{K}}(G_1 \sharp G_2)$.
\end{enumerate}
\end{theorem}

\begin{proof}
Associativity of middle contraction gives
$\langle\langle \mathbf{K}(G_1), \mathbf{K}(G_2)\rangle,
\mathbf{K}(G_3)\rangle
= \langle \mathbf{K}(G_1),
\langle \mathbf{K}(G_2), \mathbf{K}(G_3)\rangle\rangle$.
Expanding both sides via
Theorem~\ref{thm:holo-comp-gaussian-composition} and comparing
the scalar prefactors yields~\eqref{eq:holo-comp-cocycle-law}.
This is exactly the associativity condition for the central
extension. If $\mu_g$ splits the cocycle, the anomaly factors
cancel in the corrected kernels.
\end{proof}

\begin{remark}[The metaplectic extension as geometry]
\label{rem:holo-comp-metaplectic-geometry}
\index{metaplectic extension!geometric meaning}
The metaplectic splitting is not a decorative normalization. It is
the linear-model avatar of the choice of square root of the
determinant line $\Det(V_g(\cA))$. On $\overline{\mathcal{M}}_g$,
this choice corresponds to a spin structure on the Hodge bundle,
the same spin structure that enters the definition of the Mumford
isomorphism $\lambda_g^{\otimes c/2} \cong
\Det(R\pi_* \mathcal{Z}(\cA))$. The metaplectic $2$-cocycle is
therefore the sewing obstruction of the half-determinant: it
measures the failure of $\Det^{1/2}$ to be multiplicative under
gluing of Riemann surfaces.
\end{remark}

\subsection{The first nonlinear determinant anomaly}
\label{subsec:holo-comp-first-nonlinear-anomaly}

\begin{corollary}[First nonlinear determinant anomaly;
\ClaimStatusProvedHere]
\label{cor:holo-comp-first-nonlinear-anomaly}
\index{nonlinear anomaly!first determinant|textbf}
Let $G_1$, $G_2$ be composable Gaussian triples. The trace-log
expansion of the anomaly is
\begin{equation}\label{eq:holo-comp-trace-log-anomaly}
\log \alpha_g(C_1, B_2)
=
\frac{1}{2}\operatorname{Tr}(C_1 B_2)
+ \frac{1}{4}\operatorname{Tr}\!\bigl((C_1 B_2)^2\bigr)
+ \frac{1}{6}\operatorname{Tr}\!\bigl((C_1 B_2)^3\bigr)
+ \cdots.
\end{equation}
The first genuinely nonlinear determinant anomaly is the quadratic
term
\begin{equation}\label{eq:holo-comp-first-nonlinear}
\mathfrak{N}_g(G_1, G_2)
:=
\frac{1}{4}\operatorname{Tr}\!\bigl((C_1 B_2)^2\bigr).
\end{equation}
This vanishes precisely on the square-zero locus $(C_1 B_2)^2 = 0$,
where composition is controlled by the one-loop trace
$\tfrac{1}{2}\operatorname{Tr}(C_1 B_2)$ alone.
\end{corollary}

\begin{proof}
The standard power-series identity
$-\tfrac{1}{2}\log\det(1 - T) =
\tfrac{1}{2}\sum_{m \geq 1} m^{-1}\operatorname{Tr}(T^m)$
applied to $T = C_1 B_2$
gives~\eqref{eq:holo-comp-trace-log-anomaly}. The linear trace
$\tfrac{1}{2}\operatorname{Tr}(C_1 B_2)$ is already visible in the
one-loop sewing law; the first new information beyond one loop is
$\mathfrak{N}_g$.
\end{proof}

\begin{remark}[Hierarchy of Gaussian complexity]
\label{rem:holo-comp-hierarchy}
\index{Gaussian complexity!hierarchy}
The trace-log expansion~\eqref{eq:holo-comp-trace-log-anomaly}
provides a precise hierarchy of Gaussian complexity, organized
by powers of the product $C_1 B_2$:
\begin{center}
\begin{tabular}{@{}lll@{}}
\toprule
Order & Anomaly term & Geometric content \\
\midrule
$0$ & $\alpha_g = 1$ (linear sewing)
 & coevaluation only; no bending \\
$1$ & $\tfrac{1}{2}\operatorname{Tr}(C_1 B_2)$
 & one-loop; scalar curvature level \\
$2$ & $\tfrac{1}{4}\operatorname{Tr}((C_1 B_2)^2)$
 & first nonlinear; detects $\mathfrak{Q}$-level data \\
$m$ & $\tfrac{1}{2m}\operatorname{Tr}((C_1 B_2)^m)$
 & $m$-th loop order; degree-$(2m)$ shadow data \\
\bottomrule
\end{tabular}
\end{center}
The connection with the shadow obstruction tower
(\S\ref{subsec:derived-critical-locus},
Remark~\ref{rem:shadow-tower-period-correction}) is as follows.
The quadratic data $(B, C)$ of a Gaussian kernel encode the
Hessian of the complementarity potential
$S_{\cA}(x) = \tfrac{1}{2} H_{\cA}(x,x) + \cdots$\,.
The $m$-th trace-log term
$\tfrac{1}{2m}\operatorname{Tr}((C_1 B_2)^m)$ therefore probes the
$2m$-th jet of the product of complementarity potentials at a
sewing node. In particular, $\mathfrak{N}_g$ is a quartic
invariant: it sits at the same degree as the quartic contact
shadow $\mathfrak{Q}_{\cA}$ and, for uniform-weight algebras
(or any multi-generator family at genus~$1$) where
$\mathrm{obs}_g = \kappa \cdot \lambda_g$ \textup{(}UNIFORM-WEIGHT\textup{)} is proved, is controlled by the same data.
This is the mechanism by which the Gaussian sewing calculus
recovers algebra-dependent modular information beyond $\kappa$.
\end{remark}

\begin{remark}[Connection to the archetype classification]
\label{rem:holo-comp-archetype-connection}
The vanishing or nonvanishing of the first nonlinear anomaly
$\mathfrak{N}_g$ partitions the standard families by Gaussian
complexity:
\begin{itemize}
\item \emph{Heisenberg.}\;
 $B = C = 0$ (the complementarity potential is quadratic).
 The anomaly vanishes identically: $\mathfrak{N}_g = 0$ at all
 genera. Gaussian composition is exact and linear.
\item \emph{Affine $\widehat{\mathfrak{sl}}_2$.}\;
 The shadow obstruction tower terminates at cubic order, so $C_1 B_2$ is
 nilpotent of order $\leq 2$ on the complementarity Hessian.
 The square $(C_1 B_2)^2$ typically survives, giving a nonzero
 $\mathfrak{N}_g$.
\item \emph{$\beta\gamma$.}\;
 The shadow obstruction tower terminates at quartic order; $\mathfrak{N}_g$ is
 the leading nonlinear anomaly and captures the full quartic
 contact invariant $\mathfrak{Q}^{\mathrm{contact}}_{\beta\gamma}$.
\item \emph{Virasoro, $\mathcal{W}_N$.}\;
 The shadow obstruction tower is infinite. All trace-log terms
 $\operatorname{Tr}((C_1 B_2)^m)$ contribute, and the anomaly
 series does not truncate. The full metaplectic cocycle carries
 infinite-order modular data.
\end{itemize}
This is the Gaussian-kernel restatement of the archetype trichotomy
(Theorem~\ref{thm:nms-archetype-trichotomy}).
\end{remark}

\begin{remark}[Shadow tower as $A_\infty$ coproduct corrections (Vol~III)]
\label{rem:shadow-ainfty-coproduct-vol3}
In the Vol.~III coproduct normalization, the shadow invariants $S_k$ are
the coefficients of the $A_\infty$ corrections $\delta^{(k)}$ to the
Yangian coproduct $\Delta_z$:
\[
 \Delta_z^{A_\infty}
 =
 \Delta_z^{\mathrm{Yangian}}
 + \sum_{k \geq 3} \hbar^{k-1}\delta^{(k)}.
\]
Here $\delta^{(3)}$ has coefficient $\alpha=2$ from
$m_3(T,T,T)=-2T$, and $\delta^{(4)}$ has coefficient $S_4=10/27$
from $m_4(T,T,T,T)=(40/27)T$. For class~$G$, all
$\delta^{(k)}$ vanish and the Yangian coproduct is exact. For
class~$M$, the coproduct carries infinitely many corrections with
coefficients $\{S_k\}$. Thus the shadow tower both classifies the
averaged chiral algebra and supplies the $A_\infty$ correction
coefficients in the chiral quantum group coproduct; cf. Vol.~III,
Remark~\texttt{rem:ainfty-coproduct-shadow}.
\end{remark}

\begin{theorem}[Critical scalar curvature; \ClaimStatusConditional]
\label{thm:critical-dimension}
\index{critical dimension|textbf}
\index{shadow connection!flatness}
\index{Virasoro!critical dimension}
\index{W-algebra@$\mathcal{W}$-algebra!critical dimension}
The modular shadow connection
$\nabla^{\mathrm{hol}}_{g,n} =
d - \operatorname{Sh}_{g,n}(\Theta_\cA)$
is flat: $(\nabla^{\mathrm{hol}})^2 = 0$ by the MC equation. This
flatness is the full Maurer--Cartan statement; it is not equivalent
to the vanishing of the scalar characteristic.

For a one-parameter family $\cA_c$, the scalar diagonal of the
Koszul-dual curved fiber identity
$\dfib^{\,2}(\cA_c^!) = \kappa(\cA_c^!) \cdot \omega_g$
vanishes at the \emph{critical value}~$c_{\mathrm{crit}}$
where $\kappa(\cA_c^!) = 0$:
\begin{equation}\label{eq:critical-dimension-table}
\begin{array}{c|ccc}
\cA_c & \kappa(\cA_c) & \kappa(\cA_c^!)
 & c_{\mathrm{crit}}
\\[3pt]\hline
\rule{0pt}{12pt}
\mathrm{Vir}_c & c/2 & (26 - c)/2 & 26
\\[6pt]
V_k(\fg) &
 \dfrac{(k+h^\vee)\dim\fg}{2h^\vee}
 & \dfrac{(-k-h^\vee)\dim\fg}{2h^\vee}
 & k = -h^\vee
\end{array}
\end{equation}
At $c = c_{\mathrm{crit}}$:
\begin{enumerate}[label=\textup{(\roman*)}]
\item the scalar diagonal of the dual fiber curvature vanishes:
 $\dfib^{\,2}(\cA_{c_{\mathrm{crit}}}^!)|_{\mathrm{scal}}=0$;
\item the top-Hodge scalar obstruction of the dual vanishes:
 $\operatorname{obs}^{\mathrm{scal}}_g(\cA_{c_{\mathrm{crit}}}^!)
 =0$ for every $g\geq1$;
\item the complementarity splitting
 $\mathbf{Q}_g(\cA) \oplus \mathbf{Q}_g(\cA^!)
 \simeq H^*(\overline{\cM}_g, \cZ_\cA)$
 remains the Verdier eigenspace splitting of the strict flat ambient
 complex; it requires the C0/C1 hypotheses and is not a consequence
 of scalar vanishing alone.
\end{enumerate}
No termination of the higher shadow tower follows from
$\kappa(\cA_c^!)=0$ unless the higher shadow tensors
$\operatorname{Sh}_r(\cA_c^!)$ vanish by an independent computation.
\end{theorem}

\begin{proof}
For Kac--Moody, $\kappa(\cA) + \kappa(\cA^!) = 0$
(Theorem~\ref{thm:genus-universality}(ii)).
For~$\mathrm{Vir}_c$: $\kappa = c/2$,
$\kappa^! = \kappa(\mathrm{Vir}_{26-c}) = (26-c)/2$,
and $\kappa + \kappa^! = 13$
(Remark~\ref{rem:vir-vs-km-complementarity}).
Then $\kappa^! = 0$ at $c = 26$.

For~$V_k(\fg)$: $\kappa = (k+h^\vee)\dim\fg/(2h^\vee)$ and
$\kappa^! = (-k-h^\vee)\dim\fg/(2h^\vee)$; $\kappa^! = 0$
at $k = -h^\vee$ (critical level), where $\kappa = 0$ as well.
The scalar curvature identity gives
$\dfib^{\,2}|_{\mathrm{scal}}=\kappa^!\omega_g=0$ at the listed
critical value. The top-Hodge statement follows by applying
Proposition~\ref{prop:hodge-curvature-scalar-projection} to the dual:
\[
\operatorname{obs}^{\mathrm{scal}}_g(\cA^!)
=\kappa(\cA^!)\lambda_g.
\]
The final clause is Theorem~\ref{thm:quantum-complementarity-main}:
the direct sum is the idempotent splitting of the represented
Verdier involution on $\mathbf C_g(\cA)$, not a scalar curvature
calculation.
\end{proof}

\begin{remark}[Complementarity at higher degrees]
\label{rem:higher-degree-complementarity}
The scalar complementarity $\kappa(\cA) + \kappa(\cA^!) = 13$ (for Virasoro) extends to the cubic shadow: $S_3(\mathrm{Vir}_c) + S_3(\mathrm{Vir}_{26-c}) = 2 + 2 = 4$. At degree~$4$ and beyond, the duality sum $D_r(c) = S_r(c) + S_r(26{-}c)$ acquires nontrivial $c$-dependence, with denominators involving both the Lee--Yang factor $(5c{+}22)$ and its Koszul dual $(152{-}5c)$. The discriminant complementarity $\Delta(\cA) + \Delta(\cA^!) = 6960/[(5c{+}22)(152{-}5c)]$ (Theorem~\ref{thm:shadow-connection}(vi)) is the deepest degree-independent statement: its numerator is a universal constant.
\end{remark}

\begin{remark}[Complementarity numerator]
\label{rem:complementarity-numerator}
\index{complementarity!constant numerator}
The numerator $6960 = 2^4 \cdot 3 \cdot 5 \cdot 29 = 40 \cdot 174$ in the discriminant complementarity formula $\Delta + \Delta' = 6960/[(5c{+}22)(152{-}5c)]$ is a \emph{universal constant}: it depends on neither $c$ nor the specific algebra within the Virasoro family. The factor $40$ is the $c$-independent numerator of $\Delta(c) = 40/(5c{+}22)$, which arises from $8\kappa(c) Q^{\mathrm{contact}}(c) = 8 \cdot (c/2) \cdot 10/[c(5c{+}22)] = 40/(5c{+}22)$. At self-duality $c = 13$ this gives $\Delta(13) = 40/87$, and $174 = (5c{+}22 + 152{-}5c)|_{c=13} = 87 + 87$ is the sum of the two Lee--Yang denominators at self-duality.
\end{remark}

\begin{remark}[Complementarity for $\mathcal{W}$-algebras:
algebraic origin and Slodowy-slice perspective]
\label{rem:complementarity-slodowy}
\index{complementarity!Slodowy slice perspective}
\index{complementarity!Feigin--Frenkel origin}
\index{Slodowy slice!complementarity}
For $\mathcal{W}$-algebras obtained by DS reduction, the
complementarity sum
$\kappa(\mathcal{W}^k) + \kappa(\mathcal{W}^{k'})
= \mathrm{const}$
(Proposition~\ref{prop:ds-package-functoriality}(iv))
has a purely \emph{algebraic} origin in the Feigin--Frenkel
involution.

The central charge of $\mathcal{W}_N^k = \mathcal{W}^k(\fg,
f_{\mathrm{prin}})$ is
(Theorem~\ref{thm:wn-obstruction})
\begin{equation}\label{eq:wn-central-charge-complement}
c(t) \;=\; (N{-}1)\Bigl[1 -
\frac{N(N{+}1)(t{-}1)^2}{t}\Bigr],
\qquad t = k + h^\vee,
\end{equation}
and Feigin--Frenkel duality sends $t \mapsto -t$. Since
$c(t) + c(-t)$ eliminates the odd powers of~$t$:
\[
c(t) + c(-t) \;=\; 2(N{-}1)(2N^2{+}2N{+}1),
\]
the curvature sum
$\kappa(\mathcal{W}^k) + \kappa(\mathcal{W}^{k'})
= \varrho(\fg) \cdot [c(t) + c(-t)]$ is a
constant determined by~$\fg$ alone.
This is an identity of \emph{rational functions in~$t$},
not a dimensional or metric identity on~$\fg^*$:
the ambient curvature $\kappa(\widehat{\fg}_k)
= \dim(\fg)(k{+}h^\vee)/(2h^\vee)$ is itself
level-dependent.

Through the PBW identification
$\operatorname{gr}_F \mathcal{W}^k(\fg,f) \cong
\bC[J_\infty(S_f)]$
(Remark~\ref{rem:jet-slodowy}), the curvature
$\kappa(\mathcal{W}^k)$ lives on the arc-space geometry of
the Slodowy slice~$S_f$. The complementarity then says:
the genus-$1$ obstructions on the two Feigin--Frenkel-dual
arc spaces $J_\infty(S_f)$ at levels $t$ and $-t$ sum to
a constant, an algebraic shadow of the $t \mapsto -t$
symmetry on the space of Poisson parameters.
\end{remark}

\subsection{Polyakov's anomaly cancellation and complementarity}
\label{subsec:polyakov-complementarity}

\begin{remark}[Complementarity as Polyakov cancellation]
\label{rem:complementarity-polyakov-cancellation}
\index{Polyakov anomaly!complementarity interpretation}
\index{complementarity!Polyakov cancellation}
The degree-$2$ complementarity identity
$\kappa(\cA) + \kappa(\cA^!) = 13$ specialized to the Virasoro
family reads $\kappa(\mathrm{Vir}_c) + \kappa(\mathrm{Vir}_{26-c})
= c/2 + (26{-}c)/2 = 13$. This is the algebraic form of Polyakov's
conformal anomaly cancellation \cite{Polyakov1981}: at the critical
dimension $c = 26$, the matter sector (central charge~$26$,
realized by $26$ free bosons) contributes
$\kappa_{\mathrm{matter}} = 13$, while the $bc$ ghost system
(central charge $c_{\mathrm{ghost}} = -26$) contributes
$\kappa_{\mathrm{ghost}} = -13$, and their sum vanishes. In
bar-cobar language, the ghost sector \emph{is} the Koszul
dual~$\cA^!$: the BRST differential is the bar differential, ghost
number is the cohomological degree. Two distinct identities are at
play:
the complementarity sum
$\kappa(\cA) + \kappa(\cA^!) = 13$ holds for \emph{all}~$c$
(Proposition~\ref{prop:complementarity-landscape}(iii));
the physical anomaly cancellation
$\kappa_{\mathrm{eff}} := \kappa(\mathrm{matter}) +
\kappa(\mathrm{ghost}) = 0$ holds at the critical dimension $c = 26$
alone, where $\kappa(\mathrm{matter}) = 13$ and
$\kappa(\mathrm{ghost}) = -13$.
\end{remark}

\begin{proposition}[Non-critical complementarity and the Liouville
sector; \ClaimStatusConditional]
\label{prop:non-critical-liouville}
\index{Liouville sector!complementarity}
\index{Virasoro algebra!non-critical curvature}
For $c \neq 0$, the Virasoro bar complex is curved:
$d_{\barB}^2 = \kappa(\mathrm{Vir}_c) \cdot \omega_g$ with
$\kappa(\mathrm{Vir}_c) = c/2$. In Polyakov's non-critical
string theory, the Liouville sector of central charge
$c_L = 26 - c$ provides the complementary curvature
$\kappa(\mathrm{Vir}_{c_L}) = (26{-}c)/2$. The total matter +
Liouville curvature is
\[
\kappa(\mathrm{Vir}_c) + \kappa(\mathrm{Vir}_{c_L})
\;=\; \frac{c}{2} + \frac{26 - c}{2}
\;=\; 13,
\]
which is the residual anomaly absorbed by the ghost sector
$\cA^! = \mathrm{Vir}_{26-c}$. The ambient complementarity
theorem \textup{(}Theorem~\textup{\ref{thm:quantum-complementarity-main})}
encodes Polyakov's total anomaly cancellation
$c_{\mathrm{matter}} + c_{\mathrm{Liouville}} + c_{\mathrm{ghost}}
= 0$: the matter sector contributes
$Q_g(\mathrm{Vir}_c)$, the Liouville sector contributes
$Q_g(\mathrm{Vir}_{26-c})$, and complementarity forces
$Q_g(\mathrm{Vir}_c) + Q_g(\mathrm{Vir}_{26-c})
= H^*(\overline{\mathcal{M}}_g, \mathcal{Z}(\mathrm{Vir}_c))$.
\end{proposition}

\begin{proof}
The fiberwise curvature identity $\dfib^{\,2} = \kappa \cdot \omega_g$
applied to $\cA = \mathrm{Vir}_c$ gives
$\kappa(\mathrm{Vir}_c) = c/2$. The Koszul dual satisfies
$\mathrm{Vir}_c^! = \mathrm{Vir}_{26-c}$
(Example~\ref{ex:virasoro-koszul-dual}), hence
$\kappa(\mathrm{Vir}_c^!) = (26{-}c)/2$. The Liouville sector at
$c_L = 26 - c$ has $\kappa(\mathrm{Vir}_{c_L}) = (26{-}c)/2
= \kappa(\mathrm{Vir}_c^!)$, confirming that the Liouville field
is the curvature-dual of the matter field. Summing:
$\kappa_{\mathrm{matter}} + \kappa_{\mathrm{Liouville}} = 13$,
the Virasoro self-dual value, and application of
Theorem~\ref{thm:quantum-complementarity-main} at genus~$g$ yields
the full decomposition.
\end{proof}

\begin{corollary}[Degree-$4$ discriminant cancellation;
\ClaimStatusConditional]
\label{cor:complementarity-discriminant-cancellation}
\index{discriminant!complementarity}
\index{Polyakov anomaly!degree-four generalization}
The complementarity of discriminants
\[
\Delta(\mathrm{Vir}_c) + \Delta(\mathrm{Vir}_{26-c})
\;=\;
\frac{6960}{(5c{+}22)(152{-}5c)}
\]
\textup{(}Theorem~\textup{\ref{thm:shadow-connection}(vi))} is the
degree-$4$ generalization of Polyakov's anomaly cancellation: at
degree~$2$ the curvatures add to~$13$\textup{;}
at degree~$4$ the critical discriminants
$\Delta = 8\kappa S_4$ satisfy a complementarity with universal
numerator $6960 = 2^4 \cdot 3 \cdot 5 \cdot 29$. The constant
numerator reflects the same structural rigidity as the central
charge sum $c + c^! = 26$: the total degree-$4$ anomaly is
independent of the position within the Virasoro family.
\end{corollary}

\begin{proof}
Immediate from Theorem~\ref{thm:shadow-connection}(vi) and the
factorization $6960 = 40 \cdot 174$ established in
Remark~\ref{rem:complementarity-numerator}. The degree-$2$
specialization recovers $\kappa(c) + \kappa(26{-}c) = c/2 +
(26{-}c)/2 = 13$.
\end{proof}

\subsection{Conjectural bulk enhancement}
\label{subsec:holo-comp-bulk-conjecture}

\begin{conjecture}[Bulk enhancement and universal modular class;
\ClaimStatusConjectured]
\label{conj:holo-comp-bulk-enhancement}
\index{bulk factorization algebra!enhancement conjecture|textbf}
Let $\cA$ be a modular Koszul chiral algebra on~$X$. Then there
exists an enhanced bulk package $\mathfrak{Bulk}(\cA)$ (a
factorization algebra on a genuine holomorphic-topological bulk
space with boundary~$X$) whose protected boundary restriction
identifies with $\mathbb{H}_X(\cA)$, whose full higher-genus tower
is governed by a universal Maurer--Cartan element $\Theta_{\cA}$
lifting the uniform-weight scalar shadow $\kappa(\cA)\sum_{g \geq 1}\lambda_g$ \textup{(UNIFORM-WEIGHT)},
and whose quadratic truncation is exactly the metaplectic trace-log
anomaly series of
Theorem~\textup{\ref{thm:holo-comp-metaplectic-cocycle}}. After
choosing a compatible metaplectic structure, the Gaussian semigroup
of this section extends to a nonlinear kernel calculus closed under
modular operadic gluing.
\end{conjecture}

\begin{remark}[The proved package and the remaining frontier]
\label{rem:holo-comp-proved-and-frontier}
The content of this section assembles entirely from the proved core.
The protected dual transform $\mathbb{H}_X$ is a contravariant
involution recovering the Koszul dual in the coderived factorization
category. The genus-$g$ ambient complex is a shifted cotangent object
with palindromic spectral data under the protected self-duality
hypothesis. The linearized phase space admits canonical Heisenberg,
Fock, Fourier, and Weyl quantizations. Gaussian kernels compose by
Schur complement with a determinant anomaly whose metaplectic
extension satisfies a $2$-cocycle law. The first nonlinear anomaly
$\mathfrak{N}_g = \tfrac{1}{4}\operatorname{Tr}((C_1 B_2)^2)$ is
the first obstruction to naive sewing, sitting at quartic order in
exact alignment with the shadow obstruction tower.

The all-degree Maurer--Cartan element $\Theta_{\cA}$ exists by the
bar-intrinsic construction
(Theorem~\ref{thm:mc2-bar-intrinsic};
Theorem~\ref{thm:recursive-existence}).
What remains conjectural is the passage from this Gaussian calculus
to the full nonlinear modular kernel category: the construction of
$\mathfrak{Bulk}(\cA)$ as a three-dimensional bulk object.
The shadow obstruction tower provides the finite-order approximations; the bulk
enhancement provides the geometric realization.
\end{remark}

\begin{remark}[Cotangent geometry and the modular factorization adjunction]
\label{rem:cotangent-geometry-platonic}
\index{complementarity potential!cotangent interpretation}
\index{factorization adjunction!complementarity}
The complementarity potential~$S_{\cA}$
(Definition~\ref{def:complementarity-potential}) and its Legendre
dual~$S_{\cA^!}$
(Proposition~\ref{prop:legendre-duality-potentials}) realize dual
theories as \emph{Legendre-dual exact Lagrangians} in the ambient
$(-1)$-shifted symplectic moduli problem. In the factorization
envelope
(Theorem~\ref{thm:platonic-adjunction}), this upgrades
complementarity from a decomposition statement
(Theorem~\ref{thm:quantum-complementarity-main}) to true cotangent
geometry: the deformation space is locally the graph
$p = \partial S_{\cA}/\partial q$ inside a shifted cotangent
bundle, and passage to the dual theory is the formal Legendre
transform. The shadow obstruction tower provides the Taylor expansion of
this potential: $\kappa$ is the Hessian, $\mathfrak{C}$ the
cubic jet, $\mathfrak{Q}$ the quartic jet, and $\Theta$ the
all-orders completion.

For independent sums $L = L_1 \oplus L_2$ with vanishing mixed
OPE, the complementarity potential splits:
$S_L = S_{L_1} + S_{L_2}$
(Proposition~\ref{prop:independent-sum-factorization}),
and each shadow jet factors correspondingly. The Legendre-dual
Lagrangian geometry is therefore functorial across independent
orthogonal sums.
\end{remark}

\section{Scalar conductor separations}
\label{sec:hgc-supplement}

\begin{proposition}[B-family scalar conductor comparison; \ClaimStatusConditional]
\label{rem:hgc-B-family}
Assume the Vol.~III Mukai--Heisenberg witness is imported with
Mukai lattice
\[
  H^*(K3,\mathbb Z) \simeq II_{4,20},
  \qquad c_+\!\left(II_{4,20}\right)=4,
\]
and with the scalar normalization
\[
  K^{\kappa_{\mathrm{ch}}}\!\left(\mathcal H_{\mathrm{Muk}}(K3)\right)
  =2\,c_+\!\left(II_{4,20}\right).
\]
Then
\[
  K^{\kappa_{\mathrm{ch}}}\!\left(\mathcal H_{\mathrm{Muk}}(K3)\right)=8,
  \qquad
  \{0,13,250/3,98/3\}\cup\{8\}
  =\{0,8,13,250/3,98/3\}.
\]
If, in addition, the Vol.~III Humbert--Lusztig normalizations
\[
  \hbar^2 K^{\kappa_{\mathrm{ch}}}=-1,
  \qquad
  \ell=K^{\kappa_{\mathrm{ch}}}
\]
are imposed, then $\hbar^2=-1/8$ and $\ell=8$. These are scalar
normalization comparisons. They do not identify the ambient center
complex $\mathbf C_g(\cA)$, the derived chiral center
$Z^{\mathrm{der}}_{\mathrm{ch}}(\cA)$, a K3 sigma-model, a
Mukai--Heisenberg algebra, a Borcherds--Kac--Moody algebra, a Lusztig
root-of-unity category, or a Humbert monodromy category.
\end{proposition}

\begin{proof}
Mukai's K3 lattice convention gives
$H^*(K3,\mathbb Z)\simeq II_{4,20}$, so the positive rank is
$c_+(II_{4,20})=4$ \cite{Mukai1987}. Under the displayed scalar
normalization this gives $K^{\kappa_{\mathrm{ch}}}=2\cdot4=8$, and
adjoining this scalar to the four standard archetype constants gives
the stated five-element bucket. The equations
$\hbar^2K^{\kappa_{\mathrm{ch}}}=-1$ and
$\ell=K^{\kappa_{\mathrm{ch}}}$ then give $\hbar^2=-1/8$ and
$\ell=8$ by substitution. No derived-centre or Hochschild statement
follows from these equalities: Theorem~\ref{thm:quantum-complementarity-main}
acts on a center local system over $\overline{\mathcal M}_g$ together
with a represented Verdier involution, and neither datum is constructed
in this proposition for the K3, Humbert, Lusztig, or Borcherds--Kac--Moody
objects.
\end{proof}

\begin{proposition}[Separating-node automorphic scalar; \ClaimStatusConditional]
\label{rem:hgc-sepnode}
Assume the automorphic character identity
\[
  \mathrm{Ch}(\mathbf H_{\Delta_5})(Z)
  =\frac{1}{\Phi_{10}^{\mathrm{un}}(Z)}
\]
with the Gritsenko--Nikulin normalization
$\Phi_{10}^{\mathrm{un}}=\Delta_5^2$ and with the leading
Fourier--Jacobi coefficient fixed. At a separating degeneration
$Z\to\operatorname{diag}(\tau,\rho)$ the leading boundary term is
\[
  \frac{1}{\Phi_{10}^{\mathrm{un}}(Z)}
  \sim
  \frac{1}{\eta(\tau)^{24}}\cdot\frac{1}{\phi_{10,1}(\rho)}.
\]
On this automorphic Borcherds--Kac--Moody lane,
\[
  \kappa_{\mathrm{BKM}}(\mathfrak g_{\Delta_5})
  =\frac{c_\Lambda(0)}{2}
  =\operatorname{wt}(\Delta_5)
  =5.
\]
This scalar $5$ is not the B-family scalar $8$ of
Proposition~\ref{rem:hgc-B-family}, and neither scalar is a
Theorem~\ref{thm:quantum-complementarity-main} computation of a
derived center.
\end{proposition}

\begin{proof}
The displayed separating-node expansion is precisely the assumed
Fourier--Jacobi boundary normalization. Gritsenko--Nikulin identify the
primitive denominator with the weight-$5$ form $\Delta_5$, and the
chosen unnormalized genus-$2$ form is its square
\cite{GritsenkoNikulin1998}. With the Borcherds denominator convention
$\kappa_{\mathrm{BKM}}=c_\Lambda(0)/2$, the same normalization gives
$c_\Lambda(0)/2=\operatorname{wt}(\Delta_5)=5$. The argument is
automorphic and scalar: it contains no center local system over
$\overline{\mathcal M}_g$, no represented Verdier involution, and no
Hochschild complex. It therefore does not prove a derived-centre
complementarity theorem for the K3 or Borcherds--Kac--Moody object.
\end{proof}
